\documentclass[11pt,a4paper]{article}

%──────────────────────────────────────────────────────
%  Codificación, idioma y fuente
%──────────────────────────────────────────────────────
\usepackage[utf8]{inputenc}
\usepackage[T1]{fontenc}
\usepackage[spanish,mexico]{babel}
\decimalpoint

\usepackage{geometry}
\geometry{margin=2.5cm, top=3cm, bottom=3cm}

\usepackage{lmodern}
\renewcommand{\familydefault}{\sfdefault}
\usepackage{microtype}

%──────────────────────────────────────────────────────
%  Matemáticas, tablas, figuras
%──────────────────────────────────────────────────────
\usepackage{amsmath,amssymb}
\usepackage{booktabs}
\usepackage{tabularx}
\usepackage{array}
\usepackage{multirow}
\usepackage{graphicx}
\usepackage{float}
\usepackage{caption}
\captionsetup{font=small, labelfont=bf}

%──────────────────────────────────────────────────────
%  Listas y diseño
%──────────────────────────────────────────────────────
\usepackage{enumitem}
\usepackage{parskip}

%──────────────────────────────────────────────────────
%  Encabezados, secciones, cajas, hipervínculos
%──────────────────────────────────────────────────────
\usepackage{fancyhdr}
\usepackage{titlesec}
\usepackage{tcolorbox}
\usepackage{hyperref}
\usepackage[colorinlistoftodos]{todonotes}

%──────────────────────────────────────────────────────
%  Paleta editorial UAM-A
%──────────────────────────────────────────────────────
\definecolor{uamred}{RGB}{200,45,35}
\definecolor{uamlightred}{RGB}{248,232,232}
\definecolor{uamblue}{RGB}{28,28,28}
\definecolor{uamgray}{RGB}{100,100,100}
\definecolor{okgreen}{RGB}{56,142,60}
\definecolor{codebg}{RGB}{250,250,250}

%──────────────────────────────────────────────────────
%  Secciones
%──────────────────────────────────────────────────────
\titleformat{\section}
  {\large\bfseries\color{uamblue}}
  {\textcolor{uamred}{\thesection}}{0.8em}{}
  [\vspace{3pt}\color{uamred}\hrule height 0.5pt\vspace{2pt}]
\titleformat{\subsection}
  {\normalsize\bfseries\color{uamblue}}{\thesubsection}{0.8em}{}
\titleformat{\subsubsection}
  {\normalsize\itshape\color{uamgray}}{\thesubsubsection}{0.8em}{}
\titlespacing*{\section}      {0pt}{3.5ex plus 1ex minus .5ex}{1.5ex}
\titlespacing*{\subsection}   {0pt}{2.5ex plus .5ex}{.8ex}
\titlespacing*{\subsubsection}{0pt}{1.8ex plus .3ex}{.4ex}

%──────────────────────────────────────────────────────
%  Cabecera y pie
%──────────────────────────────────────────────────────
\pagestyle{fancy}\fancyhf{}
\fancyhead[L]{\small\color{uamgray}Eficiencia Terminal --- División CBI · UAM-A}
\fancyhead[R]{\small\color{uamgray}Reporte Técnico}
\fancyfoot[C]{\small\thepage}
\renewcommand{\headrulewidth}{0.4pt}
\setlength{\headheight}{15pt}

%──────────────────────────────────────────────────────
%  Cajas
%──────────────────────────────────────────────────────
\tcbuselibrary{skins,breakable}
\newtcolorbox{notabox}{enhanced,breakable,
  colback=uamlightred,colframe=uamlightred,
  borderline west={3pt}{0pt}{uamred},
  sharp corners,left=10pt,right=8pt,top=5pt,bottom=5pt,
  before upper={\textbf{\color{uamred}Nota}\enspace}}
\newtcolorbox{warnbox}{enhanced,breakable,
  colback=black!4,colframe=black!4,
  borderline west={3pt}{0pt}{uamgray},
  sharp corners,left=10pt,right=8pt,top=5pt,bottom=5pt,
  before upper={\textbf{Advertencia}\enspace}}
\newtcolorbox{okbox}{enhanced,breakable,
  colback=green!6,colframe=green!6,
  borderline west={3pt}{0pt}{okgreen},
  sharp corners,left=10pt,right=8pt,top=5pt,bottom=5pt}

%──────────────────────────────────────────────────────
%  Metadatos e hipervínculos
%──────────────────────────────────────────────────────
\hypersetup{
  pdftitle={Eficiencia Terminal en los Planes 2020 --- División CBI UAM-A},
  pdfauthor={Dr. César Simón López-Monsalvo, UAM-Azcapotzalco},
  colorlinks=true,linkcolor=uamblue,urlcolor=uamred}

%──────────────────────────────────────────────────────
%  Macros locales
%──────────────────────────────────────────────────────
\newcommand{\plan}[1]{\textsc{#1}}
\newcommand{\uea}[1]{\texttt{#1}}
\newcommand{\ET}[1]{\mathrm{ET}_{(#1)}}
\newcommand{\ETtit}{\mathbb{E}[T\,|\,\mathrm{tit}]}
\newcommand{\Pacc}{P_{\mathrm{acc}}}

%══════════════════════════════════════════════════════
\begin{document}
%══════════════════════════════════════════════════════

%──────────────────────────────────────────────────────
\begin{titlepage}
\raggedright\vspace*{0pt}
\includegraphics[width=0.54\textwidth,trim={48bp 64bp 473bp 35bp},clip]%
  {/Users/nowhere/Claude_Projects/Personal/logo_UAM.png}\par
\vspace{0.3cm}
{\footnotesize\color{uamgray}\MakeUppercase{División de Ciencias Básicas e Ingeniería}}\par
\vfill
{\color{uamred}\rule{\linewidth}{1pt}}\par\vspace{0.9cm}
{\fontsize{28}{34}\selectfont
  Eficiencia Terminal en los\\[0.05em]
  Planes de Estudio 2020\\[0.25em]
  de la División CBI}\par\vspace{0.5cm}
{\large\color{uamgray}
  Análisis de seriaciones, modelo probabilístico y hallazgos}\par\vspace{0.9cm}
{\color{uamred}\rule{\linewidth}{1pt}}\par\vspace{1cm}
{\normalsize
  Reporte técnico \quad|\quad
  %Destinatarios: jefes de departamento, Junta de CBI\\[0.3em]
  Datos: cohortes UAM-A 16I--25O;\enspace matrícula activa 25-O}\par
\vfill
\begin{minipage}[b]{0.6\linewidth}
  {\small\color{uamgray}
    Dr.\ César Simón López-Monsalvo\\[0.2em]
    Dr.\ Román Anselmo Mora Gutiérrez\\[0.2em]
    División de Ciencias Básicas e Ingeniería\\[0.2em]
    \texttt{cslm@azc.uam.mx},
    \texttt{mgra@azc.uam.mx}
    }
\end{minipage}%
\hfill
\begin{minipage}[b]{0.35\linewidth}
  \raggedleft{\small\color{uamgray}Versión 5.2 (revisión crítica)\\[0.2em]\today}
\end{minipage}\par\vspace{0.4cm}
\end{titlepage}

\tableofcontents
\clearpage

%══════════════════════════════════════════════════════
\section{Introducción}
\label{sec:introduccion}
%══════════════════════════════════════════════════════

\subsection*{Motivación}

La eficiencia terminal universitaria en México registra una de las brechas
más pronunciadas entre el plazo del plan de estudios y el tiempo real de egreso en el
conjunto de la Organización para la Cooperación y el Desarrollo Económicos (OCDE~\cite{OCDE2024}).
En las ingenierías, donde la carga crediticia y las cadenas de seriación
son especialmente complejas, esta brecha tiene consecuencias directas: retrasa
la incorporación al mercado laboral, aumenta el costo personal y familiar de
la formación, erosiona la cobertura efectiva del sistema y representa un
indicador crítico de la calidad institucional ante organismos acreditadores.
En la División de Ciencias Básicas e Ingeniería (CBI) de la Universidad Autónoma Metropolitana Azcapotzalco (UAM-A), el tiempo promedio de egreso de
los 4\,029 estudiantes que completaron su licenciatura entre 2018 y 2025 fue
de \textbf{20.6 trimestres} ---72\% por encima del plazo del plan de
12 trimestres--- con Electrónica en el extremo desfavorable (23.2 trim) e
Industrial en el favorable (19.4 trim).
Esta dispersión no es aleatoria: responde a diferencias estructurales en los
planes de estudio que pueden medirse, modelarse y, en consecuencia, corregirse.

La aprobación de los Planes 2020 supuso una oportunidad de modernización
curricular, pero no estuvo acompañada de un análisis cuantitativo del impacto
de las cadenas de seriación sobre las trayectorias de los estudiantes.
El presente reporte cubre esa brecha con un diagnóstico basado en datos
históricos propios de la institución, un modelo de simulación calibrado y
un conjunto de herramientas de análisis reproducibles.

\subsection*{Marco normativo}

La Ley General de Educación Superior (LGES~\cite{LGES2021}, Arts.~10 y~40--41) establece
la obligación de las instituciones de educación superior de garantizar
condiciones que permitan a los estudiantes concluir sus estudios con calidad
y pertinencia.
El Reglamento de Estudios Superiores de la UAM (RES~UAM~\cite{RESUAM}, Art.~58)
establece un plazo mínimo de ocho trimestres y máximo de diez años para las
licenciaturas; cada Plan~2020 de CBI fija en doce trimestres la duración nominal
del programa ---el \mbox{<<plazo del plan>>} en la práctica institucional---.
Los Arts.~16.V y~44.XI del RES~UAM imponen la baja definitiva tras la acumulación
de cinco evaluaciones No Acreditadas (NA) en una misma Unidad de Enseñanza-Aprendizaje (UEA).
Un currículo cuya combinación de carga crediticia y cadenas de seriación
hace estadísticamente improbable egresar en el plazo establecido
---como lo demuestran los modelos de este reporte--- genera una tensión
normativa que la División tiene la obligación de atender.
El análisis jurídico detallado de cada hallazgo y su respaldo legal se
desarrolla en el documento complementario \texttt{07\_analisis\_juridico.pdf};
el presente reporte se concentra en la evidencia técnica y cuantitativa.

\subsection*{Antecedente directo}

Este reporte se desprende del análisis del Tronco General Nuevo  (TGNA~\cite{TGNA2026}), que demostró por primera vez que las cadenas de seriación
del Tronco General actual reducen la fracción de la cohorte con acceso a los
cursos del tramo final a entre el 2\% y el 5\%.
Ese hallazgo, circunscrito a las once UEAs del Tronco General, motivó extender
el análisis al currículo completo de los diez programas de ingeniería:
1\,598 UEAs, 1\,203~aristas de seriación estructurales (Tabla~\ref{tab:grafos}),
499 evaluadas individualmente para impacto marginal sobre~$\ETtit$,
registros de evaluación 16I--25O y 6\,562~estudiantes activos como horizonte de proyección.

\subsection*{Principales resultados}

El hallazgo central del documento es estructural: el currículo 2020 fue diseñado
para un estudiante que inscribe $\approx 45$~cr/trim, pero la población real
inscribe históricamente $\approx 27$--$33$~cr/trim.
Esa brecha, anterior a cualquier análisis de seriaciones, es condición suficiente
para hacer estadísticamente improbable el egreso a 12~trimestres en los diez
programas.
La observación histórica lo confirma: las cohortes 2020--2021 registraron
$\ET{12}=1.52\%$ ($44/2\,892$; IC~95\%: $[1.1\%,\,2.0\%]$), coherente con
el orden de magnitud predicho por el modelo bajo toda la gama de parámetros
plausibles de calibración ($\ET{12}^{\mathrm{MC}}\approx 0.000\%$ en los diez
planes; los dos intervalos no se solapan: véase~\S\ref{sec:hallazgo}).
La incertidumbre combinada del resultado es $\ET{12}\in[0.000\%,\,2.04\%]$ con
confianza del~95\%.
El~62.5\% de los 6\,562~estudiantes activos egresará, pero entre los trimestres
17 y~28.

El modelo identifica dos mecanismos por los que el desajuste de carga se expresa
en el grafo de seriaciones.
En todos los planes la restricción dominante es la carga crediticia: incluso
eliminando todas las seriaciones, el diagnóstico principal no cambia.
La profundidad topológica máxima es de 9~niveles (Computación y Electrónica),
imponiendo una cota de 10~trimestres que permanece por debajo del umbral de~12.

El análisis de impacto marginal revela que la distribución de beneficios sobre las
499~aristas de seriación sigue una ley de saturación convexa: cada plan alcanza
el~80\% de su ganancia máxima eliminando entre~2 y~50~aristas ($N_{80}^{\mathrm{obs}}$,
según la Tabla~\ref{tab:dosis_respuesta}), y la arista de mayor impacto individual
(EDO $\to$ Circuitos Eléctricos~I) aporta $-1.07$~trimestres en Electrónica.
La ganancia máxima al eliminar todas las seriaciones oscila entre~0.3~trimestres
(Computación) y~2.2~trimestres (Electrónica); para los planes con restricción
dominante de carga, esta ganancia no supera los~0.5~trimestres.
En dos planes (Computación y Física) las curvas dosis-respuesta presentan tramos
no-monótonos: eliminar seriaciones más allá de un umbral empeora el pronóstico
por interacciones de grafo.
En todos los planes la remoción simultánea de seriaciones es subaditiva: la suma
de los beneficios individuales sobreestima el beneficio conjunto en un factor de
2--3 (índice de aditividad $\mathcal{A} \in [0.23, 0.89]$).
El optimizador de portfolio resuelve este problema identificando el subconjunto
óptimo de aristas con mayor beneficio neto tras descontar las interacciones
(véase la Tabla~\ref{tab:dosis_respuesta} y la \S\ref{sec:portfolio}).

Los análisis de correlación estadística confirman que el número de seriaciones
predice la tasa de egreso ($r_s = -0.80$, $p = 0.006$), y la fracción de UEAs
con tasa de aprobación inferior a~0.65 predice el rezago ($r_s = +0.80$ con
$\mathrm{Trim}_{50}$, $p = 0.005$), ambos robustos al controlar por la escala
del plan.
A nivel de UEA, el riesgo se concentra: el~18\% de las 1\,598~UEA-nodos
(aquellos con $a_i < 0.65$) concentra la mayor parte del riesgo acumulado de
rezago por plan.

El análisis del cuerpo docente (1\,810~profesores UAM-A activos en 16I--25O)
añade un nivel de resolución independiente: el~31.3\% de los~738~docentes
CBI opera con una tasa de aprobación media inferior al umbral de riesgo~0.65,
el Departamento de Ciencias Básicas registra la mediana departamental más baja
de la División y cubre todos los nodos críticos del Tronco General,
y la heterogeneidad docente media de un plan correlaciona con el retraso en
el egreso ($r_s = +0.78$, $p = 0.043$).
La varianza docente explica aproximadamente un~5\% de la varianza de~$\ETtit$
por encima del modelo de parámetros homogéneos.

El análisis del índice de desempeño vigente~$P$ (Acuerdo~551.4.4.1-2015)
demuestra que su componente~$\mathrm{ACT}$ (razón de créditos aprobados sobre
créditos inscritos) genera un incentivo racional al subregistro: inscribir menos
créditos eleva~$\mathrm{ACT}$ aunque el avance absoluto en el plan disminuya.
El índice alternativo de base entrópica $h=(\eta_t+\langle\eta\rangle)/2$,
donde $\eta_t=C_{A,t}/C_{\max}$ es la fracción acumulada de créditos aprobados
sobre el total del plan, extingue ese incentivo porque su único camino de mejora
es aprobar créditos absolutos.
La sustitución de~$P$ por~$h$ extingue el incentivo al subregistro estratégico
---la ventaja transitoria $\Delta P(2)=+0.111$ bajo~$P$ es estrictamente negativa
bajo~$h$ ($\Delta h(2)=-0.043$)--- y alinea el comportamiento racional con la
velocidad de avance en el plan; sin embargo, no elimina la penalización topológica
que~$P$ impone a los estudiantes bloqueados en planes de mayor profundidad.

Finalmente, el análisis global de complejidad sistémica sintetiza cinco
dimensiones del currículo ---topológica, probabilística, evaluativa, docente
y de recuperación--- en un índice compuesto $C\in[0,1]$ que clasifica los
diez planes en tres rangos: Alta (Comunicación, Química, Mecánica), Moderada
(seis planes) y Baja (Metalurgia).
El análisis de acoplamiento entre dimensiones revela una trampa auto-reforzante:
la dimensión de Recuperación (facilidad de superar UEAs reprobadas) correlaciona
negativamente con las dimensiones Probabilística ($r_s=-0.673$) y Docente
($r_s=-0.576$), lo que indica que los planes con mayor dificultad de aprobación
tampoco ofrecen mecanismos efectivos de recuperación.
Los diez planes se clasifican en cuatro arquetipos de complejidad estructural
(trampa docente, trampa evaluativa, plan frágil, plan asistido) que orientan
la selección de instrumentos de intervención específicos para cada caso.

El diagnóstico converge en la necesidad de actuar sobre cinco frentes de forma
simultánea: el topológico, el de aprobación por UEA, el de carga efectiva de
inscripción, el docente y el evaluativo.
Los cinco son necesarios; ninguno es suficiente de forma aislada.

\subsection*{Organización del documento}

El documento está estructurado en diecisiete secciones, un glosario final
y un apéndice de scripts.

La Sección~\ref{sec:hallazgo} resume el diagnóstico estructural en forma
compacta, con las cotas numéricas del modelo y la observación histórica.
La Sección~\ref{sec:contexto} presenta los dos datos de anclaje (eficiencia
terminal histórica y tiempo promedio real de egreso) y el antecedente del
Tronco General Nuevo Actualizado (TGNA~\cite{TGNA2026}).
La Sección~\ref{sec:sistema} describe la cadena de análisis: fuentes de datos,
1\,859~UEAs en el banco de datos (689~obligatorias CBI), 1\,203~aristas de
seriación estructurales de las cuales 499~fueron evaluadas individualmente
para impacto marginal (las~704~restantes presentan impacto negligible en el modelo),
y 6\,562~estudiantes activos.
Las Secciones~\ref{sec:cascada} y~\ref{sec:modelo} desarrollan los dos modelos
complementarios: la cascada estática (cota inferior analítica) y el modelo
topológico-estocástico (simulación Monte Carlo con $S=20\,000$ trayectorias
por plan).
La Sección~\ref{sec:resultados} reporta los resultados de eficiencia terminal
por plan e incluye la robustez al parámetro de fragilidad latente.
La Sección~\ref{sec:mecanismos} diagnostica los dos mecanismos del desajuste
curricular (restricción de carga y profundidad topológica) e incluye la
descomposición analítica de la carga efectiva histórica~$c_s$.
La Sección~\ref{sec:efecto} cuantifica el efecto marginal de cada seriación
individual y establece los escenarios globales de reforma.
La Sección~\ref{sec:correlaciones} analiza los determinantes estadísticos del
desempeño a nivel de plan, Unidad de Enseñanza-Aprendizaje (UEA), arista
y factores latentes; incluye la estructura factorial de los atributos de UEA.
La Sección~\ref{sec:perfil_docente} caracteriza los 1\,810~profesores UAM-A,
su relación con los nodos de mayor riesgo y la varianza docente explícita sobre
las proyecciones de~$\ETtit$.
La Sección~\ref{sec:dosis_respuesta} evalúa las curvas dosis-respuesta de
remoción combinada, la subaditividad de los efectos, las curvas no-monótonas,
el optimizador de portfolio y el modelo predictivo del índice de subaditividad;
concluye con la clasificación estratégica de los planes.
La Sección~\ref{sec:indice_desempenio} analiza el índice de desempeño vigente~$P$
y el índice alternativo~$h$: déficit de ritmo, incentivo al subregistro, doble
penalización topológica de~$P$, y extinción del incentivo bajo~$h$.
La Sección~\ref{sec:complejidad} construye el índice compuesto de complejidad
sistémica~$C$ sobre cinco dimensiones, clasifica los planes en tres rangos y
cuatro arquetipos, y analiza el acoplamiento entre dimensiones.
La Sección~\ref{sec:limitaciones} discute los límites del modelo e incluye las
cláusulas de alcance jurídico.
La Sección~\ref{sec:discusion} sintetiza en prosa continua todos los hallazgos
del documento y las acciones recomendadas sobre los cinco frentes de
intervención.
La Sección~\ref{sec:conclusiones} enuncia los doce hallazgos principales en
forma proposicional y formula la recomendación estructural sobre los cinco
frentes de intervención simultáneos.
El glosario al final del documento define todas las abreviaturas,
símbolos y factores del análisis factorial.

%══════════════════════════════════════════════════════
\section{Hallazgo central}
\label{sec:hallazgo}
%══════════════════════════════════════════════════════

El currículo 2020 fue diseñado para un estudiante que inscribe aproximadamente
40--45 créditos por trimestre ---la carga normal establecida en cada Plan de
Estudios---.
La población estudiantil real de CBI ha inscrito históricamente alrededor de
27--33~créditos por trimestre.
Esa brecha entre el estudiante supuesto en el diseño y el estudiante real es
la causa estructural primaria del retraso en el egreso: antes de analizar
una sola seriación, la aritmética de créditos hace que el egreso a
12~trimestres sea estadísticamente improbable bajo condiciones históricas.

\begin{okbox}
\textbf{Diagnóstico estructural.}
El currículo 2020 de CBI incorpora entre 376 y 419~créditos obligatorios.
A la carga efectiva histórica ($c_s \approx 27$--$33$~cr/trim), completar
esos créditos requiere entre 18 y 21~trimestres, independientemente de las
seriaciones.
El modelo Monte Carlo calibrado ($S=20\,000$ trayectorias por plan) predice
\begin{equation}
  \ET{12} = 0.000\% \quad
  \bigl(\text{IC~MC }95\%:\;[0.000\%,\;0.018\%]\bigr)
  \label{eq:ET12_modelo}
\end{equation}
en los diez programas bajo parámetros centrales de calibración.
La observación histórica (cohortes 2020--2021) arroja
\begin{equation}
  \ET{12}^{\,\mathrm{obs}} = 1.52\%
  \quad (44/2\,892;\;\text{IC~95\% Clopper-Pearson (CP): }[1.11\%,\;2.04\%]).
  \label{eq:ET12_obs}
\end{equation}
Los dos intervalos no se solapan.
El 62.5\% de los 6\,562~estudiantes activos egresará, pero entre los trimestres~17 y~28.
\end{okbox}

La brecha de 1.5~puntos porcentuales (pp) entre~(\ref{eq:ET12_modelo}) y~(\ref{eq:ET12_obs}) no
es error del modelo: es la huella de la heterogeneidad individual de carga,
que el modelo no captura al asumir $c_s$ homogéneo.
Los 44~estudiantes que egresaron en $\leq 12$~trimestres son, con alta
probabilidad, aquellos que operaron consistentemente con $c_s\approx 45$~cr/trim
---la carga de diseño del plan--- durante toda su licenciatura.
Su fracción en la cohorte ($\approx 1.5\%$) cuantifica con precisión
la distancia entre la población real y el alumnado implícito del currículo.

El análisis de sensibilidad a $c_s\pm 4$~cr/trim
(figura~\ref{fig:incertidumbre_ET12}, panel inferior)
confirma que $\ET{12}$ permanece en $0.000\%$ en todos los planes
para todo el rango de calibración plausible: el resultado es insensible a
la incertidumbre paramétrica en~$c_s$.
La incertidumbre combinada del resultado es:
\begin{equation}
  \ET{12} \;\in\; [0.000\%,\;2.04\%]
  \quad(\text{confianza }95\%),
  \label{eq:ET12_rango}
\end{equation}
donde el límite inferior es la cota MC del modelo y el límite superior es el
IC empírico de la observación histórica.

\begin{figure}[htbp]
  \centering
  \includegraphics[width=\textwidth]{../../Output/figuras/fig_incertidumbre_ET12.pdf}
  \caption{Robustez de~$\ET{12}$ (fracción de la cohorte que egresa en a lo sumo
    12~trimestres) frente a dos fuentes de incertidumbre.
    \emph{Panel superior:} intervalo de confianza (IC) al~95\% exacto de
    Clopper-Pearson sobre~$\ET{12}$ por plan, obtenido con $S=20\,000$
    trayectorias Monte Carlo; la banda gris señala la cota empírica histórica
    $[1.11\%,\,2.04\%]$.
    \emph{Panel inferior:} sensibilidad de~$\ET{12}$ a la carga efectiva
    histórica~$c_s$ variada en~$[26,34]$~cr/trim~($\pm4$~cr/trim respecto al
    valor calibrado por plan); en todos los planes el resultado es insensible a
    esta incertidumbre paramétrica.}
  \label{fig:incertidumbre_ET12}
\end{figure}

El análisis identifica dos mecanismos por los cuales el desajuste de carga
se expresa en el grafo de seriaciones, con pesos distintos según el plan.
En todos los planes, la restricción dominante es la carga crediticia: incluso
eliminando todas las seriaciones, el diagnóstico principal no cambia.
La profundidad topológica máxima es de 9~niveles (Computación y Electrónica),
imponiendo una cota de 10~trimestres que no alcanza el umbral de~12.
Las secciones que siguen documentan ambos mecanismos, el modelo que los
cuantifica, y el impacto marginal de las seriaciones individuales
---que constituye la segunda capa de intervención, relevante una vez
atendido el desajuste primario de carga.

%══════════════════════════════════════════════════════
\section{Contexto y datos de anclaje}
\label{sec:contexto}
%══════════════════════════════════════════════════════

\subsection{Antecedente: el Tronco General como caso piloto}

El análisis del Nuevo Tronco General ~\cite{TGNA2026} demostró que las
cadenas de seriación del Tronco General (TG) actual producen un filtro estadístico que lleva
la fracción de la cohorte con acceso a los cursos del tramo final a valores
del orden del 2--5\%.
El presente reporte extiende ese análisis al currículo completo de los diez
programas de ingeniería CBI.

\subsection{Ancla~1: eficiencia terminal histórica al plazo del plan de estudios}

\begin{table}[H]\centering
\caption{Egresos en el plazo establecido de 12~trimestres
         (Cuadro~1 de~\cite{TGNA2026})}
\label{tab:historico}
\small
\begin{tabular}{lccc}
\toprule
Cohorte & Ingreso & Egresados en $t\leq12$ & $\ET{12}$ \\
\midrule
2020 & 1\,446 & 42 & 2.9\% \\
2021 & 1\,446 & \phantom{0}2 & 0.14\% \\
\midrule
Pooled & 2\,892 & 44 & \textbf{1.52\%} \\
\bottomrule
\end{tabular}
\end{table}

El modelo \emph{no} utiliza este valor como ancla de calibración.
Lo compara como verificación independiente.
El modelo predice $\ET{12}=0.000\%$
(IC~MC~95\%: $[0.000\%,\,0.018\%]$) bajo los parámetros centrales
de calibración (módulo~30).
El valor histórico pooled es $1.52\%$ (IC~95\%~CP: $[1.11\%,\,2.04\%]$).
Los dos intervalos no se solapan: la diferencia de $1.5\,\text{pp}$ refleja la
heterogeneidad individual de carga no modelada.
La heterogeneidad entre cohortes ---$2.9\%$ en 2020 y $0.14\%$ en 2021---
es consistente con la variabilidad estocástica en la fracción de estudiantes
de alta carga en cada cohorte; el valor pooled de $1.52\%$ es la mejor
estimación del fenómeno bajo condiciones históricas ordinarias.

\subsection{Ancla~2: tiempo promedio real de egreso}

\begin{table}[H]\centering
\caption{Trimestres promedio reales a egreso, $\ETtit^{\,\mathrm{ofic}}$,
         período 2018--2025.
         Fuente: \texttt{Trimestres para terminar las licenciaturas 2018-2025.docx}.}
\label{tab:et_oficial}
\small
\begin{tabular}{llrr}
\toprule
Clave & Programa & Egresados & $\ETtit^{\,\mathrm{ofic}}$ \\
\midrule
amb & Ambiental    & 482 & 20.56 \\
civ & Civil        & 454 & 20.34 \\
com & Computación  & 456 & 20.31 \\
ele & Eléctrica    & 174 & 21.41 \\
elo & Electrónica  & 327 & 23.24 \\
fis & Física       & 366 & 20.25 \\
ind & Industrial   & 529 & 19.40 \\
mec & Mecánica     & 436 & 21.18 \\
met & Metalurgia   & 253 & 22.26 \\
qui & Química      & 552 & 19.41 \\
\midrule
\textbf{CBI} & & \textbf{4\,029} & \textbf{20.60} \\
\bottomrule
\end{tabular}
\end{table}

Esta es la ancla de calibración: el modelo ajusta un parámetro libre por plan
para reproducir $\ETtit^{\,\mathrm{ofic}}$.
Son medias de cohorte; el mínimo anual registrado es 17.2~trim
(Ingeniería Industrial, 2024), lo cual confirma que hay estudiantes
individuales que egresan antes de los 18~trimestres, aunque en proporción
muy reducida.

%══════════════════════════════════════════════════════
\section{Sistema de análisis}
\label{sec:sistema}
%══════════════════════════════════════════════════════

El sistema de análisis está compuesto por 36~scripts: 21~módulos en
Python~3.14, 6~en R~\cite{RCoreTeam2024}, 5~de gestión institucional
(\texttt{delivery/}) y 4~utilidades de calidad textual (\texttt{utils/});
organizados en cinco fases analíticas.
La cadena es completamente reproducible: partiendo de los datos crudos,
una corrida completa de las fases~1 y~2 (simulación y análisis estadístico)
tarda aproximadamente 4~minutos en hardware convencional.
Las fases~3--5 (extensiones cuantitativas, índices y complejidad) añaden
otros 6~minutos en los diez planes.

\subsection*{Fuentes de datos}

\begin{tabularx}{\linewidth}{@{}l X r@{}}
\toprule
\textbf{Fuente} & \textbf{Contenido} & \textbf{Escala} \\
\midrule
\texttt{diccionario\_ueas.json} &
  Historial de evaluación 16I--25O: tasas de aprobación~$a_i$, renuncia,
  reprobación, recuperación y vectores de probabilidad acumulada por
  intento~$P(k\leq k')$ para cada UEA &
  1\,859 UEAs \\[4pt]
\texttt{diccionario\_profesores.json} &
  Indicadores docentes 16I--25O: tasa de aprobación media~$a_d$, eficacia,
  eficiencia, rendimiento, probabilidad de expulsión~$P_{\mathrm{exp}}$ y
  lista de UEAs impartidas por profesor (clave \texttt{NO\_ECO}) &
  1\,810 docentes \\[4pt]
\texttt{planEstudios\_*.pdf} &
  Planes~2020 por programa: columna SERIACIÓN con códigos directos,
  prefijos \texttt{C} (corregistro), umbrales de créditos acumulados
  y carga normal/máxima por trimestre &
  10 planes \\[4pt]
\texttt{Matrícula activa\ldots.docx} &
  Inscritos activos por plan en el trimestre~25-O &
  6\,562 est. \\[4pt]
\texttt{Trimestres para terminar\ldots.docx} &
  Trimestres reales a egreso por plan y cohorte, período 2018--2025
  (ancla de calibración del modelo) &
  4\,029 egres. \\
\bottomrule
\end{tabularx}

\begin{tcolorbox}[colback=gray!10, colframe=gray!50, title=Nota sobre recuento de UEAs]
La cifra de 1\,598 UEAs mencionada en el Antecedente directo es la
\emph{suma} de UEAs por plan (contando cada UEA una vez por programa
en que aparece); el grafo unificado contiene 689 UEAs únicas.
El banco de datos \texttt{diccionario\_ueas.json} registra 1\,859 UEAs
(universo UAM-A completo). De las 689 obligatorias únicas, 639 cuentan
con datos completos de evaluación ordinaria y recuperación, y 606
disponen de las quince variables requeridas por el análisis factorial
de la \S\ref{sec:factores_uea}.
\end{tcolorbox}

\subsection*{Fases del análisis}

\textit{Fase~1 --- Extracción y modelado central (scripts 01--10, 30).}
El script~01 parsea los diez PDFs de planes de estudio, extrae las
relaciones de seriación y corregistro, y construye un
\texttt{networkx.DiGraph}~\cite{Hagberg2008} por plan con atributos de
nodo ($a_i$, créditos, nombre) y de arista (tipo de relación).
El script~02 implementa el modelo de cascada estática
(ecuación~\ref{eq:cascada}); los scripts~03 y~04 generan diagramas y
referencias analíticas.
El script~05 es el núcleo de simulación: Monte Carlo trimestre a trimestre
($S=20\,000$ trayectorias) con inscripción proporcional al grafo de
seriaciones, fragilidad latente compartida~(\ref{eq:fragilidad}) y
calibración por bisección sobre~$\ETtit^{\,\mathrm{ofic}}$.
Los scripts~06, 07 y~10 producen figuras, impacto marginal por arista
y curvas dosis-respuesta.
El script~30 cuantifica la incertidumbre de~$\ET{12}$ mediante IC~Monte
Carlo exacto (Clopper-Pearson) y análisis de sensibilidad a~$c_s$.

\textit{Fase~2 --- Análisis estadístico multivariado (scripts 08, 09, 11--19, 23R).}
Los scripts~08 y~09 realizan las correlaciones de Spearman a tres niveles
(plan, UEA, arista) con pruebas de permutación y producen las figuras
estadísticas.
Los scripts~11 y~12 ajustan modelos de saturación sobre las curvas
dosis-respuesta, clasifican los planes estratégicamente y producen las
figuras de portfolio.
Los scripts~13--17 calculan recuperabilidad, brecha de recuperación,
heterogeneidad docente, probabilidad de expulsión y el índice de riesgo
topológico enriquecido~$\rho_v$ (\S\ref{sec:riesgo_uea}).
El script~18 produce las estadísticas del cuerpo docente.
El script~19 profundiza con correlaciones parciales y análisis factorial
exploratorio (rotación oblimin~\cite{Horn1965}, $N=689$~UEAs, 5~factores).
El bloque de aditividad en \texttt{02\_remocion.R} (que consolida el original
\texttt{23R}; véase Apéndice~A) ajusta el modelo predictivo de~$\mathcal{A}(N)$
en R para verificación cruzada.

\textit{Fase~3 --- Extensiones cuantitativas (scripts 20--24, 21, 22, 23py).}
El script~20 descompone~$c_s$ en componente estructural y conductual;
el~21 incorpora varianza docente explícita en el modelo estocástico;
el~22 implementa el optimizador de portfolio (greedy, backtracking,
simulated annealing); el~23py ajusta el modelo OLS de subaditividad
en Python; el~24 genera las figuras de ajuste de saturación.

\textit{Fase~4 --- Índices de desempeño y prelación (scripts 31--36).}
Los scripts~31 y~33 simulan trayectorias de~$P$ y~$h$ bajo tres y cuatro
estrategias de inscripción respectivamente; los scripts~32, 34 y~36
producen las figuras del análisis de incentivos, la trampa ACT y la
comparación $P$~vs.~$h$.
El script~35 simula las cuatro estrategias bajo la lógica de~$h$ y
demuestra la extinción del incentivo al subregistro.

\textit{Fase~5 --- Complejidad sistémica (scripts 37--38).}
El script~37 construye el índice~$C$ sobre cinco dimensiones normalizadas
por min-max inter-plan, clasifica los planes en tres rangos y cuatro
arquetipos, y calcula la matriz de acoplamiento entre dimensiones.
El script~38 produce las cinco figuras del análisis de complejidad.

\subsection*{Infraestructura de software}

Los grafos de seriación se construyen y analizan con
\texttt{networkx}~\cite{Hagberg2008}.
Las simulaciones Monte Carlo utilizan \texttt{numpy} y \texttt{scipy}
para el modelo de fragilidad latente~$\theta_s\sim\mathcal{N}(0,\sigma^2)$.
Las figuras Python se producen con \texttt{matplotlib} y
\texttt{seaborn}; los análisis estadísticos multivariados y las figuras
de R con \texttt{ggplot2}~\cite{Wickham2016}, \texttt{dplyr},
\texttt{patchwork} y \texttt{psych}~\cite{Revelle2024}.
Los archivos de datos externos se leen con \texttt{python-dotenv}
(credenciales) y \texttt{pdfplumber}/\texttt{pdfminer} (extracción de PDFs).

\subsection*{Estructura de salidas}

Todos los artefactos se escriben en \texttt{Output/}:
las figuras de publicación en \texttt{Output/figuras/},
las figuras estadísticas en \texttt{Output/figuras\_correlaciones/},
los CSV de indicadores en \texttt{Output/R\_analisis/},
y los archivos de riesgo legal en \texttt{Output/legal/}.
Los grafos de seriación por plan se almacenan en \texttt{Grafos/} como
JSON con el esquema canónico de \texttt{networkx}.
El reporte \LaTeX\ (\texttt{06\_reporte\_v2.tex}) consume las figuras
directamente desde estas rutas mediante rutas relativas.

%══════════════════════════════════════════════════════
\section{Modelo de cascada estática}
\label{sec:cascada}
%══════════════════════════════════════════════════════

\subsection{Formulación}

El modelo de cascada estática calcula la fracción $\Pacc(v)$ de la cohorte
con acceso regular a la UEA~$v$ mediante la recursión
\begin{equation}
  \Pacc(v)
  \;=\; \prod_{u\to v}^{\!\text{seriación}} \pi_u
        \;\cdot\;
        \prod_{w}^{\!\text{corregistro}} \Pacc(w),
  \qquad
  \pi_v = \Pacc(v)\cdot p_v,
  \label{eq:cascada}
\end{equation}
donde $p_v$ es la tasa histórica de aprobación (sección de grupos de evaluación global ---GEG--- del diccionario)
de la UEA~$v$.
El primer producto de~(\ref{eq:cascada}) recorre los prerrequisitos directos~$u$
de~$v$ y usa $\pi_u$ (fracción que \emph{aprueba}); el segundo recorre las
raíces de corregistro~$w$ y usa $\Pacc(w)$ (fracción que llega, no que aprueba),
respetando la diferencia semántica entre ambas relaciones.

La eficiencia terminal queda acotada superiormente por
\begin{equation}
  \ET{t_{\min}} \;\leq\; \min_{v\in G}\,\Pacc(v),
  \label{eq:cota_ET}
\end{equation}
que es la cota $\ET{t_{\min}}\leq\min_{v\in G}\Pacc(v)$ del Reporte
TGNA-UAM~\cite{TGNA2026} extendida a los diez planes.

\subsection{Validación sobre el Tronco General}

El mapa \texttt{blocks\_v4.tex} registra
$\Pacc(\text{Métodos Numéricos}) = 4.68\,\text{\%}$.
La recursión~(\ref{eq:cascada}) con tasas GEG (grupos de evaluación global) sobre el plan Civil produce
\begin{equation}
  \Pacc(\uea{1151039}) = 0.0473,
  \label{eq:validacion}
\end{equation}
con un error del 1.1\% respecto al valor de referencia.

\begin{table}[H]\centering
\caption{Cascada estática del Tronco General actual.
         Tres pasos reducen la cohorte del 100\% al 13.4\%;
         lo que sigue es profundización sucesiva.}
\label{tab:cascada_tg}
\small
\begin{tabular}{llccc}
\toprule
Código & UEA & $p_v$ & $\Pacc(v)$ & $\pi_v$ \\
\midrule
\uea{1112042} & Introducción al Cálculo          & 0.505 & 1.000 & 0.505 \\
\uea{1112043} & Cálculo Diferencial              & 0.504 & 0.505 & 0.255 \\
\uea{1112029} & Cálculo Integral                 & 0.526 & 0.255 & 0.134 \\
\uea{1112030} & Ec.\ Diferenciales (corregistro) & 0.551 & 0.134 & 0.074 \\
\uea{1151039} & \textbf{Métodos Numéricos}       & 0.757 & \textbf{0.047} & 0.036 \\
\uea{1111081} & \textbf{Intro.\ Electrostática}  & 0.637 & \textbf{0.024} & 0.015 \\
\bottomrule
\end{tabular}
\end{table}

\subsection{Estructura de los grafos}

\begin{table}[H]\centering
\caption{Estadísticas de los grafos de seriación por plan.
         Todos los grafos son acíclicos.
         El total de~1\,203 aristas de seriación es la suma por plan; aristas
         compartidas entre planes (p.\,ej.\ las del Tronco General) se cuentan
         una vez por plan.
         De estas, 499~fueron evaluadas individualmente para impacto marginal
         sobre~$\ETtit$ (las restantes presentan impacto considerablemente menos relevante en el modelo
         y se excluyen del análisis de efecto; véase~\S\ref{sec:efecto}).}
\label{tab:grafos}
\small
\begin{tabular}{lcccc}
\toprule
Plan & UEAs oblig. & Seriaciones & Corregistros & Créditos \\
\midrule
Ambiental    & 53 & 108 & 17 & 411 \\
Civil        & 53 & 136 & 16 & 394 \\
Computación  & 52 & 157 & 22 & 406 \\
Eléctrica    & 57 &  94 & 26 & 410 \\
Electrónica  & 58 & 139 & 35 & 408 \\
Física       & 52 & 141 & 36 & 376 \\
Industrial   & 56 &  83 & 15 & 405 \\
Mecánica     & 54 & 129 & 29 & 379 \\
Metalurgia   & 58 & 129 & 29 & 419 \\
Química      & 56 &  98 & 25 & 411 \\
\midrule
\textbf{Total} & \textbf{551} & \textbf{1\,203} & \textbf{250} & --- \\
\bottomrule
\end{tabular}
\end{table}

Cada arista del grafo proviene de una entrada en la columna \textsc{seriación} del PDF del plan de estudios.
Un código sin prefijo (p.\,ej.\ \texttt{1112043}) genera una \emph{seriación directa}: la UEA prerrequisito debe estar aprobada antes de que la UEA pueda inscribirse.
Un código con prefijo \texttt{C} (p.\,ej.\ \texttt{C1112042}) genera un \emph{corregistro}: la UEA prerrequisito debe estar aprobada o inscribirse de forma simultánea en el mismo trimestre.
Las condiciones de umbral crediticio (p.\,ej.\ \texttt{300 Créditos}) establecen un mínimo de créditos acumulados como requisito de elegibilidad; no generan una seriación entre UEAs y no contribuyen al conteo de la Tabla~\ref{tab:grafos}.
Las entradas compuestas con el operador \texttt{y} generan tantas aristas como prerrequisitos listen; cada arista se contabiliza de forma independiente.

El conteo de la Tabla~\ref{tab:grafos} es por plan, de modo que el total divisional de 1\,203 seriaciones y 250 corregistros es la suma de los totales individuales y no el recuento de aristas únicas del grafo unificado.
De las 1\,203 seriaciones, 499 presentan un impacto marginal distinguible de cero sobre~$\ETtit$ y fueron evaluadas individualmente; las 704 restantes se excluyen del análisis de efecto (\S\ref{sec:efecto}).

La UEA de mayor alcance es Ecuaciones Diferenciales Ordinarias (\uea{1112030}):
19~aristas de salida y 380~descendientes, obligatoria en los diez planes.

%══════════════════════════════════════════════════════
\section{Modelo topológico-estocástico}
\label{sec:modelo}
%══════════════════════════════════════════════════════

La cascada estática da cotas analíticas, pero no responde cuánto tarda
en egresar quien logra completar el plan ni qué fracción de la matrícula
egresará eventualmente.
Para eso se construyó una simulación Monte Carlo trimestre por trimestre.

\subsection{Dinámica de inscripción y resultado}

El modelo inicializa $S$~estudiantes con el conjunto de UEAs obligatorias del
plan pendientes.
En cada trimestre, cada estudiante intenta inscribir en orden topológico las
UEAs elegibles que quepan en su carga $c_s$~cr/trim.

Una UEA es elegible cuando: (a) todos sus prerrequisitos están aprobados;
(b) todas sus raíces de corregistro están aprobadas o se inscriben
simultáneamente este trimestre; y (c) cualquier umbral de créditos acumulados
ya se ha alcanzado.

El resultado de la UEA inscrita se sortea con
\begin{equation}
  U_{s,i}\sim\mathrm{Uniforme}(0,1),
  \label{eq:sorteo}
\end{equation}
partiendo el intervalo en tres regiones:
\begin{align}
  U_{s,i} < r_i
    &\;\Longrightarrow\; \text{renuncia (baja semana~5, sin NA),}
    \label{eq:renuncia}\\
  r_i \leq U_{s,i} < r_i + (1-r_i)\,a_{s,i}
    &\;\Longrightarrow\; \text{aprueba,}
    \label{eq:aprueba}\\
  U_{s,i} \geq r_i + (1-r_i)\,a_{s,i}
    &\;\Longrightarrow\; \text{reprueba (NA),}
    \label{eq:reproba}
\end{align}
donde $r_i$ es la tasa histórica de baja en semana~5
(\texttt{Promedio\_RENUNCIAS}) y $a_{s,i}$ es la tasa efectiva del
estudiante~$s$, descrita a continuación.
La renuncia no consume oportunidad; la reprobación sí.
Al acumular cinco evaluaciones No Acreditadas (NA) en una UEA el estudiante causa baja del programa
(reglamento UAM).

\subsection{Fragilidad latente compartida}
\label{sec:fragilidad}

Un estudiante con dificultades generales tiende a reprobar varias materias
en el mismo trimestre, no una aleatoriamente: sus resultados en distintas
UEAs están correlacionados a través de su habilidad personal.
Para capturar esta correlación intra-estudiante se introduce una variable
de habilidad latente $\theta_s\sim\mathcal{N}(0,1)$, sorteada una sola vez
por estudiante al inicio de cada trimestre, y se define la tasa de
aprobación individualizada
\begin{equation}
  a_{s,i}
  = \Phi\!\left(
      \Phi^{-1}(a_i)\sqrt{1+\sigma^2} \;+\; \sigma\,\theta_s
    \right),
  \label{eq:fragilidad}
\end{equation}
donde $\Phi:\mathbb{R}\to(0,1)$ es la función de distribución acumulada
de la normal estándar, $\Phi^{-1}$ su inversa (función probit),
$a_i\in(0,1)$ es la tasa histórica de aprobación de la UEA~$i$, y
$\sigma\geq 0$ es el \textbf{parámetro de fragilidad latente}, que
controla la intensidad de la correlación entre UEAs dentro del mismo
trimestre: a mayor $\sigma$, el desempeño del estudiante en todas las
materias del trimestre queda más determinado por su habilidad personal
$\theta_s$ y menos por la dificultad específica de cada UEA.

La ecuación~(\ref{eq:fragilidad}) opera en el espacio probit: la
dificultad de la UEA~$i$ se representa por el nivel
$\mu_i = \Phi^{-1}(a_i)\sqrt{1+\sigma^2}$, y la habilidad del
estudiante entra como un desplazamiento aditivo $\sigma\,\theta_s$
sobre ese nivel.
El prefactor $\sqrt{1+\sigma^2}$ es el rescalado exigido por la
identidad probit-normal
\begin{equation}
  \mathbb{E}_{\theta_s}\!\left[
    \Phi\!\left(\mu + \sigma\,\theta_s\right)
  \right]
  = \Phi\!\left(\frac{\mu}{\sqrt{1+\sigma^2}}\right),
  \label{eq:probit_normal}
\end{equation}
que garantiza $\mathbb{E}[a_{s,i}] = a_i$ para todo $\sigma$:
sustituyendo $\mu = \Phi^{-1}(a_i)\sqrt{1+\sigma^2}$
en~(\ref{eq:probit_normal}) se recupera $\Phi(\Phi^{-1}(a_i)) = a_i$.
El modelo es, por construcción, consistente con los datos históricos
en media sin importar el valor de~$\sigma$.

El caso $\sigma=0$ reduce a $a_{s,i}=a_i$ para todos los estudiantes
(sin correlación entre UEAs), mientras que $\sigma\to\infty$ lleva
$a_{s,i}\to\mathbf{1}[\theta_s>0]$, con el resultado determinado
exclusivamente por la habilidad latente (correlación perfecta).
La dispersión de referencia $\sigma=0.6$ sitúa al modelo entre ambos
extremos; su influencia sobre las métricas de egreso se examina en
la Fig.~\ref{fig:sensibilidad}.

\subsection{Calibración}
\label{sec:calibracion}

El único parámetro libre es $\ell$ (\texttt{lento\_div}), el divisor de carga:
cada estudiante inscribe $c_s = \mathrm{cr\_total}/\ell$~cr/trim.
El valor resultante es la \emph{carga efectiva histórica} que reproduce
$\ETtit^{\,\mathrm{ofic}}$; no es la carga normal del plan (45~cr/trim) ni su máximo (63~cr/trim).
$\ell$ se determina por búsqueda binaria imponiendo que
\begin{equation}
  \ETtit^{\,\mathrm{mod}}(p) = \ETtit^{\,\mathrm{ofic}}(p)
  \qquad\forall\;p\in\{10\;\text{planes}\},
  \label{eq:calibracion}
\end{equation}
donde el lado derecho de~(\ref{eq:calibracion}) es la Tabla~\ref{tab:et_oficial}.

\begin{notabox}
$\ET{12}$ no es ancla de calibración: emerge libremente de la simulación y
se coteja con el 1.52\% de la Tabla~\ref{tab:historico} como verificación
independiente.
\end{notabox}

La Tabla~\ref{tab:calibracion} muestra los resultados de la calibración.
En ocho de los diez planes, el error en la condición~(\ref{eq:calibracion})
es $|\Delta|\leq 0.52$~trim (aproximadamente siete semanas académicas, menos de medio trimestre); el promedio
ponderado de CBI es 20.62~trim frente a 20.60~trim oficiales.

\begin{table}[H]\centering
\caption{Parámetros calibrados.
         $c_s$~es la carga efectiva histórica (no la normal de 45~cr/trim ni el máximo de 63~cr/trim).
         $\Delta = \ETtit^{\,\mathrm{mod}} - \ETtit^{\,\mathrm{ofic}}$.}
\label{tab:calibracion}
\small
\begin{tabular}{lrrrrrc}
\toprule
Plan & cr.\ total & $\ell$ & $c_s$ (cr/trim) &
      $\ETtit^{\,\mathrm{mod}}$ & $\ETtit^{\,\mathrm{ofic}}$ & $\Delta$ \\
\midrule
amb &  411 & 13.70 & 30.0 & 20.48 & 20.56 & $-0.08$ \\
civ &  394 & 13.13 & 30.0 & 20.67 & 20.34 & $+0.33$ \\
com &  406 & 13.10 & 31.0 & 20.39 & 20.31 & $+0.08$ \\
ele &  410 & 13.67 & 30.0 & 20.89 & 21.41 & $-0.52$ \\
elo &  408 & 13.60 & 30.0 & 23.54 & 23.24 & $+0.30$ \\
fis &  376 & 12.53 & 30.0 & 20.61 & 20.25 & $+0.36$ \\
ind &  405 & 13.06 & 31.0 & 19.49 & 19.40 & $+0.09$ \\
mec &  379 & 14.04 & 27.0 & 20.94 & 21.18 & $-0.24$ \\
met$^\dagger$ & 419 & 10.22 & 41.0 &
      \textit{16.25}$^{*}$ & 22.26 & --- \\
qui &  411 & 12.45 & 33.0 & 19.48 & 19.41 & $+0.07$ \\
\midrule
\textbf{CBI} & & & & \textbf{20.62} & \textbf{20.60} & $-0.02$ \\
\bottomrule
\multicolumn{7}{l}{\footnotesize
  $^\dagger$ Caso especial; véase §\ref{sec:met}.\enspace
  $^{*}$ Cota inferior; no comparable con el oficial.}
\end{tabular}
\end{table}

\begin{figure}[H]\centering
\includegraphics[width=\linewidth]{../../Output/figuras/calibracion.pdf}
\caption{Calibración del modelo probabilístico de trayectorias estudiantiles.
  Aquí $\ETtit = \mathbb{E}[T\,|\,\mathrm{tit}]$ denota el tiempo promedio de
  egreso (en trimestres) condicionado a que el estudiante eventualmente egrese;
  el superíndice \emph{mod} corresponde al valor simulado y \emph{ofic} al
  promedio oficial registrado en UAM-A (2018--2025).
  \emph{Izquierda}: $\ETtit^{\,\mathrm{mod}}$ frente a $\ETtit^{\,\mathrm{ofic}}$
  para los diez planes de CBI; la línea punteada es la identidad ($y=x$) y la
  franja gris indica la tolerancia de~$\pm 1$~trimestre.
  El símbolo triangular corresponde a Metalurgia, cuyo valor modelado es una cota
  inferior (la UEA de estancia industrial de 40~créditos impone restricciones de
  disponibilidad que el modelo no captura completamente).
  \emph{Derecha}: error de calibración $\Delta = \ETtit^{\,\mathrm{mod}} - \ETtit^{\,\mathrm{ofic}}$
  para los nueve planes comparables (excluida Metalurgia como caso especial).
  \textbf{El modelo reproduce el tiempo oficial de egreso con un error máximo de
  menos de medio trimestre en todos los planes comparables (véase tabla adjunta),
  lo que valida su uso como herramienta de diagnóstico y simulación.}}
\label{fig:calibracion}
\end{figure}

\subsection{Ingeniería Metalúrgica: el acantilado de factibilidad}
\label{sec:met}

La UEA Trabajo en Planta Metalúrgica (\uea{1145079}) tiene 40~créditos.
Con el techo estándar de búsqueda, el modelo calibraba
$c_s \approx 34.9$~cr/trim, de modo que ningún estudiante podía inscribir
esa UEA jamás, y la tasa de egreso colapsaba a cero.
La corrección implementada acota la búsqueda para garantizar
$c_s\geq 41$~cr/trim, produciendo $\ell=10.22$ y los resultados de la
Tabla~\ref{tab:calibracion}.

El valor $\ETtit^{\,\mathrm{mod}}=16.25$~trim es una cota inferior.
La diferencia con los 22.26~trim oficiales refleja restricciones de
disponibilidad de la estancia industrial que el modelo no captura.
Metalurgia se reporta con esta aclaración y no se incluye en el promedio
ponderado de CBI.

%══════════════════════════════════════════════════════
\section{Resultados}
\label{sec:resultados}
%══════════════════════════════════════════════════════

\subsection{Eficiencia terminal al plazo del plan de estudios}

El modelo se corrió con la matrícula real por plan del trimestre~25-O
($N=6\,562$~estudiantes), $\sigma=0.6$ y los parámetros de la
Tabla~\ref{tab:calibracion}.

La eficiencia terminal a 12~trimestres es estadísticamente próxima a cero
en los diez planes bajo los parámetros centrales de calibración.

Este es el resultado calibrado sobre datos históricos reales, donde
$\ETtit^{\,\mathrm{mod}}$ reproduce el promedio oficial con error menor a
medio trimestre.
El análisis de los mecanismos se presenta en la Sección~\ref{sec:mecanismos}.

\begin{figure}[H]\centering
\includegraphics[width=\linewidth]{../../Output/figuras/curvas_ET.pdf}
\caption{Curvas de egreso acumulado $\ET{k}$ (fracción de la cohorte que
  egresa en a lo sumo $k$ trimestres) para los diez planes de CBI.
  Simulación con la matrícula real del trimestre~25-O ($N=6\,562$ estudiantes)
  y parámetro de fragilidad latente $\sigma=0.6$ (donde $\sigma$ modela que
  los estudiantes con evaluaciones No Acreditadas acumuladas tienen menor
  probabilidad de reinscribir).
  La línea vertical roja indica $t=12$ (plazo del plan):
  $\ET{12}$ estadísticamente próxima a cero en los diez planes
  bajo parámetros centrales de calibración.
  A $t=18$ la fracción egresada va del 2.1\% (Electrónica) al 29.5\%
  (Industrial) y al 52.3\% (Metalurgia).
  A $t=24$ (el doble del plazo establecido) la fracción egresada oscila entre
  el 35.9\% (Electrónica) y el 63.3\% (Industrial).
  La curva punteada de Metalurgia es una cota inferior: la UEA de estancia
  industrial de 40~créditos impone restricciones de disponibilidad no capturadas
  por el modelo general.
  \textbf{Ningún plan de CBI logra que un porcentaje apreciable de su cohorte
  egrese en el plazo de doce trimestres del plan de estudios; la brecha entre el
  plazo del plan y el tiempo real es estructural, no accidental.}}
\label{fig:curvas_ET}
\end{figure}

La Figura~\ref{fig:curvas_ET} muestra que las curvas $\ET{t}$ arrancan en
cero a $t=12$ y crecen de manera continua pero lenta.
Ningún plan alcanza el 50\% de egreso antes del trimestre~18, excepto
Metalurgia.
Electrónica no despega hasta $t\approx 17$--$18$ y su curva permanece
claramente por debajo del resto a lo largo de todo el horizonte.
Industrial y Química presentan los arranques más tempranos.

\subsection{Tasas de egreso y baja}

\begin{figure}[H]\centering
\includegraphics[width=\linewidth]{../../Output/figuras/resumen_planes.pdf}
\caption{Desenlaces simulados de la cohorte activa para los diez planes de CBI.
  Las barras muestran, para cada plan, la fracción de estudiantes que: egresa
  dentro de 30~trimestres (barra teal), causa baja definitiva por acumular
  cinco evaluaciones No Acreditadas en alguna UEA (barra roja), o permanece
  sin egresar ni dar baja al llegar al trimestre~30 (barra gris).
  Simulación sobre la matrícula real del trimestre~25-O ($N=6\,562$) con
  parámetro de fragilidad latente $\sigma=0.6$.
  \textbf{La tasa de egreso oscila entre el 57.5\% (Electrónica) y el
  67.2\% (Industrial), mientras que la tasa de baja por reprobación ronda
  el 32--40\% en todos los planes; las seriaciones retrasan a quienes
  egresan pero no son la causa principal del abandono.}}
\label{fig:resumen}
\end{figure}

El modelo predice que el \textbf{62.5\%} de los 6\,562~estudiantes activos
egresará dentro del horizonte de 30~trimestres (alrededor de 4\,100~personas),
mientras el \textbf{36.3\%} causará baja por acumulación de reprobaciones
(aproximadamente 2\,380~personas).
Menos del 1.2\% permanece sin egresar ni dar baja al llegar al trimestre~30.

La Figura~\ref{fig:resumen} muestra que la tasa de egreso oscila entre el
57.5\% (Electrónica) y el 67.2\% (Industrial).
La tasa de baja es relativamente estable (32--40\%) y no correlaciona de
forma sencilla con la densidad de seriaciones: planes con muchas seriaciones
pueden tener bajas moderadas si la carga está bien distribuida.

\subsection{Distribución del tiempo de egreso}

\begin{figure}[H]\centering
\includegraphics[width=\linewidth]{../../Output/figuras/distribucion_T.pdf}
\caption{Distribución del tiempo de egreso $T$ (en trimestres) por plan.
  Para cada programa, las barras muestran el rango intercuantílico:
  barras claras = intervalo $[p_{10},\,p_{90}]$ (percentiles 10 y 90);
  barras oscuras = intervalo $[p_{10},\,p_{50}]$; la línea gruesa es la mediana
  $p_{50}$; el rombo es $\ETtit^{\,\mathrm{mod}} = \mathbb{E}[T\,|\,\mathrm{tit}]$
  (tiempo promedio de egreso simulado condicionado al egreso).
  Líneas de referencia verticales: $t=12$ (rojo, plazo del plan) y
  $t=24$ (gris, el doble del plazo establecido).
  \textbf{Las medianas de egreso van de 16~trimestres (Metalurgia)
  a 23~trimestres (Electrónica); el rango $[p_{10},p_{90}]$ es de 8--13
  trimestres dentro de cada plan, indicando que las cadenas de seriación
  afectan con mayor intensidad a los estudiantes de trayectoria más lenta.}}
\label{fig:distribucion}
\end{figure}

Las medianas de egreso (Figura~\ref{fig:distribucion}) van de 19~trim
(Industrial, Química) a 23~trim (Electrónica).
El rango $[p_{10},\,p_{90}]$ es de 8--9~trimestres en cada plan: hay
estudiantes que egresan desde el trimestre~15 y otros que tardan hasta
el~28--30.
Electrónica y Mecánica presentan las colas más pesadas ($p_{90}=28$ y
$26$~trim respectivamente), lo que indica que sus cadenas de seriación
afectan con mayor intensidad a los estudiantes del tramo lento.

\subsection{Robustez al parámetro de fragilidad latente}

\begin{figure}[H]\centering
\includegraphics[width=\linewidth]{../../Output/figuras/sensibilidad_sigma.pdf}
\caption{Sensibilidad de tres métricas de egreso al parámetro de fragilidad
  latente~$\sigma$, para el Plan Civil de referencia con carga efectiva
  histórica $c_s=30$~créditos/trimestre.
  El parámetro~$\sigma$ modela que los estudiantes con evaluaciones No
  Acreditadas acumuladas tienen menor probabilidad de reinscribir: a mayor
  $\sigma$, mayor penalización por historial de reprobación.
  Las métricas mostradas son $\ET{12}$ (fracción de la cohorte que egresa
  en a lo sumo 12~trimestres), la tasa de egreso a 30~trimestres, y
  $\ETtit = \mathbb{E}[T\,|\,\mathrm{tit}]$ (tiempo promedio de egreso
  condicionado al egreso).
  La línea punteada vertical indica $\sigma=0.6$ (valor de referencia
  calibrado sobre datos históricos UAM-A).
  \textbf{El resultado central --- $\ET{12} \approx 0$ --- es robusto para
  todo $\sigma \leq 0.9$, lo que confirma que la imposibilidad de egresar
  en el plazo del plan no es un artefacto del parámetro de fragilidad sino
  una consecuencia de la estructura del plan.}}
\label{fig:sensibilidad}
\end{figure}

La Figura~\ref{fig:sensibilidad} confirma que $\ET{12}$ es cercano a cero para
todo $\sigma\leq 0.9$ y que los resultados cualitativos son robustos para
$\sigma\in[0.3,\,1.2]$.
El resultado principal no depende de la elección exacta del parámetro de
fragilidad latente.

%══════════════════════════════════════════════════════
\section{Los dos mecanismos del desajuste curricular}
\label{sec:mecanismos}
%══════════════════════════════════════════════════════

El desajuste entre el currículo diseñado y la población real se expresa a través
de dos mecanismos distintos, con pesos distintos según el plan.
La distinción importa para el diseño de medidas correctivas: el primero opera
con independencia de la topología de seriaciones; el segundo opera como cota
secundaria en todos los planes, con efectos observables más pronunciados en
Electrónica (mayor profundidad topológica del DAG y mayor ganancia acumulada
por remoción de seriaciones) y en Metalurgia (cuya cota de carga es la más
cercana al umbral de 12~trimestres, permitiendo que la remoción de seriaciones
mueva~$\ET{12}$ de cero a $10.4\,\text{\%}$).

\subsection{Mecanismo~1: el desajuste de carga (dominante en 8 planes)}

Los Planes~2020 tienen una duración establecida de~12~trimestres y un total de
376--419~créditos obligatorios según el programa; dividiendo ambas cantidades se
obtiene la carga efectiva de diseño implícita: $\approx 45$~cr/trim, con un máximo
establecido de~63~cr/trim (sección~5 de cada plan; valores uniformes en los
diez programas CBI).
Ese parámetro implica que un estudiante que inscribe la carga de diseño con
aprobación perfecta completaría los créditos en~$\approx 9$~trimestres, con un
margen de tres trimestres para eventualidades dentro del plazo establecido.

La calibración del modelo a los datos históricos revela que la población real
opera a $c_s\approx 27$--$33$~cr/trim ---entre el 60\% y el 73\% de la
carga efectiva de diseño---.
Esa brecha es el dato central del diagnóstico: el currículo fue diseñado
asumiendo un ritmo de inscripción que la evidencia histórica muestra como
minoritario en la cohorte real.
Con la carga calibrada y una tasa neta de aprobación $\bar{a}\approx 0.745$--$0.799$,
los créditos aprobados acumulados en 12~trimestres son
\begin{equation}
  c_s\cdot\bar{a}\cdot 12 \;\approx\; 220\text{--}283\;\text{cr},
  \label{eq:acum_12}
\end{equation}
muy por debajo de los 376--419~requeridos.

La expresión~(\ref{eq:acum_12}) es condición suficiente para que $\ET{12}$
sea estadísticamente próxima a cero: incluso sin ninguna seriación, el
estudiante no puede completar el plan en 12~trimestres a la carga histórica.
Los datos de la Tabla~\ref{tab:mecanismos} lo confirman: en ocho planes, el
mínimo de acumulación crediticia
$\lceil\mathrm{cr\_total}/(c_s\cdot\bar{a})\rceil$ va de 18 a 21~trimestres,
muy lejos del umbral de 12.

\subsection{Mecanismo~2: la profundidad topológica como cota secundaria}

Las cadenas de seriación imponen un orden secuencial irreductible.
La profundidad máxima del grafo acíclico dirigido (DAG) de seriaciones ---el número mínimo de cursos
que deben completarse uno detrás del otro antes de poder inscribir el curso
terminal--- es la cota topológica pura sobre el tiempo de egreso:
\begin{equation}
  T_{\min}^{\,\mathrm{topo}} = \mathrm{depth}_{\max} + 1.
  \label{eq:tmin_topo}
\end{equation}

La Tabla~\ref{tab:mecanismos} muestra los valores de~(\ref{eq:tmin_topo})
por plan, obtenidos directamente de la búsqueda en profundidad (DFS, \emph{depth-first search}) sobre el grafo de seriaciones.
La profundidad máxima varía de 5~niveles (Ind) a 9~niveles (Com, Elo),
con mínimos topológicos de 6 a 10~trimestres.
En ningún plan la topología sola alcanza la cota de 12~trimestres:
el valor máximo, $T_{\min}^{\,\mathrm{topo}}=10$ (Com y Elo), se mantiene
por debajo del umbral del plan, y en todos los casos la cota de carga
domina ampliamente ($T_{\min}^{\,\mathrm{carga}}=15$--$21$~trim).
La restricción topológica actúa, por tanto, como cota secundaria que
limita la densidad de inscripción en los primeros trimestres, pero no
constituye el cuello de botella determinante del tiempo de egreso.

\begin{table}[H]\centering
\caption{Las dos cotas por plan.
         $\bar{a}=\overline{A_i(1-R_i)}$ (tasa neta de aprobación media).
         $T_{\min}^{\,\mathrm{topo}}=\mathrm{depth}_{\max}+1$.
         $T_{\min}^{\,\mathrm{carga}}=\lceil\mathrm{cr\_total}/(c_s\cdot\bar{a})\rceil$.
         El mecanismo dominante (mayor cota) aparece en \textbf{negrita}.}
\label{tab:mecanismos}
\small
\setlength{\tabcolsep}{5pt}
\begin{tabular}{lrrrr@{\quad}rr@{\quad}l}
\toprule
Plan &
  $c_s$ & $\bar{a}$ & cr & depth &
  $T_{\min}^{\,\mathrm{topo}}$ &
  $T_{\min}^{\,\mathrm{carga}}$ & Dominante \\
\midrule
amb &  30 & 0.735 & 411 &  8 &  9 & \textbf{19} & carga \\
civ &  30 & 0.681 & 394 &  7 &  8 & \textbf{20} & carga \\
com &  31 & 0.656 & 406 &  9 & 10 & \textbf{20} & carga \\
ele &  30 & 0.690 & 410 &  8 &  9 & \textbf{20} & carga \\
elo &  30 & 0.670 & 408 &  9 & 10 & \textbf{21} & carga \\
fis &  30 & 0.675 & 376 &  7 &  8 & \textbf{19} & carga \\
ind &  31 & 0.727 & 405 &  5 &  6 & \textbf{18} & carga \\
mec &  27 & 0.683 & 379 &  8 &  9 & \textbf{21} & carga \\
met &  41 & 0.726 & 419 &  8 &  9 & \textbf{15} & carga \\
qui &  33 & 0.716 & 411 &  8 &  9 & \textbf{18} & carga \\
\bottomrule
\end{tabular}
\end{table}

\begin{figure}[H]\centering
\includegraphics[width=\linewidth]{../../Output/figuras/profundidad_planes.pdf}
\caption{Visualización de las dos cotas estructurales que hacen imposible
  egresar en el plazo establecido de 12~trimestres.
  \emph{Izquierda}: profundidad máxima del DAG (grafo acíclico dirigido) de
  seriaciones por plan; la línea punteada marca profundidad~$=9$, valor máximo
  observado (Computación y Electrónica), con mínimo topológico
  $T_{\min}^{\,\mathrm{topo}}=10$~trimestres.
  Ningún plan alcanza la cota de 12~trimestres por restricción topológica.
  \emph{Derecha}: cota de carga mínima $T_{\min}^{\,\mathrm{carga}}$
  (mínimo de trimestres requeridos dado la carga efectiva histórica~$c_s$
  en créditos/trimestre y la tasa de aprobación media del plan) por plan,
  con línea de referencia en $t=12$.
  \textbf{Todos los planes superan los 15~trimestres de carga mínima,
  confirmando que la restricción de créditos es la cota dominante en los
  diez programas; la topología actúa como cota secundaria sin alcanzar
  por sí sola el umbral de 12~trimestres.}}
\label{fig:profundidad}
\end{figure}

\subsection{Metalurgia: el caso confirmatorio}

Metalurgia es el experimento natural que separa los dos mecanismos.
Cuando se eliminan \emph{todas} las seriaciones, $\ET{12}$ sube de 0\% a
\textbf{10.4\%}.
Esto ocurre porque la cota de carga para Metalurgia es sólo 15~trimestres
(gracias a $c_s=41$~cr/trim), sensiblemente más cercana al umbral de 12 que
la de cualquier otro plan (19--24~trim).
Al eliminar las seriaciones desaparece el orden secuencial que impide
inscribir las UEAs de mayor densidad crediticia en los primeros trimestres;
sin esa restricción, los estudiantes con alta tasa de aprobación pueden
aproximar la carga efectiva a $c_s=41$~cr/trim desde el inicio, acortando
el tiempo de egreso lo suficiente para que un~10.4\% alcance los 12~trimestres.

En los otros nueve planes, eliminar todas las seriaciones mantiene $\ET{12}$
estadísticamente próxima a cero
porque la cota de carga (18--21~trim) sigue muy por encima de 12.

Este resultado distingue con precisión la contribución de cada mecanismo.

\subsection{Descomposición de la carga efectiva histórica}
\label{sec:descomposicion_carga}

El parámetro calibrado $c_s \approx 30$~cr/trim (carga efectiva histórica de
inscripción) colapsa en un solo escalar tres fuentes de restricción cualitativamente
distintas: las restricciones topológicas impuestas por el grafo acíclico dirigido
(DAG) de seriaciones, las decisiones voluntarias de carga reducida por parte del
estudiante y las restricciones operativas de cupo u oferta trimestral.
Para el diagnóstico es relevante distinguir qué fracción del retraso es
estructural-curricular ---y por tanto directamente modificable mediante reforma del
plan--- y qué fracción es operativa o conductual.

El módulo~20 (\texttt{20\_descomposicion\_carga.py}) realiza esta descomposición
mediante la identidad
\begin{equation}
  c_s \;=\; c_s^{\mathrm{estr}} + c_s^{\mathrm{cond}},
  \label{eq:descomposicion_cs}
\end{equation}
donde $c_s^{\mathrm{estr}}$ es la carga máxima que la topología del plan permite
inscribir en cada trimestre bajo acceso libre ---estimada por simulación con
aprobación perfecta y sin tope de carga--- y $c_s^{\mathrm{cond}} = c_s - c_s^{\mathrm{estr}}$
es el residual atribuible a decisiones conductuales y restricciones operativas.
El módulo corre el modelo Monte Carlo bajo tres escenarios comparados con el
\emph{baseline} calibrado: (i)~solo restricciones topológicas ($c_s^{\mathrm{estr}}$),
(ii)~carga normal del plan (45~cr/trim, sin tope conductual) y (iii)~condición
histórica de referencia.

La conclusión central es que $\ET{12}$ (eficiencia terminal a 12 trimestres)
permanece estadísticamente próxima a cero incluso en el escenario de carga libre
en nueve de los diez planes.
Esto confirma que la restricción dominante es la carga crediticia total requerida
---independientemente de cómo se distribuye entre componentes--- y no las
decisiones conductuales individuales de los estudiantes.
La figura~\ref{fig:descomposicion_carga} muestra, plan por plan, las fracciones
atribuibles a cada componente y cuantifica cuánto del retraso es, en principio,
recuperable mediante reforma curricular frente a cuánto depende de condiciones
externas al plan de estudios.

\begin{figure}[htbp]
  \centering
  \includegraphics[width=\textwidth]{../../Output/figuras/fig_descomposicion_carga.pdf}
  \caption{Descomposición de la carga efectiva de inscripción~$c_s$ en sus
    componentes topológica~$c_s^{\rm estr}$ y conductual--operativa~$c_s^{\rm cond}$
    para los diez planes de estudio 2020.
    Panel izquierdo: barras apiladas por plan con el techo topológico~$c_s^{\rm estr}$
    marcado con diamante negro; la fracción en rojo corresponde al exceso
    conductual~$c_s^{\rm cond}$, con mayor magnitud en Metalurgia.
    Panel derecho: dispersión de~$c_s$ frente a~$c_s^{\rm estr}$; en nueve de los
    diez planes el techo topológico explica la carga observada sin residuo
    conductual apreciable.}
  \label{fig:descomposicion_carga}
\end{figure}

%══════════════════════════════════════════════════════
\section{Efecto marginal de las seriaciones}
\label{sec:efecto}
%══════════════════════════════════════════════════════

Establecido que las seriaciones no son los responsables primarios del bajo
valor de~$\ET{12}$,
la pregunta práctica es: ¿cuánto tiempo de egreso cuesta cada seriación?
¿Qué gana la División si las elimina?

El módulo~07 responde esto ejecutando, para cada uno de los diez planes y
cada arista del grafo, una simulación base seguida de una simulación con
esa arista eliminada.
La métrica es el cambio en el tiempo promedio de egreso,
\begin{equation}
  \Delta\ETtit
  = \ETtit^{\,\text{(sin arista)}} - \ETtit^{\,\text{(base)}},
  \label{eq:delta_et}
\end{equation}
donde $\Delta\ETtit<0$ indica que eliminar la seriación acorta el egreso.

\subsection{Escenarios globales}

\begin{figure}[H]\centering
\includegraphics[width=\linewidth]{../../Output/figuras/efecto_escenarios_todos.pdf}
\caption{Ganancia de egreso al eliminar las seriaciones de cada plan,
  para tres escenarios de remoción total.
  \emph{Panel izquierdo}: cambio $\Delta\ETtit$ en el tiempo promedio de
  egreso condicionado al egreso $\ETtit = \mathbb{E}[T\,|\,\mathrm{tit}]$
  (en trimestres); valores negativos indican que eliminar las seriaciones
  acorta el tiempo de egreso.
  \emph{Panel derecho}: cambio en la tasa de egreso (puntos porcentuales).
  Los tres escenarios son: remoción de todas las seriaciones del plan,
  remoción solo de las seriaciones del Tronco General, y remoción solo de
  las seriaciones de tronco específico.
  \textbf{Eliminar todas las seriaciones de un plan reduce $\ETtit$ entre
  0.2~trimestres (Computación) y 2.3~trimestres (Electrónica), pero apenas
  modifica la tasa de egreso ($<3$~pp en todos los planes), lo que confirma
  que las seriaciones retrasan el egreso sin ser la causa principal del abandono.}}
\label{fig:efecto_escenarios}
\end{figure}

La Figura~\ref{fig:efecto_escenarios} resume las ganancias de los escenarios
de remoción total.
Eliminar \emph{todas} las codependencias de un plan reduciría $\ETtit$ entre
0.2~trim (Computación) y 2.3~trim (Electrónica), con una ganancia media
ponderada de aproximadamente 0.8~trim en CBI.

Nótese que la ganancia en tasa de egreso es modesta: en ningún plan supera
los 3~puntos porcentuales.
Las seriaciones retrasan a quienes egresan; no son la causa principal de que
un 36\% cause baja.

\begin{notabox}
$\ET{12}$ permanece en 0\% al eliminar todas las codependencias en nueve
de los diez planes, confirmando que la cota de carga
($T_{\min}^{\,\mathrm{carga}} = 18$--$21$~trim) es la restricción activa.
Metalurgia confirma la excepción: sin seriaciones, $\ET{12} = 10.4\,\text{\%}$
porque su cota de carga (15~trim) queda suficientemente cerca del umbral de
12~trim.
\end{notabox}

\subsection{Las aristas más costosas}

\begin{figure}[H]\centering
\includegraphics[width=\linewidth]{../../Output/figuras/efecto_elo.pdf}
\caption{Plan Electrónica (elo): impacto individual de las seriaciones sobre
  el tiempo de egreso.
  \emph{Panel izquierdo}: las doce aristas del DAG (grafo acíclico dirigido
  de seriaciones) con mayor valor absoluto de $|\Delta\ETtit|$, donde
  $\Delta\ETtit = \ETtit^{\,\text{(sin arista)}} - \ETtit^{\,\text{(base)}}$
  es el cambio en el tiempo promedio de egreso (trimestres) al eliminar esa
  seriación; cada barra está etiquetada con el valor en trimestres y el cambio
  en la tasa de egreso en puntos porcentuales.
  Barras teal = eliminar la seriación acorta el egreso ($\Delta\ETtit < 0$);
  barras rojas = eliminarla lo alarga.
  \emph{Panel derecho}: resultados de los escenarios de remoción global del plan.
  \textbf{La seriación de Ecuaciones Diferenciales Ordinarias (EDO) $\to$ Circuitos Eléctricos~I lidera la División con
  $\Delta\ETtit = -1.07$~trimestres; ningún otro plan concentra un impacto
  marginal tan grande en una sola arista.}}
\label{fig:efecto_elo}
\end{figure}

Electrónica concentra el mayor impacto marginal de la División.
La arista más costosa (Figura~\ref{fig:efecto_elo}) es la seriación de
Ecuaciones Diferenciales Ordinarias (EDO, \uea{1112030}) $\to$ Circuitos Eléctricos~I (\uea{1124001}):
$\Delta\ETtit = -1.07$~trim y $\Delta\mathrm{grad} = +2.1$~pp.
El par teoría--laboratorio de Comunicaciones Analógicas ocupa el segundo y
tercer lugar.

\begin{figure}[H]\centering
\includegraphics[width=\linewidth]{../../Output/figuras/efecto_qui.pdf}
\caption{Plan Química (qui): impacto individual de cada seriación sobre el
  tiempo promedio de egreso $\ETtit = \mathbb{E}[T\,|\,\mathrm{tit}]$.
  Cada barra representa una arista del DAG de seriaciones; su longitud es
  $|\Delta\ETtit|$ (trimestres ahorrados al eliminar esa seriación, donde
  valores negativos indican acortamiento del egreso).
  Las barras están ordenadas por impacto descendente y coloreadas por subsecuencia
  temática.
  La cadena Termodinámica $\to$ Termodinámica~Aplicada $\to$ Balance de Materia
  $\to$ Balance de Energía $\to$ Procesos de Separación~I agrupa las aristas de
  mayor impacto secuencial; las tres primeras presentan valores de
  $-0.76$, $-0.73$ y $-0.71$~trimestres respectivamente.
  \textbf{A diferencia de Electrónica, donde una sola arista domina, Química
  muestra un perfil de impactos comprimidos en una cadena secuencial: no hay
  una seriación aislada que concentre el beneficio, sino que el problema es
  la longitud de la cadena temática completa.}}
\label{fig:efecto_qui}
\end{figure}

En Química (Figura~\ref{fig:efecto_qui}) el patrón es diferente: no hay una
arista dominante sino una cadena de tres con impactos similares.
La lectura práctica es que eliminar sólo la primera arista ganaría 0.76~trim;
eliminar las tres simultáneamente ganaría más que la suma por efectos de
composición.
La cuantificación precisa de estos efectos de composición ---los beneficios
de remover simultáneamente las $N$~aristas de mayor impacto--- se analiza
en la Sección~\ref{sec:dosis_respuesta}.

Las figuras equivalentes para los diez planes están en
\texttt{Output/figuras/efecto\_\{plan\}.pdf}.

%══════════════════════════════════════════════════════
\section{Correlaciones estadísticas: determinantes estructurales de la~ET}
\label{sec:correlaciones}
%══════════════════════════════════════════════════════

Las secciones anteriores caracterizan la eficiencia terminal plan a plan y evalúan el impacto de cada seriación en aislamiento.
Esta sección aborda la pregunta inversa: ¿qué rasgos estructurales del plan o de la UEA predicen los resultados?
La respuesta permite orientar las decisiones de reforma hacia los factores de mayor palanca antes de analizar cada arista por separado.

El análisis opera en tres niveles jerárquicos: el plan ($N=10$), la UEA
($N=606$ nodos con datos completos) y la arista de seriación ($N=499$).
El perfil del cuerpo docente ---el cuarto nivel de resolución--- se presenta en la
\S\ref{sec:perfil_docente}.
Las correlaciones son de Spearman (sin supuesto de normalidad).
Los modelos de regresión se reportan con la advertencia de que los coeficientes
de determinación son bajos, lo que constituye en sí mismo un hallazgo.

\subsection{Nivel plan: número de seriaciones y fracción de UEAs difíciles}

Con sólo diez observaciones, la potencia estadística es limitada.
Las correlaciones significativas a $p<0.05$ se interpretan como señales cualitativas, no como estimaciones de precisión.

\begin{table}[H]\centering
\caption{Correlaciones de Spearman entre rasgos estructurales del plan y resultados ($|r_s|>0.60$, $N=10$).
  Las dos variables de mayor poder predictivo aparecen en \textbf{negrita}.
  $\mathrm{Trim}_{50}$ = trimestre en que el 50\% de la cohorte ha egresado.}
\label{tab:corr_plan}
\small
\begin{tabular}{llccc}
\toprule
Resultado & Predictor & Dirección & $r_s$ & $p$ \\
\midrule
$\ET{21}$ & \textbf{Fracc.\ UEAs $a_i<0.65$} & alta $\Rightarrow$ menor & $-0.84$ & $0.004$ \\
$\mathrm{Trim}_{50}$ & \textbf{Fracc.\ UEAs $a_i<0.65$} & alta $\Rightarrow$ tardío & $+0.80$ & $0.005$ \\
Tasa egreso & \textbf{N seriaciones} & más $\Rightarrow$ menor & $-0.80$ & $0.006$ \\
$\mathrm{Trim}_{50}$ & Media $a_i$ & alta $\Rightarrow$ pronto & $-0.78$ & $0.008$ \\
$\ET{24}$ & Media $a_i$ & alta $\Rightarrow$ mayor & $+0.76$ & $0.016$ \\
Tasa baja & Profundidad topológica & mayor $\Rightarrow$ mayor & $+0.70$ & $0.023$ \\
Tasa egreso & Profundidad topológica & mayor $\Rightarrow$ menor & $-0.66$ & $0.039$ \\
\bottomrule
\end{tabular}
\end{table}

Dos variables concentran la mayor parte del poder predictivo (Tabla~\ref{tab:corr_plan}):
la \emph{fracción de UEAs con tasa de aprobación $a_i<0.65$} y el
\emph{número absoluto de seriaciones del plan}.
Planes con muchas UEAs difíciles tardan más en que su cohorte egrese;
planes con más seriaciones presentan menores tasas de egreso y mayor tasa de baja.
La profundidad topológica y la tasa media de aprobación son predictores secundarios pero consistentes.

\begin{figure}[H]\centering
\includegraphics[width=\linewidth]{../../Output/figuras_correlaciones/fig_02_scatter_plan_lm.pdf}
\caption{Las cuatro correlaciones de Spearman más fuertes entre rasgos
  estructurales de los planes y resultados de egreso ($N=10$~programas CBI).
  Cada punto es un plan; la banda sombreada es el intervalo de confianza al
  95\% de la regresión lineal.
  Aquí $a_i$ denota la tasa de aprobación histórica de la UEA (probabilidad
  de acreditar en un intento), ``UEAs difíciles'' son aquellas con $a_i<0.65$,
  $\mathrm{Trim}_{50}$ es el trimestre en que el 50\% de la cohorte ha
  egresado, y $\ET{21}$ es la fracción egresada a los 21~trimestres.
  \emph{Superior izquierda}: número de seriaciones vs.\ tasa de egreso
  ($r_s=-0.80$, $p=0.006$); planes con más de 130~seriaciones concentran
  las menores tasas de egreso.
  \emph{Superior derecha}: fracción de UEAs difíciles vs.\ $\mathrm{Trim}_{50}$
  ($r_s=+0.80$, $p=0.005$).
  \emph{Inferior izquierda}: tasa media de aprobación vs.\ $\ET{21}$
  ($r_s=+0.72$, $p=0.024$).
  \emph{Inferior derecha}: profundidad topológica del DAG (grafo acíclico
  dirigido de seriaciones) vs.\ tasa de baja ($r_s=+0.70$, $p=0.023$).
  \textbf{Las dos variables con mayor poder predictivo --- número de seriaciones
  y fracción de UEAs difíciles --- son directamente intervenibles mediante
  reforma curricular; la asociación persiste al controlar por el tamaño del plan.}}
\label{fig:scatter_plan}
\end{figure}

La Figura~\ref{fig:scatter_plan} ilustra las cuatro relaciones más fuertes.
El plan Industrial ocupa el extremo favorable en prácticamente todos los paneles:
pocas seriaciones (83), alta tasa media de aprobación ($\bar{a}=0.80$), menor profundidad topológica (5) y la menor tasa de baja de la División (32.5\%).
Electrónica ocupa el polo opuesto: profundidad~9, 139~seriaciones y la tasa de egreso más baja (57.5\%).
Computación tiene el mayor número absoluto de seriaciones (157) y la segunda tasa de egreso más baja (59.6\%).

Para descartar que la correlación $n_\text{seriaciones} \to$ tasa de egreso sea un artefacto de planes más grandes ---que naturalmente tienen más seriaciones y más UEAs difíciles---, se calculan correlaciones parciales de Spearman con dos controles independientes.
El efecto persiste al controlar por $n_\text{UEAs}$: $r_\text{parcial} = -0.756$ ($p=0.018$); y también al controlar por $\bar{a}_i$: $r_\text{parcial} = -0.649$ ($p=0.058$).
Análogamente, el efecto de la profundidad topológica sobre la tasa de baja es robusto al control por $n_\text{UEAs}$: $r_\text{parcial} = +0.746$ ($p=0.021$).
Esto descarta el confusor de escala del plan y es consistente con una relación estructural entre la arquitectura de seriaciones ---más que el simple tamaño del currículo--- y las restricciones de acceso a la trayectoria de egreso.

\begin{figure}[H]\centering
\includegraphics[width=\linewidth]{../../Output/figuras_correlaciones/fig_03_paralelas.pdf}
\caption{Perfil estructural y de resultados de los diez planes de CBI en
  coordenadas paralelas.
  Cada línea es un plan; el color codifica la tasa de egreso (amarillo claro
  = alta, morado = baja).
  Todas las variables se normalizan a $[0,1]$ y se orientan de modo que
  \emph{arriba} represente siempre el resultado más favorable.
  \emph{Zona rosada} (tres ejes): variables de estructura curricular --- número
  de seriaciones (invertido), fracción de UEAs con tasa de aprobación $a_i < 0.65$
  (invertida), y profundidad topológica del DAG (invertida).
  \emph{Zona azul} (cuatro ejes): indicadores de egreso --- tasa de egreso,
  $\ET{21}$ (fracción egresada a los 21~trimestres), $\mathrm{Trim}_{50}$
  (trimestre mediano de egreso, invertido) y tasa de baja (invertida).
  \textbf{Industrial (amarillo claro) ocupa la posición superior en ambas
  zonas; Electrónica (morado) muestra el perfil más adverso y la mayor
  separación entre las zonas de estructura y de resultados, lo que ilustra
  visualmente cómo la complejidad curricular se traduce en peores trayectorias
  de egreso.}}
\label{fig:paralelas}
\end{figure}

\subsection{Nivel UEA: el mapa de riesgo topológico}
\label{sec:riesgo_uea}

La pregunta a nivel de UEA es qué nodos del DAG de seriaciones condicionan más la trayectoria de los estudiantes.
La respuesta exige combinar dos dimensiones que son independientes entre sí:
la \emph{dificultad intrínseca} del nodo (capturada por $1-a_i$) y su \emph{posición topológica}.

La posición topológica se caracteriza por tres métricas: el número de descendientes en el DAG (\texttt{desc}), el grado de salida (\texttt{out\_deg}) y la centralidad de intermediación normalizada (\texttt{betw}).
De las tres, la que mejor predice el impacto marginal de eliminar las seriaciones salientes de una UEA es \texttt{desc} ($r_s=0.77$), seguida de \texttt{out\_deg} ($r_s=0.68$) y \texttt{betw} ($r_s=0.65$).
La tasa de aprobación $a_i$ por sí sola no correlaciona con el impacto ($r_s\approx -0.05$): una UEA difícil en posición periférica del DAG tiene impacto marginal mínimo.

Se define el \emph{índice de riesgo topológico enriquecido} de la UEA~$v$ como
\begin{equation}
  \rho_v
  \;=\; (1 - a_i)\;\cdot\;
        \frac{\mathrm{betw}_v}{\displaystyle\max_{u}\,\mathrm{betw}_u}
        \;\cdot\;
        \ln\!\bigl(1 + \mathrm{desc}_v\bigr)
        \;\cdot\;
        (1 + \mathrm{CV}_v)
        \;\cdot\;
        (1 + P_{\mathrm{exp},v}),
  \label{eq:risk}
\end{equation}
donde $\mathrm{CV}_v$ es el coeficiente de variación de la eficacia entre los
profesores que han impartido~$v$ (heterogeneidad docente intra-UEA) y
$P_{\mathrm{exp},v}$ es la probabilidad histórica de acumular cinco~\textsc{na}
consecutivos (expulsión).
Los dos últimos factores enriquecen el índice original
$(1-a_i)\cdot\mathrm{betw}_{\mathrm{norm}}\cdot\ln(1+\mathrm{desc})$ con
dimensiones no capturadas por la topología del DAG.
La correlación de Spearman entre el ranking original y el enriquecido es
$r_s=0.9997$: el orden de los nodos es prácticamente idéntico, pero la
\emph{caracterización} del riesgo cambia significativamente.

\begin{table}[H]\centering
\caption{Top-10 UEAs por índice de riesgo enriquecido~$\rho_v$~(\ref{eq:risk}).
  $a_i$: tasa de aprobación.
  Desc: descendientes en el DAG unificado.
  CV: coeficiente de variación de eficacia entre profesores.
  $P_{\mathrm{exp}}$: probabilidad de acumular 5~\textsc{na}.
  Tipo: D+I = docente + intrínseca; PD = principalmente docente.}
\label{tab:top10_riesgo}
\small
\begin{tabular}{llcrrcc}
\toprule
Código & UEA & $a_i$ & Desc & CV & $P_{\mathrm{exp}}$ & Tipo \\
\midrule
1112029 & Cálculo Integral                            & 0.54 & 416 & 0.42 & 2.2\% & D+I \\
1111081 & Dinámica del Cuerpo Rígido                  & 0.68 & 306 & 0.39 & 0.7\% & PD  \\
1112043 & Cálculo Diferencial                         & 0.51 & 488 & 0.47 & 4.3\% & D+I \\
1112030 & Ecuaciones Diferenciales Ordinarias         & 0.57 & 380 & 0.37 & 1.4\% & D+I \\
1111079 & Cinemática y Dinámica de Partículas         & 0.57 & 308 & 0.38 & 1.9\% & D+I \\
1112013 & Complementos de Matemáticas                & 0.54 & 231 & 0.39 & 2.2\% & D+I \\
1145054 & Ingeniería de los Materiales                & 0.65 &  75 & 0.32 & 1.0\% & PD  \\
1113046 & Termodinámica                               & 0.72 & 165 & 0.31 & 0.5\% & PD  \\
1142006 & Mecánica de Sólidos~I                       & 0.59 &  54 & 0.44 & 0.9\% & D+I \\
1111083 & Intro.\ Electrostática y Magnetostática     & 0.65 &  80 & 0.34 & 0.8\% & PD  \\
\bottomrule
\end{tabular}
\end{table}

Las cuatro primeras posiciones las ocupan las UEAs de la cadena de Cálculo
y Mecánica: Cálculo Diferencial $\to$ Cálculo Integral $\to$ Ecuaciones Diferenciales
Ordinarias, con Cinemática y Dinámica de Partículas en posición paralela.
Estas UEAs comparten cuatro características: tasas de aprobación bajas
($a_i = 0.51$--$0.57$), obligatorias en los diez planes, cientos de UEAs
descendientes, \emph{y} alta heterogeneidad docente intra-UEA
($\mathrm{CV} = 0.37$--$0.47$, categoría \emph{Docente + intrínseca}).
Este último rasgo revela que parte de la dificultad de estas UEAs no es
inherente al contenido sino variable según el docente asignado,
lo que abre una vía de intervención complementaria a la revisión estructural
del plan: la homogeneización de la práctica docente.
Son los nodos donde la inversión en apoyo académico produce el mayor retorno
sistémico para la División, tanto mediante remoción de seriaciones como mediante
la homogeneización de la práctica pedagógica.
Las Figuras~\ref{fig:ranking_enriquecido} y~\ref{fig:delta_ranking} visualizan
el ranking completo y el efecto de incorporar las dimensiones docente y de
expulsión sobre las posiciones individuales.

\begin{figure}[H]\centering
\includegraphics[width=\linewidth]{../../Output/figuras_correlaciones/fig_23_ranking_enriquecido.pdf}
\caption{Ranking de las UEAs de la División CBI por índice de riesgo topológico
  enriquecido $\rho_v$~(\ref{eq:risk}), que integra cinco factores multiplicativos:
  dificultad de acceso~$(1-a_i)$, centralidad de intermediación normalizada
  $\mathrm{betw}_{\mathrm{norm}}$, riqueza estructural $\ln(1+\mathrm{desc})$,
  heterogeneidad docente~$\mathrm{CV}$ y probabilidad de expulsión~$P_{\mathrm{exp}}$.
  Cada barra es una UEA con al menos un descendiente en el DAG (grafo acíclico
  dirigido de seriaciones); las UEAs se ordenan de mayor a menor~$\rho_v$ y el
  color identifica el plan de estudios al que pertenece (amb = Ambiental,
  civ = Civil, com = Computación, ele = Eléctrica, elo = Electrónica,
  fis = Física, ind = Industrial, mec = Mecánica,
  met = Metalurgia, qui = Química).
  \textbf{Las cinco UEAs de mayor riesgo --- Cálculo Diferencial, Cálculo
  Integral, Ecuaciones Diferenciales Ordinarias, Cinemática y Dinámica de
  Partículas, y Complementos de Matemáticas --- están presentes en los diez
  planes y concentran simultáneamente baja aprobación, alta centralidad
  topológica y alta varianza docente, lo que las convierte en el objetivo
  principal de cualquier política de intervención divisional.}}
\label{fig:ranking_enriquecido}
\end{figure}

\begin{figure}[H]\centering
\includegraphics[width=\linewidth]{../../Output/figuras_correlaciones/fig_24_delta_ranking.pdf}
\caption{Desplazamiento de posición en el ranking al pasar del índice de riesgo
  topológico original $(1-a_i)\cdot\mathrm{betw}_{\mathrm{norm}}\cdot\ln(1+\mathrm{desc})$
  al índice enriquecido~$\rho_v$~(\ref{eq:risk}), que añade la heterogeneidad
  docente~$\mathrm{CV}$ y la probabilidad de expulsión~$P_{\mathrm{exp}}$.
  Cada punto es una UEA; el eje horizontal es la posición ordinal en el ranking
  original y el eje vertical la posición en el ranking enriquecido
  (posición~1 = mayor riesgo en ambos ejes); puntos por encima de la diagonal
  ascienden al enriquecerse el índice (resultaban más riesgosos de lo que
  indicaba la topología pura) y puntos por debajo descienden.
  Los puntos de mayor desplazamiento se etiquetan con el nombre de la UEA.
  \textbf{La correlación de Spearman entre ambos rankings es $r_s = 0.9997$:
  el diagnóstico estratégico no se altera, pero Ingeniería Geotécnica
  (CV$=0.93$) asciende del rango~34 al~19 y tres UEAs con $P_{\mathrm{exp}}>0.30$
  emergen como focos de pérdida estudiantil directa, invisibles en el análisis
  topológico puro.}}
\label{fig:delta_ranking}
\end{figure}

\begin{figure}[H]\centering
\includegraphics[width=\linewidth]{../../Output/figuras_correlaciones/fig_05_bottleneck.pdf}
\caption{Mapa de riesgo topológico de las UEAs (Unidades de Enseñanza-Aprendizaje)
  de CBI que actúan como prerequisitos en el DAG (grafo acíclico dirigido de
  seriaciones).
  Cada punto es una UEA con al menos un descendiente en el DAG.
  \emph{Eje horizontal}: tasa de aprobación $a_i$ (probabilidad de acreditar
  la UEA en un intento), escala lineal.
  \emph{Eje vertical}: centralidad de intermediación normalizada (\emph{betweenness}),
  en escala logarítmica; mide qué fracción de los caminos más cortos del DAG
  pasan por ese nodo.
  \emph{Tamaño}: logaritmo del número de UEAs descendientes (cursos que requieren
  directa o indirectamente haber aprobado esta UEA).
  \emph{Color}: índice de riesgo enriquecido $\rho_v \propto (1-a_i)
  \cdot \mathrm{betw}_{\mathrm{norm}} \cdot \ln(1+\mathrm{desc})
  \cdot (1+\mathrm{CV}) \cdot (1+P_{\mathrm{exp}})$, donde CV es el
  coeficiente de variación de la eficacia entre los docentes que han impartido
  la UEA y $P_{\mathrm{exp}}$ es la probabilidad histórica de acumular cinco
  evaluaciones No Acreditadas.
  La zona sombreada (cuadrante superior izquierdo) agrupa las UEAs de alto
  riesgo: baja aprobación \emph{y} alta centralidad simultáneamente.
  Los umbrales delimitadores son $a_i=0.70$ (línea vertical) y la mediana de
  \texttt{betw} (línea horizontal).
  Las etiquetas identifican las UEAs en el percentil~95 de~$\rho_v$.
  \textbf{Cálculo Diferencial, Cálculo Integral y Ecuaciones Diferenciales
  Ordinarias dominan el cuadrante de alto riesgo por combinar baja aprobación,
  centralidad máxima y centenares de UEAs descendientes; son los nodos donde
  la inversión en apoyo académico produce el mayor retorno sistémico.}}
\label{fig:bottleneck}
\end{figure}

La Figura~\ref{fig:bottleneck} pone de manifiesto que la gran mayoría de las
689~UEAs del grafo unificado tiene centralidad nula o próxima a cero (no están
en caminos críticos del DAG).
El contraste estadístico entre el 20\% de UEAs con mayor~$\rho_v$ (cuellos de
botella) y el 80\% restante es significativo en todas las variables relevantes
(test de Wilcoxon, $p<0.001$): los cuellos de botella tienen tasa de aprobación
mediana de 0.69 frente a 0.82 del resto, y generan un impacto marginal medio de
0.066~trim frente a 0.000~trim del resto.

\subsection{Concentración del impacto marginal}
\label{sec:pareto}

\begin{notabox}
Los predictores evaluados en esta sección operan a nivel de \emph{arista
individual}, a diferencia de la sección anterior (\S\ref{sec:riesgo_uea})
donde operan a nivel de \emph{nodo agregado}.
La discrepancia en el poder predictivo de \texttt{desc}
($r_s=0.77$ a nivel de nodo vs.\ $r_s=-0.017$ a nivel de arista) es
consecuencia directa de la asimetría de la distribución de impactos: unas
pocas aristas concentran casi todo el efecto, y el nodo fuente predice
bien el impacto \emph{acumulado} de todas sus aristas salientes, pero no el
de cada arista individual.
\end{notabox}

\noindent\textbf{Nota sobre el universo de aristas.} El análisis de
impacto marginal (\S\ref{sec:efecto}) evaluó 499
aristas con efecto no nulo sobre $\ETtit$. El
análisis de correlación de esta sección opera sobre ese mismo universo
($N=499$ aristas con datos estructurales completos para la regresión),
confirmado por el archivo \texttt{Output/R\_analisis/edge\_features.csv}
(499 filas de datos).

La distribución del impacto marginal $|\Delta\ETtit|$ sobre las 499~seriaciones
evaluadas es fuertemente asimétrica.

\begin{figure}[H]\centering
\includegraphics[width=\linewidth]{../../Output/figuras_correlaciones/fig_08_pareto_aristas.pdf}
\caption{Concentración del impacto de las seriaciones sobre el tiempo de egreso
  $\ETtit = \mathbb{E}[T\,|\,\mathrm{tit}]$ (tiempo promedio condicionado al
  egreso, en trimestres).
  \emph{Izquierda}: las 30~aristas del DAG (grafo acíclico dirigido de
  seriaciones) de mayor impacto absoluto $|\Delta\ETtit|$ en toda la División,
  ordenadas de mayor a menor y coloreadas por plan.
  La seriación Ecuaciones Diferenciales Ordinarias $\to$ Circuitos Eléctricos~I
  (plan Electrónica) lidera con $|\Delta\ETtit|=1.07$~trimestres.
  \emph{Derecha}: curva de Lorenz del impacto acumulado sobre las 499~seriaciones
  evaluadas; el eje horizontal es la fracción de aristas (ordenadas de mayor a
  menor impacto) y el eje vertical es la fracción acumulada del impacto total.
  \textbf{El 26.5\% de las aristas de seriación concentra el 80\% del
  impacto total sobre el tiempo de egreso; este perfil Pareto implica que una
  reforma curricular acotada a un cuarto de las seriaciones capturaría la gran
  mayoría del beneficio alcanzable.}}
\label{fig:pareto}
\end{figure}

El 26.5\% de las aristas evaluadas (132 de 499) concentra el 80\% del impacto total
sobre $\ETtit$ (Figura~\ref{fig:pareto}).
Los modelos de regresión múltiple que intentan predecir $|\Delta\ETtit|$ a partir
de rasgos estructurales del nodo fuente alcanzan $R^2\approx 2$--$4\,\text{\%}$:
el impacto no es bien predicho por características individuales.
El análisis de Spearman sobre las 499~aristas confirma esta baja predictabilidad:
el número de descendientes del nodo fuente ($r_s = -0.017$, $p = 0.71$) y su
centralidad de intermediación ($r_s = 0.004$, $p = 0.93$) no predicen el impacto
de una arista específica, aunque ambas métricas sí predicen el impacto
\emph{agregado} de un nodo (\S\ref{sec:riesgo_uea}).
Esto es una consecuencia directa de la asimetría extrema: unas pocas aristas
concentran casi todo el efecto, y sus características estructurales no difieren
suficientemente del resto para que ningún modelo las identifique a priori.

El predictor más fuerte del impacto sobre~$\ETtit$ a nivel de arista individual
es el impacto de esa misma arista sobre la tasa de baja
($r_s = 0.569$, $p < 10^{-42}$): las seriaciones que más retrasan el egreso
son también las que más aumentan el abandono, lo que confirma que el ranking de
las 132~aristas prioritarias optimiza simultáneamente ambos resultados sin
necesidad de ponderar los objetivos.
En segundo lugar, los créditos del nodo posterior son un predictor positivo
pero modesto ($r_s = 0.158$, $p = 0.0004$): aristas hacia UEAs de mayor peso
crediticio tienen un impacto marginal ligeramente superior.
Esta asociación es accionable: al diseñar portfolios quirúrgicos conviene dar
preferencia a aristas hacia UEAs de alto valor crediticio cuando el impacto
marginal individual está empatado.

La implicación práctica es directa: el listado de las 132~aristas de mayor impacto
(disponible en \texttt{Output/R\_analisis/edge\_features.csv}) es más accionable
que cualquier predictor basado en características de los nodos.
El índice~$\rho_v$~(\ref{eq:risk}) sirve para identificar qué \emph{nodos} merecen
apoyo académico prioritario; el ranking de impacto marginal, para seleccionar qué
\emph{aristas} conviene revisar en una reforma curricular.

\subsection{Dimensiones adicionales: recuperabilidad, heterogeneidad docente
            y riesgo de expulsión}
\label{sec:dimensiones_adicionales}

Los análisis de las secciones~\ref{sec:riesgo_uea} y~\ref{sec:pareto} utilizan
como insumo principal la tasa de aprobación~$a_i$ y la posición topológica de
cada UEA.
El diccionario de datos históricos contiene tres dimensiones adicionales que
enriquecen el diagnóstico sin alterar el ranking de riesgo
($r_s=0.9997$ entre el índice original y el enriquecido).

\subsubsection{Recuperabilidad}

El campo \texttt{Vector\_Prob\_Aprobacion}$=[P(\leq k)]_{k=1}^{5}$ reporta la
probabilidad acumulada de aprobar en a lo sumo $k$~intentos.
El \emph{spread} $\Delta P = P(\leq 5)-P(\leq 1)$ mide cuánto mejora el
estudiante con intentos adicionales.

De las 639~UEAs obligatorias de CBI con datos completos, \textbf{164}
(25.7\%) se encuentran en el cuadrante <<rescatable>> ($P(\leq1)<0.65$,
$\Delta P > 0.25$): son materias difíciles en el primer intento, pero
recuperables con oportunidades adicionales.
Ninguna UEA del conjunto cae en el cuadrante de <<dificultad persistente>>
($\Delta P < 0.15$): la cota de aprobación a cinco intentos supera~0.97 en
todos los casos.
Esto implica que el problema de primer acceso en CBI no es de incapacidad
acumulada sino de contexto de inscripción: los estudiantes que no pueden
cursar una UEA en el trimestre óptimo por restricciones de seriación
pierden intentos antes de estar académicamente listos.

Las cinco UEAs de mayor riesgo son todas <<rescatables>>
($\Delta P = 0.36$--$0.49$) con spread que correlaciona positivamente con
la centralidad topológica ($r_s = 0.25$): remover sus seriaciones salientes
tiene efecto doble --- libera la restricción de acceso \emph{y} permite al
estudiante llegar a la UEA con más margen de intentos disponibles.

\subsubsection{Gap de recuperación académica}

Cada UEA con grupos de recuperación permite contrastar la evaluación
global ordinaria (\texttt{grupos\_eva\-lua\-cion\_glo\-bal}) con los
grupos de recuperación (\texttt{recu\-pe\-ra\-cion\_ueas}).
La diferencia $\delta a = a_i(\text{rec})-a_i(\text{eval})$ mide si la
segunda oportunidad institucional mejora la tasa de aprobación.

CBI es la \emph{única} división con $\delta a$ medio positivo
($+0.015$, frente a $-0.064$ en la División de Ciencias y Artes para el Diseño (CYAD) y $-0.126$ en la División de Ciencias Sociales y Humanidades (CSH)).
El 41.7\% de las UEAs CBI tienen $\delta a < 0$ --- la recuperación empeora
la tasa ---, y el 58.3\% la mejora.
Entre las UEAs de mayor riesgo topológico, Cálculo Integral, Cálculo
Diferencial y EDO muestran $\delta a \in [+0.07, +0.11]$: una proporción
no despreciable (27\,--\,33\% de los inscritos en evaluación normal)
llega a recuperación y, en estos tres casos, la supera con mayor tasa que en
la evaluación ordinaria.
La centralidad topológica correlaciona negativamente con $\delta a$
($r_s = -0.15$): precisamente los nodos más críticos son aquellos donde la
recuperación proporciona menor beneficio, lo que indica saturación del
mecanismo de segunda oportunidad en los cuellos de botella principales.
La Figura~\ref{fig:gap_recuperacion} muestra la distribución de~$\delta a_i$
en CBI con énfasis en la relación entre el beneficio de la recuperación y
la posición topológica de cada UEA.

\begin{figure}[H]\centering
\includegraphics[width=\linewidth]{../../Output/figuras_correlaciones/fig_20_gap_recuperacion.pdf}
\caption{Diferencia de tasa de aprobación entre la evaluación de recuperación
  y la evaluación global ordinaria
  ($\delta a_i = a_i(\mathrm{rec}) - a_i(\mathrm{eval})$) para las UEAs de
  la División CBI con datos disponibles en ambas modalidades
  (fuente: secciones \texttt{recuperacion\_ueas} y
  \texttt{grupos\_evaluacion\_global} del \texttt{diccionario\_ueas.json}).
  El eje horizontal es la tasa de aprobación en evaluación ordinaria~$a_i(\mathrm{eval})$;
  el eje vertical es~$\delta a_i$ (valores positivos: la recuperación mejora la
  tasa; valores negativos: la empeora).
  El tamaño de cada punto es proporcional a la fracción de inscritos que llega
  a evaluación de recuperación; el color indica el plan de estudios.
  Las UEAs de mayor centralidad topológica (Cálculo Diferencial, Cálculo
  Integral y EDO) se etiquetan explícitamente.
  \textbf{CBI es la única división con $\delta a$ medio positivo ($+0.015$),
  pero la correlación entre centralidad topológica y~$\delta a$ es negativa
  ($r_s = -0.15$): los nodos más críticos del DAG son donde la segunda
  oportunidad produce el menor beneficio diferencial, lo que confirma que
  la restricción estructural de trayectoria --- y no la incapacidad de
  aprendizaje --- es el factor dominante en los cuellos de botella de CBI.}}
\label{fig:gap_recuperacion}
\end{figure}

\subsubsection{Heterogeneidad docente intra-UEA}

El \textsc{cv} de la eficacia entre los profesores que han impartido una
UEA mide la componente docente del riesgo.
De las 521~UEAs CBI con datos de tres o más profesores,
el \textsc{cv} mediano es 0.21 y el 34\% tiene \textsc{cv}$>0.30$
(alta heterogeneidad).

El \textsc{cv} correlaciona positivamente con el índice de riesgo~$\rho_v$
($r_s = 0.22$): los cuellos de botella topológicos también presentan
alta varianza docente.
La clasificación resultante distingue:
\begin{itemize}[leftmargin=2em]
  \item \emph{Docente + intrínseca} (\textsc{cv}$>0.30$,\;$a_i<0.65$):
        Cálculo Diferencial, Cálculo Integral, EDO, Cinemática y
        Dinámica de Partículas, Complementos de Matemáticas.
        La intervención óptima combina reducción de seriaciones
        \emph{y} homogeneización pedagógica.
  \item \emph{Principalmente docente} (\textsc{cv}$>0.30$,\;$a_i\geq0.65$):
        Dinámica del Cuerpo Rígido, Termodinámica, Ingeniería de los Materiales.
        El contenido no es intrínsecamente difícil para quien accede; la
        varianza la induce la práctica docente diferenciada.
\end{itemize}
Una UEA nueva entra al top-20 de riesgo cuando se incluye el \textsc{cv}:
Ingeniería Geotécnica (\textsc{cv}$=0.93$, rango original~34, enriquecido~19),
invisible en el análisis topológico puro por su menor betweenness.
La Figura~\ref{fig:heterogeneidad_docente} muestra la distribución divisional
del~\textsc{cv} y su co-ocurrencia con la posición topológica crítica.

\begin{figure}[H]\centering
\includegraphics[width=\linewidth]{../../Output/figuras_correlaciones/fig_21_heterogeneidad_docente.pdf}
\caption{Heterogeneidad de la eficacia docente intra-UEA, medida por el
  coeficiente de variación (\textsc{cv} $= \sigma/\mu$) de la eficacia entre
  todos los profesores que han impartido la UEA en los periodos~16I--25O.
  \emph{Panel izquierdo}: distribución del~\textsc{cv} sobre las 521~UEAs
  de CBI con datos de tres o más profesores; la línea de trazos vertical
  marca el umbral de alta heterogeneidad (\textsc{cv}$>0.30$).
  \emph{Panel derecho}: \textsc{cv} vs.\ centralidad de intermediación
  normalizada ($\mathrm{betw}_{\mathrm{norm}}$) para las UEAs con al menos
  un descendiente en el DAG; el color indica la tasa de aprobación~$a_i$
  y el tamaño el número de descendientes~$\mathrm{desc}$.
  Las UEAs del Tronco General se etiquetan por nombre.
  \textbf{El 34\% de las UEAs de CBI presenta alta heterogeneidad docente
  (CV${}>0.30$), y esta varianza co-ocurre sistemáticamente con alta
  centralidad topológica: los nodos que más penalizan la trayectoria de
  egreso son también los que muestran mayor dispersión de resultados
  entre los docentes que los imparten, lo que abre una vía de intervención
  pedagógica complementaria a la reforma estructural del plan de estudios.}}
\label{fig:heterogeneidad_docente}
\end{figure}

\begin{figure}[H]\centering
\includegraphics[width=\linewidth]{../../Output/figuras_correlaciones/fig_gap_eficacia_riesgo.pdf}
\caption{Brecha de eficacia docente en las UEAs (Unidades de Enseñanza-Aprendizaje)
  de mayor riesgo sistémico en CBI.
  \emph{Eje horizontal}: coeficiente de variación (\textsc{cv}) de la eficacia
  entre todos los profesores que han impartido esa UEA en el periodo 16I--25O;
  el \textsc{cv} mide la heterogeneidad de la práctica docente intra-UEA.
  \emph{Eje vertical}: brecha relativa entre la eficacia del mejor docente y la
  eficacia mediana de los docentes de esa UEA (escala fraccional, de 0 a 1).
  \emph{Tamaño}: número de UEAs descendientes en el DAG (grafo acíclico dirigido)
  unificado de seriaciones (indicador de alcance sistémico del nodo).
  \emph{Color}: tasa de aprobación histórica $a_i$ (probabilidad de acreditar
  en un intento); tonos fríos = baja aprobación.
  Las UEAs del cuadrante superior derecho (\textsc{cv}~alto y brecha~alta)
  son las de mayor potencial de mejora mediante intervención docente: el
  \emph{mejor} profesor ya logra resultados cualitativamente superiores a la
  mediana, de modo que la dificultad no es intrínseca al contenido sino
  atribuible al docente asignado.
  \textbf{Cálculo Diferencial y Cálculo Integral acumulan simultáneamente
  alta heterogeneidad docente, centenares de UEAs descendientes y baja
  aprobación, lo que los convierte en los nodos de mayor retorno posible
  para una política combinada de reforma curricular y homogeneización pedagógica.}}
\label{fig:gap_eficacia}
\end{figure}

\subsubsection{Probabilidad de expulsión}

$P_{\mathrm{exp},v}$ es la probabilidad histórica de acumular cinco~\textsc{na}
en la UEA~$v$, lo que en el reglamento UAM-A desencadena la expulsión
definitiva de la licenciatura.
El modelo estocástico (\S\ref{sec:modelo}) ya implementa este límite
(\texttt{MAX\_INTENTOS}$=5$): los resultados de egreso reportados son,
por tanto, coherentes con la distribución empírica de~$P_{\mathrm{exp},v}$.

En CBI, la mediana es $P_{\mathrm{exp}}=0.0011$ y el 17\% de UEAs supera
el umbral del~1\%.
Tres UEAs presentan valores críticos:
UEA~1145092 ($P_{\mathrm{exp}}=0.51$, $a_i=0.12$),
1136028 ($P_{\mathrm{exp}}=0.33$, $a_i=0.60$) y
1154027 ($P_{\mathrm{exp}}=0.30$, $a_i=0.33$).
Ninguna de las tres ocupa posición central en el DAG, por lo que su impacto
sobre~$\ETtit$ está contenido; no obstante, representan fuentes de pérdida
estudiantil directa que merecen atención pedagógica independiente de la
reforma curricular.
La Figura~\ref{fig:prob_expulsion} muestra la distribución de~$P_{\mathrm{exp}}$
en CBI y su relación con la profundidad topológica.

\begin{figure}[H]\centering
\includegraphics[width=\linewidth]{../../Output/figuras_correlaciones/fig_22_prob_expulsion.pdf}
\caption{Probabilidad histórica de expulsión definitiva
  ($P_{\mathrm{exp},v}$: probabilidad empírica de acumular cinco evaluaciones
  No~Acreditadas en la misma UEA~$v$; según el Reglamento de Estudios
  Superiores~\cite{RESUAM} Arts.~16.V y~44.XI, cinco~NA en la misma UEA
  implican baja definitiva de la licenciatura) para las UEAs de la División
  CBI con historial suficiente.
  \emph{Panel izquierdo}: distribución de~$P_{\mathrm{exp}}$ en escala
  logarítmica; la línea de trazos vertical marca el umbral de riesgo
  operativo ($P_{\mathrm{exp}} = 1\%$); las UEAs con $P_{\mathrm{exp}}>0.10$
  se etiquetan por nombre.
  \emph{Panel derecho}: $P_{\mathrm{exp}}$ vs.\ profundidad topológica del
  nodo (número de niveles desde la raíz del DAG); el color indica la tasa
  de aprobación~$a_i$.
  \textbf{La mediana divisional es $P_{\mathrm{exp}} = 0.0011$ y el 17\%
  de las UEAs supera el umbral del~1\%; aunque los tres casos críticos
  ($P_{\mathrm{exp}} > 0.30$) ocupan posiciones periféricas en el DAG y su
  impacto sobre~$\ETtit$ está contenido, constituyen fuentes de pérdida
  estudiantil directa que el modelo de cascada puede subestimar y que
  requieren intervención pedagógica urgente e independiente de la
  reforma curricular.}}
\label{fig:prob_expulsion}
\end{figure}

\begin{figure}[H]\centering
\includegraphics[width=0.90\linewidth]{../../Output/figuras_correlaciones/fig_19_recoverability.pdf}
\caption{Mapa de recuperabilidad: análisis de la capacidad de los estudiantes
  para aprobar UEAs (Unidades de Enseñanza-Aprendizaje) con intentos adicionales.
  Para cada UEA, $P(\leq k)$ denota la probabilidad acumulada de aprobar en a
  lo sumo $k$~intentos; $\Delta P = P(\leq5) - P(\leq1)$ mide cuánto mejoran
  los estudiantes con reintentos (recuperabilidad).
  \emph{Panel superior izquierdo}: $P(\leq1)$ vs.\ $\Delta P$ para todas las
  UEAs con datos; las líneas de trazos delimitan cuatro cuadrantes según la
  combinación de dificultad de primer acceso y capacidad de recuperación.
  Ninguna UEA cae en el cuadrante de ``dificultad persistente''
  ($P(\leq1) < 0.65$ y $\Delta P < 0.15$): la cota de aprobación a cinco
  intentos supera 0.97 en todos los casos.
  \emph{Panel superior derecho}: distribución de~$\Delta P$ por nivel
  topológico (profundidad del nodo en el DAG, grafo acíclico dirigido de
  seriaciones); los nodos más profundos tienden a menor recuperabilidad.
  \emph{Panel inferior izquierdo}: centralidad de intermediación
  (\emph{betweenness}) vs.\ $\Delta P$ para las UEAs difíciles
  ($P(\leq1) < 0.65$); el tamaño de cada punto es proporcional al número
  de UEAs descendientes.
  \emph{Panel inferior derecho}: curvas $P(k)$ de las ocho UEAs difíciles con
  mayor centralidad, mostrando la progresión de probabilidad de aprobación
  de $k=1$ a $k=5$~intentos.
  \textbf{El hallazgo central es que en CBI no existe dificultad intrínseca
  irreversible: las UEAs que son difíciles en el primer intento son recuperables
  con reintentos, por lo que el problema no es de incapacidad estudiantil sino
  de acceso limitado a las oportunidades de reintento que imponen las seriaciones.}}
\label{fig:recoverability}
\end{figure}

\subsection{Estructura factorial de los atributos de UEA}
\label{sec:factores_uea}

El conjunto de 15~variables disponibles por UEA presenta redundancia interna: las
métricas topológicas (\texttt{betw}, \texttt{out\_deg}, \texttt{desc}) co-varían,
y la eficacia, la eficiencia y~$a_i$ son casi collineales.
Un análisis factorial exploratorio (rotación oblimin, $N=606$~UEAs con datos
completos, 5~factores seleccionados por análisis paralelo~\cite{Horn1965}) revela la estructura
latente subyacente.

\begin{table}[H]\centering
\caption{Cargas factoriales oblimin de los 15~atributos de UEA sobre los
  5~factores latentes.
  Se omiten cargas $|c| < 0.30$.
  Varianza explicada total por los 5~factores: 66.7\%.}
\label{tab:factores_uea}
\small
\begin{tabular}{lccccc}
\toprule
Variable & MR3 & MR1 & MR2 & MR4 & MR5 \\
\midrule
eficacia                 & $+0.996$ & & & & \\
$a_i$                    & $+0.996$ & & & & \\
eficiencia               & $+0.948$ & & & & \\
spread$_{15}$            & $-0.664$ & & $+0.371$ & & \\
betw\_norm               & & $+0.998$ & & & \\
out\_deg                 & & $+0.630$ & & & \\
$n_\text{descendientes}$ & & $+0.510$ & & & \\
cv (eficacia)            & & & $+0.947$ & & \\
rango eficacia           & & & $+0.769$ & & \\
profundidad (depth)      & & & & $+0.716$ & \\
pagerank                 & & & & $+0.521$ & \\
$n_\text{planes}$        & & & & $-0.693$ & \\
frac\_rec                & & & & & $+0.792$ \\
créditos                 & & & & & $+0.486$ \\
$\delta a_i$             & & & & & \\
\bottomrule
\end{tabular}
\end{table}

Los cinco factores tienen interpretaciones claras y son mutuamente
independientes (correlaciones inter-factor $< 0.15$):

\begin{itemize}[leftmargin=2em]
  \item \textbf{MR3 — Calidad intrínseca del nodo}:
        eficacia, $a_i$ y eficiencia con carga $>0.94$.
        El spread negativo indica que los nodos de alta calidad mejoran menos
        con intentos adicionales (sus estudiantes ya aprueban al primer intento).

  \item \textbf{MR1 — Influencia topológica}:
        betweenness (carga prácticamente unitaria), grado de salida y
        número de descendientes.
        Es el factor que más predice el impacto marginal agregado de un nodo.

  \item \textbf{MR2 — Dispersión inter-docente}:
        coeficiente de variación y rango de eficacia entre profesores.
        Captura la dimensión de heterogeneidad pedagógica, \emph{ortogonal} a
        la calidad intrínseca (MR3) y a la topología (MR1).
        Su independencia respecto de MR3 implica que una UEA puede ser
        intrínsecamente difícil \emph{y} tener docentes uniformes, o
        fácil con docentes muy heterogéneos.

  \item \textbf{MR4 — Embeddedness curricular}:
        profundidad positiva, intersección de planes negativa y pagerank
        positivo.
        UEAs con carga alta en MR4 son nodos especializados:
        profundos en la secuencia, presentes en pocos planes, pero
        con alta influencia dentro de ellos.

  \item \textbf{MR5 — Recuperabilidad y peso crediticio}:
        fracción de estudiantes en recuperación y créditos de la UEA.
        UEAs de alto MR5 concentran intentos adicionales y tienen mayor peso
        en la trayectoria.
\end{itemize}

\begin{figure}[H]\centering
\includegraphics[width=\linewidth]{../../Output/figuras_correlaciones/fig_uea_factor_loadings.pdf}
\caption{Cargas factoriales de los 15~atributos de UEA sobre los 5~factores
  latentes identificados por análisis factorial exploratorio con rotación
  oblimin ($N=606$~UEAs con datos completos; varianza explicada total: 66.7\%).
  Cada panel muestra las cargas de cada variable en uno de los cinco factores;
  el color de la barra indica el signo de la carga (positivo/negativo).
  Los factores se ordenan por varianza explicada decreciente:
  MR3 captura la calidad intrínseca del nodo ($a_i$, eficacia, eficiencia);
  MR1 captura la influencia topológica (\emph{betweenness}, grado de salida,
  número de descendientes);
  MR2 captura la dispersión entre docentes (CV de eficacia, rango de eficacia);
  MR4 y MR5 capturan profundidad en el DAG y recuperabilidad, respectivamente.
  Las correlaciones inter-factor son menores a~0.15 (factores casi ortogonales).
  \textbf{La independencia entre MR3 (calidad intrínseca) y MR1 (topología)
  confirma que la dificultad de una UEA y su posición estratégica en el grafo
  son dimensiones independientes: el riesgo sistémico surge solo cuando ambas
  coinciden, lo que valida la forma multiplicativa del índice de riesgo~$\rho_v$
  frente a alternativas puramente aditivas.}}
\label{fig:factores_uea}
\end{figure}

La independencia entre MR3 y MR1 es el hallazgo estructural más relevante:
el riesgo sistémico ---que motivó la definición del índice~$\rho_v$ como
producto de dificultad intrínseca y centralidad topológica--- no es un único
factor latente sino la coincidencia de dos dimensiones ortogonales.
Este resultado valida la forma multiplicativa del índice~(\ref{eq:risk}) frente
a alternativas aditivas, que subestimarían la interacción entre las dos
dimensiones.

%══════════════════════════════════════════════════════
\section{Perfil del cuerpo docente CBI (1\,810 profesores UAM-A; 16I--25O)}
\label{sec:perfil_docente}
%══════════════════════════════════════════════════════

El análisis de correlaciones de la sección anterior opera sobre el currículo:
planes, UEAs y aristas.
Esta sección extiende el diagnóstico al nivel de los actores mediante el
diccionario de profesores, que registra 1\,810~docentes con actividad en UAM-A
en 16I--25O; 738 de ellos imparten UEAs en la División CBI.
Los indicadores disponibles por docente son: tasa de aprobación~$a_d$, eficacia,
eficiencia, rendimiento académico (escala~0--10) y probabilidad de expulsión
por UEA.

\subsection{Distribución divisional}

La mediana divisional de~$a_d$ es~$\mathbf{0.744}$
($Q_1=0.606$, $Q_3=0.828$, rango intercuartílico~0.222).
El~31.3\% de los 738~docentes CBI presenta~$a_d < 0.65$ ---el umbral
de riesgo alto en el mapa de UEAs--- y el~11.2\% cae por debajo de~0.50.
El rendimiento académico mediano es~6.32/10.
Cada profesor imparte, en promedio, 8.6~UEAs distintas en el periodo analizado.

\begin{figure}[H]\centering
\includegraphics[width=0.92\linewidth]%
  {../../Output/figuras_correlaciones/fig_25_distribucion_profesores.pdf}
\caption{Distribución de la tasa de aprobación docente~$a_d$ (fracción de
  estudiantes que aprueba al cursar con ese profesor) entre los
  738~docentes de la División CBI con actividad registrada en UAM-A
  durante los periodos 16I--25O (trimestres de invierno 2016 a otoño 2025).
  \emph{Izquierda}: histograma de~$a_d$ con líneas verticales en los
  percentiles~25, 50 y~75; el eje horizontal va de 0 a 1.
  \emph{Derecha}: distribución de los mismos 738~docentes agrupados por
  cuartil de $a_d$.
  La línea de trazos vertical marca el umbral de riesgo alto $a_d=0.65$
  (equivalente al umbral de UEAs difíciles).
  \textbf{El 31.3\% de los docentes CBI opera por debajo del umbral de
  riesgo $a_d=0.65$, lo que implica que cerca de un tercio del profesorado
  constituye, en sí mismo, un factor de fragilidad académica independiente
  de la estructura del plan de estudios.}}
\label{fig:distribucion_prof}
\end{figure}

\subsection{Heterogeneidad departamental}

La tabla~\ref{tab:dept_docentes} y la figura~\ref{fig:depto_prof} muestran
que hay diferencias sistemáticas entre departamentos.
El hallazgo más relevante es que \textbf{Ciencias Básicas} ---el departamento
que imparte las UEAs del Tronco General (Cálculo, Física, EDO, Química)---
tiene la mediana departamental más baja de CBI: $\tilde{a}_d = 0.670$,
con la mayor desviación estándar ($\sigma=0.167$) entre los departamentos con
cobertura amplia.
Por contraste, Energía ($\tilde{a}_d=0.792$) y Electrónica
($\tilde{a}_d=0.781$) operan notoriamente por encima del promedio divisional.

\begin{table}[H]\centering\small
\caption{Tasa de aprobación docente por departamento CBI
  (departamentos con $\geq 5$ profesores;
  $\tilde{a}_d$ = mediana, $\sigma$ = desviación estándar).}
\label{tab:dept_docentes}
\begin{tabular}{lrrrrr}
\toprule
Departamento & $n$ & $\tilde{a}_d$ & $\bar{a}_d$ & $\sigma$ & Eficacia med. \\
\midrule
Energía            & 154 & 0.792 & 0.774 & 0.140 & 0.792 \\
Electrónica        &  83 & 0.781 & 0.757 & 0.151 & 0.781 \\
Materiales         & 101 & 0.771 & 0.744 & 0.127 & 0.771 \\
Sistemas           &  85 & 0.745 & 0.720 & 0.162 & 0.745 \\
\textbf{Ciencias Básicas} & \textbf{272} & \textbf{0.670} & \textbf{0.649} & \textbf{0.167} & \textbf{0.670} \\
\bottomrule
\end{tabular}
\end{table}

\begin{figure}[H]\centering
\includegraphics[width=0.88\linewidth]%
  {../../Output/figuras_correlaciones/fig_26_profesores_por_depto.pdf}
\caption{Distribución de la tasa de aprobación docente~$a_d$ por departamento
  de la División CBI (solo departamentos con $\geq 5$~docentes con actividad
  en 16I--25O).
  Cada caja muestra la distribución de~$a_d$ de los profesores del departamento:
  los bigotes indican los percentiles~10 y~90, la caja el rango intercuartílico,
  y la línea central la mediana.
  La línea de trazos horizontal marca el umbral de riesgo $a_d=0.65$.
  \textbf{Ciencias Básicas --- el departamento más numeroso (272~docentes) y
  responsable de impartir las UEAs del Tronco General (Cálculo, Física,
  Ecuaciones Diferenciales, Química) --- presenta la mediana departamental
  más baja de CBI ($\tilde{a}_d = 0.670$) y la mayor desviación estándar
  ($\sigma=0.167$), concentrando así el factor de riesgo docente precisamente
  en los nodos topológicamente más críticos del currículo.}}
\label{fig:depto_prof}
\end{figure}

\subsection{Perfil docente en las UEAs de mayor riesgo}

El cruce entre el ranking de riesgo~$\rho_v$ y el campo \texttt{Profesores}
de cada UEA permite caracterizar quién imparte los nodos más críticos del DAG.

\begin{figure}[H]\centering
\includegraphics[width=0.96\linewidth]%
  {../../Output/figuras_correlaciones/fig_27_cobertura_critica.pdf}
\caption{Perfil del cuerpo docente en las 20~UEAs de mayor índice de riesgo
  enriquecido $\rho_v$ de la División CBI.
  \emph{Panel izquierdo}: eficacia mediana de los docentes asignados a cada
  UEA (fracción de estudiantes que aprueba con el docente mediano) vs.\ tasa
  de aprobación histórica~$a_i$ de la UEA (probabilidad de acreditar en un
  intento, promediada sobre todos los docentes); la línea punteada diagonal
  marca la condición de igualdad eficacia $= a_i$.
  Puntos por encima de la diagonal indican que la mediana docente supera la
  tasa global de la UEA; puntos por debajo, lo contrario.
  \emph{Panel derecho}: coeficiente de variación (\textsc{cv}) de la eficacia
  entre los docentes que han impartido cada UEA; las barras rojas superan el
  umbral de heterogeneidad alta ($\textsc{cv} > 0.30$), indicando que la
  dificultad de la UEA depende significativamente del docente asignado.
  \textbf{En las 20~UEAs de mayor riesgo, más del 60\% presenta
  $\textsc{cv} > 0.30$, lo que demuestra que la mayor parte de su dificultad
  tiene un componente docente intervenible, no solo un componente intrínseco
  de contenido.}}
\label{fig:cobertura_critica}
\end{figure}

Tres hallazgos destacan:

\begin{enumerate}
  \item \textbf{Ciencias Básicas concentra los cuellos de botella}.
        Cálculo Diferencial, Cálculo Integral, EDO, Cinemática y Dinámica de
        Partículas y Complementos de Matemáticas tienen entre~54 y~73~docentes
        CBI históricamente asignados, todos con CV de eficacia $\in[0.38, 0.48]$
        y rango 0.03--1.0 por UEA.
        El docente que un alumno recibe en cualquiera de estas UEAs determina
        una probabilidad de aprobación que varía de~3\% a~100\% dependiendo
        del asignado.

  \item \textbf{La eficacia mediana del docente en la UEA crítica es
        sistemáticamente inferior al 65\%}.
        En seis de las ocho UEAs del Tronco General de alto riesgo, la
        eficacia mediana del docente ---calculada sobre los grupos específicos
        de esa UEA, no como promedio general del profesor--- cae entre~0.50
        y~0.64.
        Esto confirma que el riesgo de estas UEAs no es exclusivamente
        intrínseco al contenido, sino parcialmente atribuible a la práctica
        pedagógica del profesorado que las cubre habitualmente.

  \item \textbf{Casos de heterogeneidad extrema}.
        Ingeniería Geotécnica (CV$=0.93$), Comportamiento Mecánico de los
        Materiales (CV$=0.53$) y Cálculo Diferencial (CV$=0.48$) concentran
        la mayor dispersión intra-UEA.
        En Ingeniería Geotécnica, el rango de eficacia docente es~0.0--0.97:
        el resultado del alumno depende casi exclusivamente del docente
        asignado, no del contenido.
\end{enumerate}

\begin{figure}[H]\centering
\includegraphics[width=0.82\linewidth]%
  {../../Output/figuras_correlaciones/fig_28_eficacia_rendimiento.pdf}
\caption{Eficacia docente vs.\ rendimiento académico para los 738~docentes CBI,
  clasificados por cuartil de tasa de aprobación~$a_d$.
  El eje horizontal es la eficacia (fracción de estudiantes que aprueba con
  ese docente en esa UEA específica) y el eje vertical es el rendimiento
  académico mediano de los estudiantes del docente (escala~0--10 de la UAM-A).
  Los puntos pequeños son docentes individuales; los puntos grandes coloreados
  son las medianas de cada cuartil (Q1 = menor $a_d$, Q4 = mayor $a_d$).
  \textbf{El Q1 (31.3\% de los docentes CBI) se concentra en la zona de
  eficacia $< 0.65$ y rendimiento $< 5.0$, mientras que el Q4 opera con
  eficacia $> 0.82$ y rendimiento $> 7.5$; la separación entre cuartiles
  confirma que la calidad docente y el rendimiento de los estudiantes son
  indicadores coherentes y que la brecha entre el cuartil inferior y superior
  es lo suficientemente grande como para ser un objetivo de política
  institucional prioritaria.}}
\label{fig:eficacia_rendimiento}
\end{figure}

\subsection{Implicaciones para la intervención}

El análisis de profesores proporciona el nivel de resolución más fino del
diagnóstico:

\begin{itemize}[leftmargin=2em]
  \item La heterogeneidad docente explica por qué el CV de eficacia de las
        UEAs críticas supera 0.37--0.48: no es solo dispersión de cohortes
        sino \emph{dispersión de docentes asignados}.
  \item El 31\% de docentes CBI con $a_d < 0.65$ representa un potencial
        de mejora: si esos docentes alcanzaran la mediana divisional (0.744),
        la tasa de aprobación efectiva de las UEAs que cubren aumentaría
        entre~7 y~12~puntos porcentuales.
  \item Ciencias Básicas requiere intervención prioritaria no solo como
        departamento que cubre los nodos topológicamente críticos, sino
        porque es el de menor eficacia y mayor dispersión de la División
        (tabla~\ref{tab:dept_docentes}).
        La plataforma de datos construida permite identificar, en cada
        periodo, qué docentes de Ciencias Básicas imparten los nodos de
        riesgo y con qué eficacia histórica.
  \item El análisis cruzado entre el perfil docente y los resultados del plan
        revela una asociación estadísticamente significativa entre la
        heterogeneidad docente media de un plan y el retraso en el egreso:
        $r_s(\text{CV eficacia},\;\mathrm{Trim}_{50}) = +0.78$
        ($p_{\mathrm{BH}} = 0.043$, $N=10$).
        Además, los planes con mayor fracción de docentes heterogéneos
        tienden a tener también más seriaciones
        ($r_s = +0.73$, $p = 0.018$), lo que sugiere que el problema
        estructural y el factor humano no son independientes: la complejidad
        curricular y la heterogeneidad pedagógica co-ocurren en los mismos
        programas.
        Esto implica que la reforma de seriaciones, si no va acompañada de
        homogeneización docente, puede subestimar el beneficio real del cambio
        curricular.
\end{itemize}

\begin{figure}[H]\centering
\includegraphics[width=0.78\linewidth]{../../Output/figuras_correlaciones/fig_cross_het_grad.pdf}
\caption{Asociación entre la heterogeneidad docente media del plan y el retraso
  de egreso de la cohorte ($N=10$~planes de CBI).
  \emph{Eje horizontal}: coeficiente de variación (CV) medio de la eficacia
  de los docentes del plan, promediado sobre todas las UEAs del programa; mide
  cuán dispersa es la práctica docente dentro del plan.
  \emph{Eje vertical}: $\mathrm{Trim}_{50}$, el trimestre en que el 50\%
  de la cohorte simulada ha egresado (a mayor $\mathrm{Trim}_{50}$, mayor
  retraso típico en titularse).
  \emph{Tamaño}: tasa de baja del plan (fracción de estudiantes que causa
  baja definitiva por acumulación de reprobaciones).
  La banda sombreada es el intervalo de confianza al 95\% de la regresión;
  la correlación de Spearman es $r_s = +0.78$ ($p = 0.043$, corregido por
  Benjamini-Hochberg~\cite{BH1995}).
  Los planes con mayor dispersión docente exhiben cohortes sistemáticamente
  más tardías; Industrial, con el menor CV docente (0.18), presenta también
  el $\mathrm{Trim}_{50}$ más bajo de la División.
  \textbf{La heterogeneidad docente y la complejidad curricular co-ocurren en
  los mismos programas ($r_s = +0.73$ entre CV docente y número de seriaciones),
  lo que indica que la reforma de seriaciones sin homogeneización pedagógica
  simultánea subestimará el beneficio real del cambio curricular.}}
\label{fig:cross_het_grad}
\end{figure}

\subsection{Varianza docente explícita en las proyecciones de~$\ETtit$}
\label{sec:varianza_docente}

El modelo estocástico descrito en la Sección~\ref{sec:modelo} emplea~$a_i$ como
propiedad fija de cada UEA: la tasa de aprobación histórica media sobre todos
los docentes que la han impartido.
Sin embargo, el análisis de heterogeneidad docente de las secciones anteriores
muestra que las UEAs de mayor riesgo topológico presentan coeficientes de
variación (CV) de eficacia intra-UEA en el rango~$0.37$--$0.48$, con rangos de
aprobación de~$3\%$ a~$100\%$ dependiendo del docente asignado.
En esas condiciones, las proyecciones de~$\ETtit$ (tiempo esperado de egreso
condicional al egreso) y los impactos marginales de seriaciones son medias sobre
una distribución docente no modelada que puede ser sustancial.

El módulo~21 (\texttt{21\_modelo\_varianza\_docente.py}) extiende el modelo
estocástico para capturar explícitamente esta varianza.
En cada trimestre~$t$, para cada UEA~$i$, se sortea un docente
$d^{(t)}\sim\mathrm{Uniforme}$ sobre el pool histórico construido desde
\texttt{diccionario\_profesores.json}, y se sustituye la tasa de aprobación
global~$a_i$ por la tasa específica del docente asignado~$a_{d,i}$.
El mecanismo de fragilidad latente compartida~(\ref{eq:fragilidad}) opera sobre
este valor individualizado:
\begin{equation}
  a_{s,i}^{(t)}
  = \Phi\!\left(
      \Phi^{-1}(a_{d^{(t)},i})\sqrt{1+\sigma^2} \;+\; \sigma\,\theta_s
    \right),
  \label{eq:varianza_docente}
\end{equation}
donde $\sigma$ es el parámetro de fragilidad latente y $\theta_s\sim\mathcal{N}(0,1)$
es la habilidad individual del estudiante~$s$ en el trimestre~$t$.

El módulo compara~$\ETtit$ bajo cuatro condiciones de asignación docente: la
condición base (tasa media~$a_i$), asignación aleatoria del pool histórico,
docente con eficacia mediana y docente del percentil~25 del pool de la UEA.
Los intervalos de confianza (IC) al~$95\%$ sobre~$\ETtit$ que resultan de esta
comparación cuantifican la incertidumbre introducida por la distribución docente
y permiten presentar proyecciones con bandas de error explícitas en lugar de
estimaciones puntuales.
Esta extensión es relevante para evaluar el beneficio de políticas de homogeneización
pedagógica: el diferencial~$\ETtit(\mathrm{mediana}) - \ETtit(\mathrm{P25})$ estima
el margen de mejora de~$\ETtit$ alcanzable únicamente con asignación docente más
favorable, sin ninguna reforma curricular.

\begin{figure}[htbp]
  \centering
  \includegraphics[width=\textwidth]{../../Output/figuras/fig_varianza_docente.pdf}
  \caption{Varianza docente intra-UEA y su impacto potencial sobre la eficiencia terminal.
    \emph{Panel izquierdo:} distribución de tasas de aprobación del cuerpo docente
    por UEA (Unidad de Enseñanza-Aprendizaje), ordenadas por amplitud intercuartílica
    (rango P25--P75); se muestran las cien unidades con mayor heterogeneidad interna.
    \emph{Panel derecho:} diferencial~$\ETtit(\mathrm{med})-\ETtit(\mathrm{P25})$,
    donde~$\ETtit=\mathbb{E}[T\,|\,\mathrm{tit}]$ es el tiempo promedio de egreso
    condicionado al egreso (en trimestres), estimado por simulación Monte Carlo
    asignando a la UEA el docente mediano y el docente del percentil~25 de eficacia;
    la diferencia cuantifica el margen de reducción del tiempo de egreso alcanzable
    mediante asignación docente más favorable, sin reforma curricular.}
  \label{fig:varianza_docente}
\end{figure}

%══════════════════════════════════════════════════════
\section{Remoción combinada: curva dosis-respuesta}
\label{sec:dosis_respuesta}
%══════════════════════════════════════════════════════

La sección anterior evaluó el impacto de cada seriación de forma
\emph{marginal}: se elimina una sola arista y se mide el cambio en~$\ETtit$
manteniendo las demás intactas.
Ese análisis responde qué aristas vale la pena quitar; esta sección responde
qué ocurre cuando se quitan varias a la vez.

La distinción es relevante para la política curricular.
Supóngase que las tres seriaciones de mayor impacto individual ahorran
0.8, 0.5 y 0.3~trimestres respectivamente: la predicción aditiva promete
$0.8 + 0.5 + 0.3 = 1.6$~trimestres de ganancia.
¿Se cumple esa promesa si las tres se eliminan simultáneamente en la siguiente
revisión de plan?
La respuesta depende de si las aristas son causalmente independientes.
Si las tres convergen en el mismo nodo fuente del DAG (por ejemplo, todas
parten de EDO), eliminar la primera libera parte del cuello de botella y las
dos siguientes aportan un beneficio ya parcialmente capturado.
El efecto combinado será positivo pero \emph{subaditivo}: menor que la suma.

Para cuantificar esta diferencia se diseña un experimento de dosis-respuesta:
la ``dosis'' es~$N$, el número de seriaciones eliminadas simultáneamente
(siempre las de mayor impacto marginal individual), y la ``respuesta'' es la
reducción observada en~$\ETtit$.

\subsection{Diseño del experimento}

Para cada plan se ordenan las seriaciones por impacto marginal descendente
$|\Delta\ETtit|$ (ranking de la sección anterior) y se construye un
\emph{grid} adaptativo de valores de~$N$ que cubre los puntos absolutos bajos
($N = 1, 2, 3$) y los percentiles del 5\% al 100\% del total de
seriaciones del plan.
En cada punto del grid se eliminan simultáneamente las $N$~aristas de mayor
impacto marginal y se ejecutan $K = 10$~réplicas Monte Carlo independientes
($S = 10\,000$~estudiantes cada una) para estimar la ganancia combinada
\begin{equation}
  G(N) \;=\; \ETtit^{\,\text{(base)}} - \ETtit^{\,\text{(sin $N$ aristas)}}
  \label{eq:ganancia}
\end{equation}
con su intervalo de confianza al 95\% por distribución~$t$ ($\mathrm{df}=9$).

La \emph{regla de selección greedy} --- tomar siempre las $N$ aristas de mayor
impacto marginal individual --- es el mejor instrumento disponible para
priorizar qué seriaciones intervenir, pero no garantiza que el subconjunto de
tamaño~$N$ sea globalmente óptimo para la remoción conjunta.
Las consecuencias de esta limitación se discuten en
\S\ref{sec:nomonotona}.

\subsection{Modelos de saturación}
\label{sec:modelos_sat}

Las curvas $G(N)$ exhiben rendimientos decrecientes en la mayoría de los
planes: las primeras aristas removidas concentran la mayor parte del beneficio.
Para parametrizar esta saturación se ajustan dos modelos.

\noindent\textbf{Modelo de Michaelis-Menten (MM)~\cite{Michaelis1913}.}
Tomado de la cinética enzimática --- donde la velocidad de reacción satura
con la concentración de sustrato ---, se adapta aquí como
\begin{equation}
  G(N) \;=\; \frac{V_{\max}\cdot N}{K_M + N},
  \label{eq:MM}
\end{equation}
donde $V_{\max}$ es la ganancia asintótica cuando $N\to\infty$ y $K_M$ es el
número de aristas para el cual $G(K_M) = V_{\max}/2$ (semisaturación).
Es importante notar que~(\ref{eq:MM}) es un modelo de ajuste fenomenológico:
su dominio natural es $N\in[0,\infty)$, mientras que el dominio físico está
acotado por el total de seriaciones del plan ($N \leq n_{\mathrm{ser}}$).
Cuando la curva no alcanza la saturación dentro de ese dominio, el valor
$V_{\max}$ extrapolado supera la ganancia físicamente posible y carece de
interpretación directa de política.
El caso más claro es Metalurgia: el modelo estima $V_{\max} = 3.50$~trim, pero
la ganancia observada al eliminar \emph{todas} sus seriaciones es solo
$G_{\max}^{\mathrm{obs}} = 1.56$~trim.
La magnitud relevante para la política es siempre $G_{\max}^{\mathrm{obs}}$;
$V_{\max}$ y $K_M$ describen únicamente la \emph{forma} de la curva.

Para el modelo MM, el número de aristas que captura una fracción~$p$ de
$V_{\max}$ es
\begin{equation}
  N_p^{\mathrm{MM}} \;=\; \frac{K_M\,p}{1-p},
  \label{eq:breakeven}
\end{equation}
de donde $N_{80}^{\mathrm{MM}} = 4K_M$ y $N_{95}^{\mathrm{MM}} = 19K_M$.

\noindent\textbf{Modelo exponencial de saturación.}
El segundo modelo expresa la ganancia como
\begin{equation}
  G(N) \;=\; A\bigl(1 - e^{-bN}\bigr),
  \label{eq:exp_sat}
\end{equation}
donde $A$ es la ganancia asintótica y $b$ la tasa de saturación.
Este modelo es más flexible en la región $N$ pequeño y su umbral de beneficio
fraccional es
\begin{equation}
  N_p^{\mathrm{Exp}} \;=\; -\,\frac{\ln(1-p)}{b},
  \label{eq:breakeven_exp}
\end{equation}
de donde $N_{80}^{\mathrm{Exp}} = \ln 5\,/\,b \approx 1.61/b$ y
$N_{95}^{\mathrm{Exp}} = \ln 20\,/\,b \approx 3.00/b$.

El mejor modelo para cada plan se selecciona por el Criterio de Información de Akaike (AIC); tres planes prefieren
el exponencial (Química, Física y Civil).
En ambos modelos, el umbral empírico $N_{80}^{\mathrm{obs}}$ --- número de
aristas del grid para las cuales $G(N)$ alcanza el 80\% de
$G_{\max}^{\mathrm{obs}}$ --- es el parámetro de política central reportado
en la Tabla~\ref{tab:dosis_respuesta}.

\begin{figure}[H]\centering
\includegraphics[width=\linewidth]{../../Output/figuras_correlaciones/fig_sat_curvas.pdf}
\caption{Ajuste de los modelos de saturación a las curvas dosis-respuesta~$G(N)$
  por plan, donde~$G(N)=\ETtit^{\,\text{(base)}}-\ETtit^{\,\text{(sin $N$ aristas)}}$
  es la reducción en trimestres del tiempo promedio de egreso~$\ETtit=\mathbb{E}[T\,|\,\mathrm{tit}]$
  al eliminar simultáneamente~$N$ seriaciones.
  Puntos: ganancia media observada en simulación Monte Carlo ($K=10$~réplicas).
  Línea continua: modelo preferido por el Criterio de Información de Akaike~(AIC);
  rojo~$=$~Michaelis-Menten~(MM), $G(N)=V_{\max}N/(K_M+N)$;
  azul~$=$~exponencial.
  Línea discontinua: modelo alternativo.
  La línea punteada horizontal señala~$V_{\max}$ cuando supera~$G_{\max}^{\mathrm{obs}}$
  en más del 10\% (extrapolación fuera del dominio físico; caso evidente en
  Metalurgia: $V_{\max}=3.50$ frente a~$G_{\max}^{\mathrm{obs}}=1.56$~trim).
  La línea gris vertical marca~$N_{80}^{\mathrm{obs}}$, el número mínimo de
  seriaciones necesario para capturar el 80\% de~$G_{\max}^{\mathrm{obs}}$.}
\label{fig:sat_curvas}
\end{figure}

\begin{figure}[H]\centering
\includegraphics[width=\linewidth]{../../Output/figuras_correlaciones/fig_sat_params.pdf}
\caption{\emph{Izquierda:} comparación entre~$N_{80}^{\mathrm{obs}}$ (número mínimo
  de seriaciones para capturar el 80\% de la ganancia máxima observada~$G_{\max}^{\mathrm{obs}}$,
  estimado directamente en la simulación) y~$N_{80}^{\mathrm{modelo}}$ predicho
  por el ajuste de menor Criterio de Información de Akaike~(AIC).
  La discrepancia en Metalurgia ($N_{80}^{\mathrm{MM}}=319$ vs.\
  $N_{80}^{\mathrm{obs}}=50$) refleja extrapolación fuera del dominio físico:
  la curva no alcanza saturación, por lo que el modelo Michaelis-Menten~(MM)
  proyecta~$N_{80}$ a un valor irreal.
  \emph{Derecha:} diferencia de AIC entre modelos
  ($\Delta\mathrm{AIC}=\mathrm{AIC}_{\mathrm{Exp}}-\mathrm{AIC}_{\mathrm{MM}}$);
  valores positivos indican preferencia por~MM; negativos, por exponencial.
  Los umbrales convencionales~$|\Delta\mathrm{AIC}|=2$ se marcan en gris;
  tres planes prefieren el modelo exponencial (Química, Física y Civil).}
\label{fig:sat_params}
\end{figure}

\subsection{Resultados: curvas y parámetros}

\begin{figure}[H]\centering
\includegraphics[width=\linewidth]{../../Output/figuras_correlaciones/fig_11_remocion_curvas.pdf}
\caption{Curvas dosis-respuesta de remoción combinada de seriaciones para los
  diez planes de CBI.
  \emph{Eje $x$}: número~$N$ de seriaciones eliminadas simultáneamente,
  seleccionadas en orden descendente de impacto marginal individual.
  \emph{Eje $y$}: ganancia combinada $G(N) = \ETtit^{\,\text{(base)}} -
  \ETtit^{\,\text{(sin $N$ aristas)}}$ en trimestres, donde
  $\ETtit = \mathbb{E}[T\,|\,\mathrm{tit}]$ es el tiempo promedio de egreso
  condicionado al egreso; valores positivos indican acortamiento del egreso.
  La franja sombreada es el intervalo de confianza al 95\%
  ($K=10$~réplicas Monte Carlo de $S=10\,000$ estudiantes cada una).
  La línea discontinua es el mejor ajuste paramétrico (modelo de
  Michaelis-Menten o exponencial, seleccionado por AIC).
  La línea vertical punteada roja marca el umbral empírico $N_{80}^{\mathrm{obs}}$
  (número mínimo de aristas para capturar el 80\% de la ganancia máxima
  observada $G_{\max}$).
  \textbf{El patrón dominante es la saturación convexa: las primeras aristas
  removidas concentran la mayor ganancia y el beneficio marginal decrece
  rápidamente; Electrónica y Química saturan en $N < 10$, mientras Metalurgia
  y Física muestran un crecimiento cuasi-lineal que no se aplana dentro del
  dominio observable.}}
\label{fig:remocion_curvas}
\end{figure}

La Figura~\ref{fig:remocion_curvas} muestra el panorama completo de la
División.
El patrón dominante es la saturación convexa: en la mayoría de los planes la
curva crece con pendiente decreciente y la franja de IC95\% se estrecha
conforme aumenta~$N$, señal de que el experimento es estadísticamente estable.
Electrónica y Química exhiben la saturación más abrupta, mientras
que Metalurgia y Física muestran curvas más graduales que no se aplanan dentro
del dominio observable.
Dos planes (Computación y Física) presentan además tramos no-monótonos donde
$G(N) < 0$; este fenómeno se analiza en \S\ref{sec:nomonotona}.

\begin{table}[H]\centering
\caption{Parámetros del modelo dosis-respuesta por plan, ordenados por
  $G_{\max}^{\mathrm{obs}}$ descendente.
  $n_{\mathrm{ser}}$: aristas del grafo de seriaciones con impacto marginal no nulo evaluadas en la curva dosis-respuesta (subconjunto del total de seriaciones del plan).
  Modelo: ajuste seleccionado por AIC.
  $V_{\max}^{\mathrm{MM}}$ y $K_M$: parámetros del modelo de Michaelis-Menten
  $G(N)=V_{\max}N/(K_M+N)$~(\ref{eq:MM})
  (solo para planes con modelo MM; cuantifican la forma, no el máximo físico).
  $G_{\max}^{\mathrm{obs}}$: mayor ganancia observada en el grid.
  $N_{80}^{\mathrm{obs}}$: número de aristas para alcanzar el 80\% de
  $G_{\max}^{\mathrm{obs}}$ (parámetro de política).
  $\mathcal{A}_{N_{80}}$: índice de aditividad evaluado en $N_{80}^{\mathrm{obs}}$.}
\label{tab:dosis_respuesta}
\small
\setlength{\tabcolsep}{5pt}
\begin{tabular}{l r r r r r r r}
\toprule
Plan & $n_{\mathrm{ser}}$ & Modelo &
  $V_{\max}^{\mathrm{MM}}$ & $K_M$ &
  $G_{\max}^{\mathrm{obs}}$ & $N_{80}^{\mathrm{obs}}$ & $\mathcal{A}_{N_{80}}$ \\
\midrule
Electrónica    & 53 & MM  & 2.19 & 1.2  & 2.18 & \phantom{0}8 (15\%) & 0.52 \\
Metalurgia     & 50 & MM  & 3.50$^\dagger$ & 79.9 & 1.56 & 50 (100\%) & 0.89 \\
Ambiental      & 49 & MM  & 0.74 & 0.7  & 1.13 & 24 (49\%) & 0.46 \\
Química        & 54 & Exp & ---  & ---  & 1.49 & \phantom{0}5 (\phantom{0}9\%) & 0.37 \\
Eléctrica      & 49 & MM  & 0.88 & 0.5  & 1.04 & 10 (20\%) & 0.31 \\
Mecánica       & 49 & MM  & 0.43 & 0.9  & 0.63 & 15 (31\%) & 0.23 \\
Física         & 46 & Exp & ---  & ---  & 0.58 & 46 (100\%) & 0.44 \\
Civil          & 47 & Exp & ---  & ---  & 0.44 & \phantom{0}9 (19\%) & 0.61 \\
Industrial     & 51 & MM  & 0.34 & 0.2  & 0.50 & \phantom{0}8 (16\%) & 0.35 \\
Computación    & 51 & MM  & 0.18 & 0.0  & 0.30 & \phantom{0}2 (\phantom{0}4\%) & 0.86 \\
\bottomrule
\multicolumn{8}{@{}p{\dimexpr\linewidth-\tabcolsep\relax}}{\footnotesize\raggedright
  $^\dagger$ $V_{\max}^{\mathrm{MM}} = 3.50$~trim supera el máximo físico
  ($G_{\max}^{\mathrm{obs}} = 1.56$~trim) porque la curva de Metalurgia no
  satura en $N \leq 50$; la extrapolación describe la forma, no la ganancia
  alcanzable.}
\end{tabular}
\end{table}

La Tabla~\ref{tab:dosis_respuesta} debe leerse en la columna
$N_{80}^{\mathrm{obs}}$, el indicador de política:
\begin{itemize}[leftmargin=2em]
\item \textbf{Electrónica} y \textbf{Química}: $N_{80} = 8$ y $5$~aristas
  (15\% y 9\% del total del plan).
  Pocas aristas concentran el 80\% del beneficio máximo --- estructura
  Pareto pronunciada, candidatos a reforma quirúrgica.
\item \textbf{Metalurgia} y \textbf{Física}: $N_{80} = n_{\mathrm{ser}}$
  (100\% del total).
  El beneficio se distribuye de forma aproximadamente uniforme entre todas sus
  seriaciones; ninguna arista individual domina.
\item \textbf{Computación}: $N_{80} = 2$~aristas (4\%), con
  $G_{\max}^{\mathrm{obs}} = 0.30$~trim.
  La estructura topológica de este plan concentra el impacto en un par de
  aristas, pero la ganancia absoluta es muy pequeña: la restricción dominante
  en Computación es la carga crediticia histórica, no las seriaciones.
\end{itemize}

\begin{figure}[H]\centering
\includegraphics[width=\linewidth]{../../Output/figuras_correlaciones/fig_14_curvas_top3.pdf}
\caption{Curvas dosis-respuesta de remoción combinada para los tres planes con
  mayor ganancia potencial.
  El eje $x$ es el número~$N$ de seriaciones eliminadas simultáneamente
  (ordenadas por impacto marginal descendente); el eje $y$ es la ganancia
  combinada $G(N)$ en trimestres sobre $\ETtit = \mathbb{E}[T\,|\,\mathrm{tit}]$
  (tiempo promedio de egreso condicionado al egreso).
  La franja sombreada es el IC95\% ($K=10$~réplicas Monte Carlo); la línea
  discontinua es el ajuste de Michaelis-Menten (MM) o exponencial según el
  plan; el rombo rojo marca $N_{80}^{\mathrm{obs}}$ (número mínimo de aristas
  para capturar el 80\% de la ganancia máxima observada~$G_{\max}$).
  \emph{Izquierda}: Electrónica, $G_{\max}=2.18$~trim, $N_{80}=8$~aristas
  (15\% del plan).
  \emph{Centro}: Química, $G_{\max}=1.49$~trim, $N_{80}=5$~aristas (9\%).
  \emph{Derecha}: Metalurgia, $G_{\max}=1.56$~trim, $N_{80}=50$~aristas (100\%).
  \textbf{La morfología radicalmente distinta ilustra dos perfiles estratégicos:
  Electrónica y Química permiten una reforma quirúrgica acotada (pocas aristas
  capturan casi todo el beneficio), mientras Metalurgia requiere intervenir la
  totalidad del plan para obtener el máximo beneficio.}}
\label{fig:curvas_top3}
\end{figure}

La Figura~\ref{fig:curvas_top3} amplía los tres planes de mayor ganancia.
En Electrónica y Química la curva se aplana de forma nítida alrededor
de~$N_{80}$ y las réplicas convergen: cada arista adicional más allá de ese
umbral aporta menos de 0.05~trim.
En Metalurgia la curva crece de manera aproximadamente lineal hasta $N = 50$
sin mostrar saturación, lo que implica que el beneficio de la reforma crece
con el alcance de la intervención y que cualquier subconjunto de seriaciones
será subóptimo.

\subsection{Portfolio quirúrgico: aristas prioritarias para intervención}
\label{sec:portfolio}

Los planes con saturación rápida son los candidatos más favorables para una
reforma quirúrgica: un conjunto acotado de seriaciones concentra el grueso del
beneficio.
La Figura~\ref{fig:portfolio} identifica esas aristas para Electrónica
(8 aristas, $N_{80} = 15\,\text{\%}$ del plan) y Química
(5 aristas, $N_{80} = 9\,\text{\%}$), y cuantifica la brecha entre la
predicción aditiva y el efecto combinado real.

\begin{figure}[H]\centering
\includegraphics[width=\linewidth]{../../Output/figuras_correlaciones/fig_15_portfolio_quirurgico.pdf}
\caption{Portfolio de intervención quirúrgica: aristas prioritarias para los
  planes Electrónica (panel superior) y Química (panel inferior).
  \emph{Panel izquierdo de cada plan}: efectos marginales individuales
  $|\Delta\ETtit_i|$ (cambio en trimestres en $\ETtit = \mathbb{E}[T\,|\,\mathrm{tit}]$
  al eliminar esa sola seriación) de las $N_{80}$~aristas del portfolio,
  etiquetadas por el nombre de la UEA destino (la UEA que deja de requerir
  la seriación directa), ordenadas de mayor a menor impacto.
  \emph{Panel derecho de cada plan}: comparación entre la ganancia predicha por
  la suma de los efectos marginales individuales (barra gris, ``predicción aditiva'')
  y la ganancia combinada efectivamente observada al eliminar simultáneamente
  todo el portfolio (barra de color).
  $\mathcal{A} = G(N)/\sum_{i}\Delta\ETtit_i$ es el índice de aditividad
  (razón entre la ganancia combinada observada y la suma de los impactos
  marginales); $\mathcal{A} < 1$ indica subaditividad, es decir, que eliminar
  varias seriaciones juntas produce menos beneficio que la suma de sus efectos
  individuales porque comparten nodos fuente en el DAG.
  \textbf{Electrónica: las 8~aristas del portfolio suman 3.74~trim en predicción
  aditiva pero producen solo 1.93~trim combinados ($\mathcal{A}=0.52$);
  Química: las 5~aristas suman 3.37~trim aditivos frente a 1.25~trim combinados
  ($\mathcal{A}=0.37$); en ambos casos, las proyecciones sin corrección por
  subaditividad sobreestiman el beneficio real en un factor de 2--3.}}
\label{fig:portfolio}
\end{figure}

En Electrónica las tres aristas de mayor peso corresponden a dos cadenas
distintas del plan.
La primera cadena es la secuencia
Cálculo Integral~(\uea{1112029}) $\to$ EDO~(\uea{1112030}) $\to$
Circuitos Eléctricos~I~(\uea{1124001}):
los dos eslabones de esta cadena aparecen en el portfolio con impactos de
$-0.53$~trim (CI $\to$ EDO) y $-1.07$~trim (EDO $\to$ CE~I) respectivamente.
La segunda cadena es la arista
Diseño de Sistemas Electrónicos~(\uea{1123043}) $\to$ Control Digital~(\uea{1124045}),
con impacto de $-0.70$~trim.
El alto solapamiento causal ($\mathcal{A} = 0.52$) se explica por la primera
cadena: EDO actúa como \emph{nodo intermedio} compartido.
Al eliminar simultáneamente CI $\to$ EDO y EDO $\to$ CE~I, el bloqueo que EDO
imponía sobre CE~I se libera completamente desde los dos extremos,
pero el beneficio combinado es menor que la suma porque ambas aristas operan
sobre el mismo cuello de botella.

En Química el portfolio comprende la cadena
Termodinámica $\to$ Termodinámica~Aplicada $\to$ Balance de Materia $\to$ Balance de Energía
$\to$ Procesos de Separación, la cadena de mayor impacto secuencial del plan
(identificada en \S\ref{sec:efecto}).
Las cinco aristas del portfolio son los candidatos directos a revisión en la
Comisión de Planes y Programas de CBI.

\subsection{Curvas no-monótonas: por qué eliminar más no siempre ayuda}
\label{sec:nomonotona}

Computación y Física presentan tramos donde $G(N) < 0$: eliminar ciertos
subconjuntos intermedios de seriaciones empeora $\ETtit$ respecto a la línea
de base.
Los puntos no-monótonos observados en los datos son los siguientes:
Computación ($N = 33$: $G = -0.51$~trim; $N = 41$: $G = -1.38$~trim),
Física ($N = 23$: $G = -0.56$~trim; $N = 37$: $G = -0.75$~trim),
con excursiones menores también en Civil ($N = 31$, $-0.04$~trim) y Mecánica.

La interpretación correcta es la siguiente.
Las tasas de aprobación $a_i$ son constantes en el modelo: eliminar una
seriación nunca cambia la probabilidad de que un estudiante apruebe ninguna
UEA, solo cambia cuándo puede inscribirla.
Por tanto, la no-monotonicidad \emph{no} puede explicarse por estudiantes que
intentan cursos sin la preparación adecuada.
La causa real es que la selección greedy --- tomar las $N$ aristas de mayor
impacto marginal individual --- no es globalmente óptima para todo valor de~$N$.
Cuando se eliminan conjuntos intermedios de aristas, algunas combinaciones
generan configuraciones de DAG donde el flujo de estudiantes hacia las UEAs
terminales es menos eficiente que con menos eliminaciones: el grafo resultante
tiene rutas con mayor concentración de UEAs de baja aprobación que el grafo
completo.
Aumentar~$N$ más allá de esos tramos revierte la situación porque aristas
adicionales liberan caminos alternativos de mayor eficiencia.

La implicación práctica es clara: en Computación y Física no debe usarse el
ranking greedy extendido a valores arbitrarios de~$N$.
En Computación el experimento identifica $N_{\mathrm{opt}} = 2$ como el
subconjunto óptimo observado; en Física el beneficio se maximiza al remover
prácticamente todas las seriaciones.
En cualquier caso, la ganancia absoluta en ambos planes es modesta
($G_{\max}^{\mathrm{obs}} < 0.6$~trim), confirmando que la restricción
dominante es la carga crediticia, no la topología.

\subsection{Subaditividad: cuándo la suma de las partes sobreestima el todo}
\label{sec:subaditividad}

\noindent\textbf{Definición formal.}
El \emph{índice de aditividad} para un subconjunto de $N$ aristas es
\begin{equation}
  \mathcal{A}(N)
  \;=\; \frac{G(N)}{\displaystyle\sum_{i=1}^{N}\Delta\ETtit_i},
  \label{eq:aditividad}
\end{equation}
donde $G(N)=\ETtit^{\,\text{(base)}}-\ETtit^{\,\text{(sin $N$ aristas)}}$
es la ganancia combinada observada~(\ref{eq:ganancia}) y
$\sum_{i=1}^{N}\Delta\ETtit_i$ es la suma de los $N$~efectos marginales
individuales.
$\mathcal{A} = 1$ indica aditividad perfecta (los efectos no se solapan);
$\mathcal{A} < 1$ indica \emph{subaditividad}: la predicción aditiva sobreestima
el beneficio real porque las aristas comparten caminos causales aguas arriba
en el DAG; eliminar la primera arista libera parte del cuello de botella y la
segunda, que comparte ese nodo fuente, aporta un beneficio ya parcialmente
capturado.

\noindent\textbf{Universalidad de la subaditividad.}
En todos los planes y para todo~$N > 1$ la subaditividad es estadísticamente
significativa.
Los índices evaluados en $N_{80}^{\mathrm{obs}}$ oscilan entre
$\mathcal{A} = 0.23$ (Mecánica) y $\mathcal{A} = 0.89$ (Metalurgia): la suma
de los impactos marginales individuales sobreestima el beneficio combinado en
un factor de 1.1 a 4.3 según el plan.

\noindent\textbf{Variación con~$N$.}
La Figura~\ref{fig:mapa_aditividad} muestra cómo evoluciona $\mathcal{A}(N)$
para cada plan conforme aumenta la fracción de seriaciones eliminadas.

\begin{figure}[H]\centering
\includegraphics[width=\linewidth]{../../Output/figuras_correlaciones/fig_16_mapa_aditividad.pdf}
\caption{Mapa de aditividad: evolución del índice de aditividad
  $\mathcal{A}(N) = G(N)/\sum_{i=1}^{N}\Delta\ETtit_i$ para cada plan
  conforme crece la fracción de seriaciones eliminadas simultáneamente.
  Aquí $G(N)$ es la ganancia combinada observada al remover $N$~seriaciones,
  $\Delta\ETtit_i$ es el efecto marginal individual de la $i$-ésima seriación,
  y $\mathcal{A} < 1$ indica subaditividad (la suma de los efectos marginales
  sobreestima la ganancia real porque las aristas comparten nodos fuente en el DAG).
  \emph{Eje $x$}: fracción $N/n_{\mathrm{ser}}$ del total de seriaciones del
  plan removidas simultáneamente.
  \emph{Eje $y$}: plan, ordenado de mayor a menor ganancia máxima observada
  $G_{\max}^{\mathrm{obs}}$.
  \emph{Color de cada celda}: azul intenso = $\mathcal{A} \ll 1$ (subaditividad
  fuerte); blanco = $\mathcal{A} \approx 1$ (prácticamente aditivo); rojo =
  $\mathcal{A} > 1$ (superaditivo; aparece solo en $N=1$ como artefacto de
  ruido Monte Carlo, ya que por definición el efecto marginal y el combinado son
  idénticos cuando $N=1$).
  El rombo rojo en cada fila marca el umbral $N_{80}^{\mathrm{obs}}$ del plan
  (número mínimo de aristas para capturar el 80\% de $G_{\max}^{\mathrm{obs}}$).
  \textbf{La subaditividad es prácticamente universal: casi todas las celdas son
  azules y su intensidad aumenta con la fracción de aristas removidas; solo para
  fracciones pequeñas ($N/n_{\mathrm{ser}} \lesssim 0.05$) el índice se aproxima
  a la unidad, indicando que las primeras aristas removidas son suficientemente
  independientes entre sí como para que la predicción aditiva sea aproximadamente
  válida.}}
\label{fig:mapa_aditividad}
\end{figure}

Tres patrones emergen del mapa.
Primero, la subaditividad es prácticamente universal: la gran mayoría de las
celdas es azul, y la intensidad aumenta con la fracción de aristas removidas.
Segundo, para fracciones pequeñas ($N/n_{\mathrm{ser}} \lesssim 0.05$) el índice
es próximo a~1, indicando que las primeras aristas removidas son
suficientemente independientes entre sí --- es la zona donde la predicción
aditiva es más fiable.
Tercero, los planes de mayor ganancia (Electrónica, Química) presentan
subaditividad más intensa a fracciones intermedias, precisamente porque sus
aristas de alto impacto comparten nodos críticos de alto alcance (EDO,
Termodinámica~Aplicada).

\noindent\textbf{Consecuencia práctica.}
La suma de impactos marginales sobreestima el beneficio de cualquier reforma
simultánea.
El beneficio esperado de eliminar un conjunto de $N$ seriaciones es
aproximadamente $\mathcal{A}(N)\cdot\sum_{i=1}^{N}\Delta\ETtit_i$,
con el factor corrector $\mathcal{A}$ disponible en
\texttt{Output/\allowbreak R\_analisis/\allowbreak remocion\_aditividad.csv}.
El ranking greedy sigue siendo el mejor criterio para \emph{seleccionar} qué
aristas intervenir, pero toda proyección de ganancia presentada como compromiso
institucional debe escalarse por~$\mathcal{A}$.

\begin{figure}[H]\centering
\includegraphics[width=\linewidth]{../../Output/figuras_correlaciones/fig_13_remocion_resumen.pdf}
\caption{Síntesis comparativa de los resultados de remoción combinada para los
  diez planes de CBI.
  \emph{(A)}: Ganancia máxima observada $G_{\max}^{\mathrm{obs}}$ (reducción
  en trimestres del tiempo de egreso $\ETtit = \mathbb{E}[T\,|\,\mathrm{tit}]$
  al eliminar todas las seriaciones del plan), ordenada descendentemente.
  \emph{(B)}: Fracción de seriaciones necesaria para capturar el 80\% de
  $G_{\max}^{\mathrm{obs}}$ ($N_{80}^{\mathrm{obs}}/n_{\mathrm{ser}}$,
  donde $N_{80}^{\mathrm{obs}}$ es el número mínimo de aristas para alcanzar
  el 80\% del máximo y $n_{\mathrm{ser}}$ es el total de seriaciones del plan),
  en función del número total de seriaciones; la ausencia de correlación indica
  que la concentración del beneficio no depende del tamaño del plan.
  \emph{(C)}: Índice de aditividad $\mathcal{A} = G(N_{80})/\sum_{i=1}^{N_{80}}\Delta\ETtit_i$
  evaluado en $N_{80}^{\mathrm{obs}}$ (razón entre la ganancia combinada observada
  y la suma de los efectos marginales); $\mathcal{A} < 1$ indica que los efectos
  se solapan y la suma de marginales sobreestima la ganancia real; la línea
  vertical en $\mathcal{A}=1$ es la referencia de aditividad perfecta.
  \emph{(D)}: Curvas de acumulación del beneficio tipo Lorenz: el eje $x$ es la
  fracción de seriaciones removidas (ordenadas de mayor a menor impacto marginal)
  y el eje $y$ la fracción acumulada del beneficio total; las curvas sobre la
  diagonal indican dominancia Pareto (pocas aristas concentran el beneficio) y
  las que siguen la diagonal indican beneficio distribuido uniformemente.
  \textbf{El cuadro completo muestra que los planes difieren tanto en magnitud de
  la ganancia alcanzable como en la eficiencia de la intervención necesaria para
  capturarla: Electrónica y Química combinan alta ganancia y concentración extrema,
  mientras Metalurgia requiere intervenir el plan completo para obtener una ganancia
  comparable.}}
\label{fig:remocion_resumen}
\end{figure}

La Figura~\ref{fig:remocion_resumen} sintetiza los resultados en cuatro
paneles.
El panel~(A) confirma la escala de las ganancias: Electrónica lidera
(2.18~trim), seguida por Química y Metalurgia ($\approx 1.5$~trim), mientras
que Computación e Industrial presentan la ganancia potencial más baja
($< 0.5$~trim).
El panel~(B) muestra que la fracción $N_{80}/n_{\mathrm{ser}}$ no guarda
relación con el tamaño total del plan: un plan grande no es necesariamente más
fácil de reformar de forma quirúrgica.
El panel~(C) confirma que todos los índices $\mathcal{A}(N_{80})$ se sitúan
por debajo de la unidad --- subaditividad universal --- con Mecánica en el
extremo inferior (0.23) y Metalurgia en el superior (0.89).
El panel~(D) es el más informativo para la política: las curvas de Lorenz de
Electrónica y Química se alejan pronunciadamente de la diagonal, señal de que
un conjunto pequeño de aristas captura casi toda la mejora y la reforma puede
ser quirúrgica; las curvas de Metalurgia y Física se aproximan a la diagonal,
señal de que el beneficio está distribuido y la reforma efectiva requiere una
intervención mucho más amplia.

\begin{notabox}
La subaditividad es una propiedad topológica inevitable: las aristas de mayor
impacto individual convergen en los mismos nodos fuente (EDO, Termodinámica)
y sus efectos se solapan cuando se eliminan juntas.
Sumar los impactos marginales individuales sobreestima el beneficio de la
remoción conjunta en un factor de 1.1 a 4.3 según el plan.
El índice $\mathcal{A}$ corrige esta sobreestimación; debe aplicarse siempre
antes de presentar proyecciones de ganancia como compromiso institucional.
\end{notabox}

\subsection{Optimizador de portfolio para planes con curvas no-monótonas}
\label{sec:optimizador_portfolio}

La selección greedy ---tomar siempre las $N$ aristas de mayor impacto marginal
individual--- es el mejor instrumento disponible para priorizar intervenciones
en la mayor parte de los planes, pero Computación y Física presentan tramos
donde $G(N) < 0$ para subconjuntos intermedios
(\S\ref{sec:nomonotona}).
Esto demuestra que el ranking greedy puede ser contraproducente si se aplica
mecánicamente más allá de ciertos valores de~$N$, en particular en comisiones
académicas donde las decisiones se toman arista por arista sin visión del
portfolio completo.

El módulo~22 (\texttt{22\_optimizador\_portfolio.py}) implementa dos estrategias
de búsqueda alternativas para los planes con curvas no-monótonas.
La primera es \emph{greedy con backtracking} ($k=5$): si agregar la siguiente
arista del ranking produce $G < G_{\mathrm{anterior}}$, se exploran las
siguientes~$k=5$ candidatas antes de decidir, evitando así los valles locales de
la curva sin el costo computacional de la búsqueda global.
La segunda es \emph{simulated annealing} sobre el espacio de subconjuntos, con
temperatura inicial calibrada para aceptar intercambios de
$\Delta G = -0.1$~trim con probabilidad~$0.5$ y enfriamiento multiplicativo
de factor~$0.95$.

Para Computación, el subconjunto óptimo observado es $N_{\mathrm{opt}}=2$~aristas;
para Física, el beneficio se maximiza al remover prácticamente la totalidad del
plan ($N_{\mathrm{opt}}\approx n_{\mathrm{ser}}$).
En ambos casos la ganancia absoluta es modesta
($G_{\max}^{\mathrm{obs}}<0.6$~trim), lo que confirma que la restricción
dominante es la carga crediticia histórica, no la topología, y que la prioridad
de reforma en estos planes no son las seriaciones sino el frente de carga.

\begin{figure}[H]\centering
\includegraphics[width=0.85\linewidth]{../../Output/figuras/fig_portfolio_optimo.pdf}
\caption{Comparación de estrategias de optimización de portfolio en los cuatro planes
  de referencia (Computación y Física: no-monótonos; Electrónica y Química: referencia).
  Para un presupuesto de~$N=10$ aristas, el \emph{greedy con backtracking} (BT)
  supera al greedy puro en todos los planes; el \emph{simulated annealing} (SA) es
  competitivo con BT en los planes no-monótonos y ligeramente inferior en los demás.
  La ganancia absoluta confirma la asimetría: en los planes de referencia (Electrónica
  y Química) la topología aporta 1.3--1.9~trim; en los no-monótonos (Computación y
  Física), menos de~0.7~trim.}
\label{fig:portfolio_optimo}
\end{figure}

\subsection{Modelo predictivo del índice de subaditividad}
\label{sec:modelo_aditividad}

El índice de aditividad $\mathcal{A}$ (ecuación~\ref{eq:aditividad}) varía
entre~$0.23$ y~$0.89$ entre planes sin que la sola inspección de la
Tabla~\ref{tab:dosis_respuesta} revele qué estructura topológica determina
si una reforma conjunta será muy subaaditiva ($\mathcal{A}\approx0.23$) o
casi aditiva ($\mathcal{A}\approx0.89$).
Sin esa comprensión, el factor corrector~$\mathcal{A}$ es empírico: aplica
a los currículos actuales pero no tiene capacidad predictiva para currículos
modificados, precisamente donde más se necesita.

El módulo~23 (\texttt{23\_modelo\_aditividad.R}) construye un modelo de regresión
de~$\mathcal{A}(N)$ sobre cinco predictores topológicos del subconjunto de aristas
eliminadas:
\begin{equation}
  \mathcal{A}(N)
  \;\approx\; f\!\bigl(
      f_{\mathrm{src}},\;
      f_{\mathrm{dst}},\;
      \mathrm{diam},\;
      \mathrm{div}_{\mathrm{src}},\;
      \overline{\mathrm{betw}}_{\mathrm{src}}
    \bigr),
  \label{eq:modelo_aditividad}
\end{equation}
donde $f_{\mathrm{src}}$ es la fracción de aristas del subconjunto que comparten
nodo fuente, $f_{\mathrm{dst}}$ la que comparte nodo destino, $\mathrm{diam}$
el diámetro del subgrafo inducido por el subconjunto, $\mathrm{div}_{\mathrm{src}}
= |\text{src únicos}|/|S|$ la diversidad de fuentes y
$\overline{\mathrm{betw}}_{\mathrm{src}}$ la centralidad de intermediación media
de los nodos fuente en el grafo acíclico dirigido (DAG) completo.
El modelo se ajusta sobre los puntos de las curvas dosis-respuesta con ganancia
combinada positiva y ratio acotado (97~puntos de los 120 totales; los restantes
corresponden a tramos no-monótonos donde el ratio es inestable).
El ajuste arroja $R^2=0.22$ (RMSE validación cruzada por plan~$=0.22$):
los predictores significativos son~$f_{\mathrm{dst}}$ ($t=2.86$, positivo,
las aristas que convergen al mismo destino interfieren menos entre sí) y
$\overline{\mathrm{depth}}_{\mathrm{src}}$ ($t=-2.13$, negativo, nodos fuente
más profundos producen mayor subaditividad).
El $R^2$ obtenido no alcanza el umbral de~$0.60$ necesario para aplicación
predictiva directa; el modelo describe la tendencia estructural pero no sustituye
a la simulación completa para currículos reformados.

\begin{figure}[H]\centering
\includegraphics[width=\linewidth]{../../Output/figuras_correlaciones/fig_aditividad_modelo.pdf}
\caption{Modelo predictivo del índice de aditividad~$\mathcal{A}(N)=G(N)/\sum_{i=1}^{N}\Delta\ETtit_i$,
  donde~$G(N)$ es la ganancia combinada al remover~$N$ seriaciones
  y~$\Delta\ETtit_i$ el efecto marginal individual de la $i$-ésima arista;
  $\mathcal{A}<1$ indica subaditividad (los efectos se solapan).
  \emph{Izquierda:} valores observados de~$\mathcal{A}(N)$ frente a los predichos
  por regresión de Mínimos Cuadrados Ordinarios~(OLS, por sus siglas en inglés)
  con cinco predictores topológicos del subconjunto; varianza explicada~$R^2=0.22$,
  color por plan.
  \emph{Derecha:} trayectorias observadas de~$\mathcal{A}(N)$ para los diez planes
  en función de la fracción de seriaciones eliminadas; la dispersión inter-plan
  es la varianza que el modelo lineal trata de capturar.}
\label{fig:aditividad_modelo}
\end{figure}

\subsection{Clasificación estratégica de los planes}
\label{sec:clasificacion}

Los resultados anteriores permiten posicionar cada plan en un espacio
bidimensional definido por dos indicadores de política: la magnitud del
beneficio alcanzable ($G_{\max}^{\mathrm{obs}}$, eje vertical) y la eficiencia
de la intervención ($N_{80}^{\mathrm{obs}}$ como porcentaje del total de
seriaciones del plan, eje horizontal).
El cuadrante superior-izquierdo agrupa a los planes de mayor retorno con
menor número de cambios curriculares.

\begin{figure}[H]\centering
\includegraphics[width=\linewidth]{../../Output/figuras_correlaciones/fig_17_clasificacion_planes.pdf}
\caption{Clasificación estratégica de los diez planes de CBI según el
  beneficio alcanzable y la eficiencia de la intervención requerida.
  \emph{Eje $x$}: $N_{80}^{\mathrm{obs}}$ expresado como porcentaje del
  total de seriaciones del plan; posición más a la izquierda indica que
  pocas aristas bastan para capturar el 80\% del beneficio máximo
  (intervención más quirúrgica).
  \emph{Eje $y$}: ganancia máxima observada $G_{\max}^{\mathrm{obs}}$ en
  trimestres (reducción en $\ETtit = \mathbb{E}[T\,|\,\mathrm{tit}]$
  al eliminar todas las seriaciones del plan); posición más alta indica
  mayor beneficio potencial.
  \emph{Tamaño del símbolo}: número total de seriaciones del plan.
  \emph{Color}: cuadrante estratégico en el que cae el plan.
  El rótulo itálico bajo cada punto es el índice de aditividad
  $\mathcal{A}(N_{80})$ (razón entre la ganancia combinada y la suma de
  efectos marginales; valores cercanos a 1 indican menor solapamiento entre
  aristas).
  Las líneas discontinuas son las medianas de cada eje sobre los diez planes.
  \textbf{El cuadrante superior izquierdo (alta ganancia + intervención
  concentrada) es el de mayor retorno institucional: Electrónica y Química
  se ubican aquí, con ganancias de 2.18 y 1.49~trimestres alcanzables con
  solo 8 y 5~seriaciones respectivamente, lo que los convierte en los
  candidatos prioritarios para una reforma quirúrgica de los planes 2020.}}
\label{fig:clasificacion}
\end{figure}

\textbf{Electrónica} y \textbf{Química} (cuadrante
\emph{quirúrgico}): ganancia alta ($G_{\max} > 1.4$~trim) con
$N_{80} < 16\,\text{\%}$ del total.
Son los candidatos prioritarios de la División para reforma focalizada.

\textbf{Ambiental}, \textbf{Eléctrica} y \textbf{Metalurgia} (cuadrante de
\emph{reforma amplia}): ganancias sustanciales (1.04--1.56~trim) pero
$N_{80} \geq 20\,\text{\%}$, lo que exige intervenir una fracción considerable
de las seriaciones.
En Metalurgia, $N_{80} = 100\,\text{\%}$ y la disponibilidad de la estancia
industrial añade una restricción extracurricular al alcance de la reforma.

\textbf{Civil}, \textbf{Computación} e \textbf{Industrial} (cuadrante
\emph{focal}): ganancia baja (0.30--0.50~trim) con $N_{80} < 20\,\text{\%}$.
En Computación, dos aristas concentran el máximo beneficio topológico
($G_{\max} = 0.30$~trim, $N_{80} = 4\,\text{\%}$); la restricción dominante
es la carga crediticia histórica.

\textbf{Física} y \textbf{Mecánica} (cuadrante \emph{distribuido}): ganancia
baja (0.58--0.63~trim) con $N_{80} \geq 31\,\text{\%}$.
En estos planes la inversión en reforma topológica produce retornos menores que
en los cuadrantes quirúrgico o de reforma amplia; la prioridad debería
orientarse a la carga académica y al apoyo en UEAs críticas.

%══════════════════════════════════════════════════════
\section{Índices de prelación: incentivos de inscripción y penalización topológica}
\label{sec:indice_desempenio}
%══════════════════════════════════════════════════════

El acceso a los grupos de mayor demanda no es aleatorio: lo ordena un índice
de desempeño cuya evolución a lo largo de la trayectoria refleja, de manera
no evidente, la estructura de seriaciones del plan.
Esta sección examina el índice vigente~$P$ y la propuesta alternativa~$h$
de forma paralela ---primero la mecánica de cada uno, luego los incentivos de
inscripción que generan--- y concluye con una comparación directa de su
desempeño bajo las mismas cohortes simuladas y su respuesta común a la
restricción topológica del grafo de seriaciones.

%──────────────────────────────────────────────────────
\subsection{El índice~\texorpdfstring{$P$}{P} vigente: fórmula y déficit crónico de ritmo}
\label{sec:P_vigente}

El acuerdo~551.4.4.1 del Consejo Divisional (febrero de~2015) establece el
procedimiento de cálculo del índice de desempeño~$P$ y su aplicación en el
proceso de reinscripción, vigente desde el trimestre~15-P.
Cada solicitud de inscripción a un grupo ofertado por la División se ordena de
mayor a menor según~$P$; cuando dos o más estudiantes empatan, el desempate
primario es el grado de avance~$\mathrm{GA}$ y el secundario el
aprovechamiento del último trimestre~$\mathrm{ACT}$; si persiste el empate se
asigna de forma aleatoria.
El índice se define como
\begin{equation}
  P = a(\mathrm{RA}) + b(\mathrm{GA}) + c(\mathrm{PN})
      + d(\mathrm{ACT}) + e(\mathrm{PTN}),
  \quad 0 \le P \le 1,
  \label{eq:indice_P}
\end{equation}
donde $\mathrm{RA}$~es el ritmo de avance, $\mathrm{GA}$~el grado de avance,
$\mathrm{PN}$~el promedio de calificaciones general normalizado,
$\mathrm{ACT}$~el aprovechamiento de los créditos inscritos en el trimestre
lectivo inmediato anterior y $\mathrm{PTN}$~el promedio del mismo trimestre
normalizado, con pesos fijos $c=0.1$, $d=0.4$ y $e=0.1$ para los componentes
de trayectoria global y de último trimestre respectivamente.
Los pesos~$a$ y~$b$ son variables en función de los créditos contabilizados
acumulados $\mathrm{CC}$:
\begin{equation}
  a =
  \begin{cases}
    0.35 & \text{si } \mathrm{CC} < 100, \\
    0.30 & \text{si } 100 \le \mathrm{CC} < 250, \\
    0.25 & \text{si } 250 \le \mathrm{CC} < 350, \\
    0.05 & \text{si } \mathrm{CC} \ge 350,
  \end{cases}
  \qquad
  b = 0.40 - a,
  \label{eq:pesos_ab}
\end{equation}
de modo que $(a+b)+c = d+e = 0.5$: dado que $c=0.10$ es
constante, se cumple $a+b=0.40$ para todo valor de~$CC$, y
la suma de los pesos de la trayectoria acumulada $(a+b+c=0.50)$
iguala la de los indicadores del último trimestre $(d+e=0.50)$.
Los componentes individuales son
\begin{equation}
  \mathrm{RA} = \min\!\left(\frac{C_A}{38\,T_T},\,1\right),
  \quad
  \mathrm{GA} = \frac{\mathrm{CC}}{C_P},
  \quad
  \mathrm{ACT} = \frac{C_{A,T-1}+C_{A,T-2}}{C_{I,T-1}+C_{I,T-2}},
  \label{eq:componentes_P}
\end{equation}
donde $C_A$~son los créditos aprobados acumulados desde el ingreso, $T_T$~el
número de trimestres lectivos transcurridos, $C_P$~los créditos mínimos del
plan de estudios, $C_{A,T-1}$~y~$C_{I,T-1}$~los créditos aprobados e
inscritos en la evaluación global del trimestre inmediato anterior,
y~$C_{A,T-2}$, $C_{I,T-2}$~los correspondientes a evaluaciones de
recuperación o especial en el trimestre anterior a ese.
El ritmo normal de referencia se fija en~$38$~cr/trim, correspondiente a una
trayectoria de trece trimestres con carga promedio de~$C_P/13\approx38$~cr/trim.
Las calificaciones se normalizan al intervalo~$[0,1]$ aplicando la
equivalencia~MB$\to10$, B$\to8$, S$\to6$ (acuerdo~91.8 del Colegio Académico).

Nótese que la fórmula de~$\mathrm{ACT}$ integra el resultado de la evaluación
de recuperación del trimestre~$T-2$: un estudiante que aprueba en recuperación
mejora directamente su~$\mathrm{ACT}$ y con ello su~$P$.
Esto constituye un incentivo institucional ya existente a presentar la
recuperación, lo que es coherente con el hallazgo de que CBI es la única
división con tasa neta de recuperación positiva
($\delta a = +0.015$; véase~\S\ref{sec:dimensiones_adicionales}).

\subsection{Desajuste estructural de ritmo y doble penalización topológica de~\texorpdfstring{$P$}{P}}
\label{sec:RA_cronico}

El denominador de~$\mathrm{RA}$ fija en~$38$~cr/trim el criterio de avance
normal.
La calibración del modelo topológico-estocástico
(\S\ref{sec:modelo}) revela que la carga efectiva histórica oscila
entre~$c_s=27$~cr/trim en Mecánica y~$c_s=41$~cr/trim en Metalurgia;
para los ocho planes restantes se concentra en el rango~$30$--$33$~cr/trim.
Con la única excepción parcial de Metalurgia
($c_s=41 > 38$), todos los planes operan en el régimen $c_s < 38$, de modo
que la razón~$C_A/(38\,T_T)$ satisface~$\mathrm{RA}<1$ de manera estructural
para la práctica totalidad de la cohorte.

La simulación del modelo con parámetros calibrados cuantifica esta
discrepancia: la fracción de estudiantes con~$\mathrm{RA}<1$ es igual a~$1.000$
en nueve planes y~$0.986$ en Metalurgia durante toda la ventana de
observación~$T\in[2,24]$~trim (Figura~\ref{fig:RA_cronico}).
\textbf{El índice de desempeño vigente coloca estructuralmente a la totalidad
del alumnado en situación de déficit de ritmo, no porque los estudiantes
avancen mal sino porque el denominador de~$\mathrm{RA}$ está calibrado a un
parámetro de diseño ---13~trimestres a~38~cr/trim--- que el propio documento
oficial reconoce como inalcanzable para la mayoría de la cohorte desde~1987.}

\begin{figure}[H]\centering
\includegraphics[width=0.92\linewidth]{../../Output/figuras_correlaciones/fig_32b_RA_cronico.pdf}
\caption{Déficit crónico de Ritmo de Avance por plan.
  El Ritmo de Avance se define como~$\mathrm{RA}=\min(C_A/(38\cdot T),\,1)$,
  donde~$C_A$ son los créditos acumulados,~$T$ el número de trimestres cursados
  y~$38$~cr/trim el ritmo establecido en el plan de estudios; $\mathrm{RA}<1$
  significa que el estudiante acumula créditos más lento de lo establecido.
  Las curvas muestran la fracción de la cohorte simulada con~$\mathrm{RA}<1$
  en cada trimestre, por plan.
  En todos los planes salvo Metalurgia ($c_s=41>38$~cr/trim) más del~98\% de
  los estudiantes tienen~$\mathrm{RA}<1$ desde el tercer trimestre.
  La línea discontinua marca el umbral del~50\%.
  El déficit es universal y estructural: la carga efectiva histórica~$c_s$
  cae por debajo de~38~cr/trim porque las seriaciones bloquean la inscripción
  de créditos disponibles, no por un déficit individual del alumnado.}
\label{fig:RA_cronico}
\end{figure}

\paragraph{La doble penalización topológica.}
\label{sec:doble_penalizacion}

Las cadenas de seriación generan una penalización compuesta sobre el índice~$P$
que opera a través del componente~$\mathrm{GA}$.
El mecanismo es el siguiente: cuando un estudiante no supera una UEA que actúa
como prerrequisito de un subconjunto de UEAs del plan, los créditos de ese
subconjunto no se acumulan en~$C_A$ durante los trimestres de bloqueo;
$\mathrm{GA}=C_A/C_P$ decrece por debajo de lo que el avance neto en otras
ramas del grafo aportaría; y el índice~$P$ resultante, dominado en etapas
avanzadas por~$b\,\mathrm{GA}$ con~$b=0.35$ (véase~(\ref{eq:pesos_ab})),
se deprime precisamente cuando la necesidad de acceso preferente a los grupos
del cuello de botella es mayor.
La seriación no solo retardal egreso ---efecto cuantificado en las
secciones precedentes--- sino que adicionalmente degrada el instrumento que
regula el acceso al grupo donde el problema podría resolverse.

La simulación del modelo para el plan de \textbf{Ingeniería Electrónica}
---el de mayor impacto topológico sobre~$\ETtit$--- permite aislar este efecto
con precisión.
Tomando como cuello de botella la UEA Cálculo Integral (\uea{1112029}),
inicio de la cadena Cálculo Integral~$\to$ EDO~$\to$ Circuitos Eléctricos~I
que concentra el~65\% del beneficio del portfolio quirúrgico del plan
(\S\ref{sec:clasificacion}), se clasifica a los~$3\,000$ estudiantes simulados
en dos grupos en cada trimestre: los que aún no han acreditado Cálculo Integral
(bloqueados) y los que ya la han superado (desbloqueados).
En el trimestre~12, los estudiantes bloqueados registran un índice~$P$ medio
de~$0.445$ frente a~$0.655$ de los desbloqueados ---una brecha de~$0.21$
puntos--- con una diferencia análoga en el grado de avance:
$\mathrm{GA}_{\mathrm{bloq}}=0.350$ frente a~$\mathrm{GA}_{\mathrm{nobl}}=0.572$.
\textbf{El bloqueo en la primera UEA de la cadena deprime el índice~$P$ en más
de~20~puntos porcentuales respecto al grupo que logró superarla, justo en el
período en que el acceso preferente a los grupos de Cálculo Integral sería más
valioso para revertir el rezago.}

\begin{figure}[H]\centering
\includegraphics[width=\linewidth]{../../Output/figuras_correlaciones/fig_32e_bloqueados.pdf}
\caption{Doble penalización topológica del bloqueo curricular.
  Se comparan trayectorias medias del índice de desempeño~$P$
  (Acuerdo~551.4.4.1-2015, UAM-A), del Grado de Avance~$\mathrm{GA}=C_A/C_P$
  (fracción de créditos del plan acreditados) y del índice alternativo
  entrópico~$h=(\eta_t+\langle\eta\rangle)/2$ (donde~$\eta_t=C_{A,t-1}/45$ es
  la fracción del máximo inscribible de créditos aprobada el trimestre anterior)
  para dos grupos del plan Electrónica simulado con~$S=3\,000$ estudiantes:
  aquellos cuya cadena de seriaciones pasa por Cálculo Integral como cuello de
  botella (bloqueados, línea discontinua roja) y quienes acceden sin restricción
  a las UEAs de los trimestres siguientes (desbloqueados, línea continua azul).
  La brecha en~$P$ y en~$\mathrm{GA}$ se abre desde el quinto trimestre y
  alcanza~$\Delta P\approx0.21$~puntos en~$T=12$: el bloqueo deprime el índice
  directamente vía~$\mathrm{GA}$ e indirectamente al reservar acceso preferente
  a la UEA bloqueante a quienes ya la han acreditado.}
\label{fig:bloqueados}
\end{figure}

La magnitud de esta doble penalización varía entre planes según la profundidad
topológica de sus cuellos de botella.
La Figura~\ref{fig:doble_penalizacion} muestra el incremento en~$P$ que se
obtiene al eliminar la totalidad de las seriaciones de cada plan
($\Delta P = P(\text{sin seriaciones}) - P(\text{baseline})$) en los
trimestres~$T\in\{8,12,16\}$.
Los planes de perfil quirúrgico ---Electrónica y Química--- son los que
experimentan la mayor recuperación de~$P$ al eliminar las seriaciones, en
correspondencia con el mayor impacto topológico que sus cadenas ejercen sobre
$\ETtit$ (véase la clasificación estratégica de~\S\ref{sec:clasificacion}).

\begin{figure}[H]\centering
\includegraphics[width=0.9\linewidth]{../../Output/figuras_correlaciones/fig_32c_doble_penalizacion.pdf}
\caption{Magnitud de la penalización topológica sobre el índice de desempeño~$P$
  (Acuerdo~551.4.4.1-2015, UAM-A).
  Cada barra muestra~$\Delta P=P(\text{sin seriaciones})-P(\text{base})$, la
  ganancia en el índice al eliminar la totalidad de las seriaciones del plan,
  evaluada en tres trimestres representativos ($T\in\{8,12,16\}$).
  $\Delta P>0$ significa que las cadenas de prerequisito deprimen~$P$ respecto
  al escenario sin restricciones topológicas; el canal principal es el Grado de
  Avance~$\mathrm{GA}=C_A/C_P$ (fracción de créditos del plan acreditados):
  las seriaciones reducen~$C_A$ disponible y arrastran~$\mathrm{GA}$ hacia abajo.
  Los planes Electrónica y Química, de perfil quirúrgico con cadenas largas y
  concentradas, acumulan la mayor penalización.
  Los planes con restricción dominante de carga (Civil, Computación, Industrial)
  muestran incrementos menores porque el cuello de botella es la carga global
  $c_s$, no la topología.}
\label{fig:doble_penalizacion}
\end{figure}

%──────────────────────────────────────────────────────
\subsection{Estrategias de inscripción bajo~\texorpdfstring{$P$}{P}: la trampa ACT}
\label{sec:trampa_ACT}

La dependencia del índice~$P$ del componente~$\mathrm{ACT}$ ---con peso~$d=0.40$,
el mayor de todos--- suscita una pregunta de política: ¿puede el alumnado
incrementar~$P$ de manera sostenida mediante la inscripción estratégica de pocas
UEAs de alta~$a_i$?
De confirmarse ese incentivo, el parámetro~$c_s\approx30$~cr/trim observado
históricamente sería al menos parcialmente explicable como respuesta racional al
diseño del propio índice.
Para responder, se simulan tres estrategias sobre la misma cohorte~($S=3\,000$
estudiantes/plan, mismo vector~$\theta_s$), de modo que la comparación es
dentro-estudiante.

\paragraph{Definición de las tres estrategias.}
\label{sec:estrategias_definicion}

Las tres estrategias difieren en la carga máxima inscrita por trimestre y en el
criterio de selección de UEAs cuando hay varias elegibles simultáneamente.
La \emph{estrategia agresiva} impone el límite establecido~$c^\ast=38$~cr/trim
y selecciona las UEAs en orden topológico ascendente de profundidad en el grafo
de seriaciones, replicando el escenario al que aspira el diseño curricular de~13
trimestres.
La \emph{estrategia baseline} utiliza la carga calibrada~$c_s$ de cada plan
---obtenida del ajuste al anclaje histórico~$\mathrm{E}[T\mid\text{tit}]$--- con
el mismo orden topológico; corresponde al comportamiento efectivo histórico y
sirve de referencia.
La \emph{estrategia conservadora} fija~$c=24$~cr/trim (~$3$ UEAs de~$8$~cr
en promedio) y selecciona las UEAs elegibles en orden descendente de tasa de
aprobación~$a_i$, maximizando la probabilidad de que la UEA inscrita sea aprobada
en ese trimestre y, con ello, el numerador de~$\mathrm{ACT}$ en el trimestre
siguiente.
La estrategia conservadora encarna el incentivo ACT en su forma más pura:
inscribir menos, pero inscribir lo más fácil disponible.

\paragraph{El horizonte de ventaja en~$P$.}
\label{sec:horizonte_ventaja}

La Figura~\ref{fig:ventaja_P} muestra~$\Delta P(t) = P_\mathrm{cons}(t) - P_\mathrm{base}(t)$
como función del trimestre, desagregada por plan.
En el trimestre~$T=2$, la estrategia conservadora exhibe una ventaja media de
$\Delta P(2)=+0.111$ puntos sobre el baseline, atribuible íntegramente al componente
$\mathrm{ACT}$: en~$T=1$ los estudiantes conservadores inscribieron menos créditos
pero con mayor probabilidad de aprobación, de modo que su tasa~$\mathrm{ACT}$
en~$T=2$ supera en~$0.268$ puntos a la del baseline.
Sin embargo, esta ventaja se agota con rapidez: en~$T=4$ la diferencia media es
$\Delta P(4)=-0.013$, y para~$T=8$ ha caído a~$\Delta P(8)=-0.070$.
El cruce de cero ocurre en~$T\approx2.5$, consistentemente en todos los planes
(Figura~\ref{fig:ventaja_P}).

El mecanismo de extinción de la ventaja tiene dos causas concurrentes.
En primer lugar, la estrategia conservadora agota en pocas UEAs el conjunto de
materias de alta~$a_i$ que no requieren prerrequisitos; a partir de~$T=3$,
las UEAs elegibles son estructuralmente las mismas que en el baseline, pues la
topología del grafo impone el mismo cuello de botella.
En segundo lugar, la restricción de carga~$c=24$~cr/trim acumula créditos más
lentamente que el baseline ($c_s\in\{27\text{--}41\}$~cr/trim), de modo que
$\mathrm{GA}=C_A/C_P$ diverge progresivamente:
$\Delta\mathrm{GA}(12)=-0.210$ puntos en promedio.
Dado que el peso~$b$ de~$\mathrm{GA}$ crece con~$C_A$ hasta alcanzar~$b=0.35$
para~$C_A\geq350$~cr, el déficit en grado de avance domina
sobre la ventaja residual en~$\mathrm{ACT}$ a partir del cuarto trimestre.

\begin{figure}[H]\centering
\includegraphics[width=0.96\linewidth]{../../Output/figuras_correlaciones/fig_33d_ventaja_P.pdf}
\caption{Incentivo al subregistro del componente~$\mathrm{ACT}$ en el índice de desempeño~$P$
  (Acuerdo~551.4.4.1-2015, UAM-A).
  Se representa la diferencia~$\Delta P(t)=P_{\mathrm{cons}}(t)-P_{\mathrm{base}}(t)$
  entre dos estrategias de inscripción simuladas: la \emph{conservadora}
  (24~cr/trim, comenzando por las UEAs de mayor tasa de aprobación~$a_i$) y el
  \emph{baseline} (carga histórica~$c_s$ calibrada, en orden topológico).
  Líneas finas coloreadas: cada uno de los diez planes; línea gruesa azul: media.
  La zona sombreada roja delimita la ventana~$T\in[1,2.5]$ en que la estrategia
  conservadora supera al baseline gracias a que el componente
  $\mathrm{ACT}=C_{A,T-2}/C_{I,T-2}$ (razón de créditos aprobados sobre inscritos
  en los dos últimos trimestres) se eleva al reducir el denominador.
  A partir de~$T\approx3$, el Grado de Avance~$\mathrm{GA}=C_A/C_P$ (fracción de
  créditos del plan ya acreditados) domina: la menor acumulación de~$C_A$ penaliza
  $P$ de forma sostenida y permanente en todos los planes.}
\label{fig:ventaja_P}
\end{figure}

\paragraph{La trampa: mayor~$P$ a corto plazo, menor tasa de egreso a largo plazo.}
\label{sec:trampa_titulacion}

La Figura~\ref{fig:P_estrategias} muestra la trayectoria completa del índice~$P$
bajo las tres estrategias para cuatro planes representativos.
La estrategia conservadora produce un~$P$ superior durante los primeros dos
trimestres y cae por debajo del baseline desde~$T=3$, sin recuperarse en ningún
plan ni en ningún trimestre posterior.
La estrategia agresiva, en cambio, exhibe la trayectoria de~$P$ más alta a partir
de~$T=6$: la mayor velocidad de acumulación de créditos eleva~$\mathrm{GA}$ y
$\mathrm{RA}$ de manera que supera la penalización por~$\mathrm{ACT}$ moderado.

\begin{figure}[H]\centering
\includegraphics[width=0.96\linewidth]{../../Output/figuras_correlaciones/fig_33a_P_estrategias.pdf}
\caption{Trayectorias del índice de desempeño~$P$ (Acuerdo~551.4.4.1-2015, UAM-A)
  para cuatro planes representativos bajo tres estrategias de inscripción.
  Las curvas muestran la media con intervalo intercuartílico (IQR, rango P25--P75,
  zona sombreada) obtenidos sobre~$S=3\,000$ estudiantes/plan.
  Las estrategias son: \emph{agresiva} (línea punteada azul claro: inscripción al
  tope establecido de~$c^\ast=38$~cr/trim, seleccionando UEAs en orden topológico),
  \emph{baseline} (línea sólida azul oscuro: carga histórica~$c_s$ calibrada, orden
  topológico) y \emph{conservadora} (línea discontinua roja: 24~cr/trim, comenzando
  por las UEAs de mayor tasa de aprobación~$a_i$).
  La conservadora supera al baseline en~$T=1$--$2$ gracias al componente~$\mathrm{ACT}$
  y cae por debajo de manera permanente desde~$T=3$.
  La agresiva alcanza los valores más altos de~$P$ a partir de~$T=6$--$8$.
  La línea vertical gris marca el cruce~$T\approx2.5$.}
\label{fig:P_estrategias}
\end{figure}

La paradoja de la trampa ACT queda expuesta con precisión en la
Figura~\ref{fig:trampa_scatter}: la estrategia conservadora presenta un~$P$
medio en~$T=8$ comparable al del baseline en varios planes, pero la tasa de
egreso cae a~$0.207$ en promedio (rango: $0.08$ en Electrónica, $0.29$ en
Física), frente a~$0.596$ del baseline y~$0.628$ de la agresiva.
El tiempo promedio al egreso condicionado a egresar se extiende a
$\mathrm{E}[T\mid\text{tit}]_\mathrm{cons}=24.4$~trimestres ---más de tres años
adicionales respecto al baseline--- precisamente porque la lentitud acumulada
genera nuevos bloqueos en las cadenas de seriación que el estudiante no puede
sortear dentro de la ventana~$T_\mathrm{max}=26$~trim.

\begin{figure}[H]\centering
\includegraphics[width=0.80\linewidth]{../../Output/figuras_correlaciones/fig_33c_trampa_ACT_scatter.pdf}
\caption{La trampa ACT: mayor índice~$P$ (Acuerdo~551.4.4.1-2015) a mitad de carrera
  no garantiza un egreso más rápido.
  El eje horizontal es el valor medio del índice~$P$ en el trimestre~8; el eje
  vertical es el tiempo promedio de egreso condicionado al egreso
  $\mathrm{E}[T\mid\text{tit}]$ (en trimestres; a mayor valor, mayor retraso).
  Cada punto es un plan, bajo una de las tres estrategias simuladas:
  \emph{agresiva} (triángulos azul claro: $c^\ast=38$~cr/trim, orden topológico),
  \emph{baseline} (círculos azul oscuro: carga histórica~$c_s$, orden topológico)
  y \emph{conservadora} (cuadrados rojos: 24~cr/trim, UEAs de mayor~$a_i$ primero).
  La estrategia conservadora exhibe~$P(T=8)$ comparable al baseline en varios planes
  pero egresa entre~$4$ y~$5$~trimestres más tarde: el componente~$\mathrm{ACT}$
  de~$P$ puede inflarse reduciendo el denominador sin que ello acelere el egreso.
  Metalurgia se excluye porque en ese plan $c^\ast=38 < c_s=41$~cr/trim, de modo
  que ningún estudiante egresa dentro del horizonte de~26~trimestres bajo las
  estrategias agresiva y conservadora.}
\label{fig:trampa_scatter}
\end{figure}

\paragraph{Respuesta a la pregunta causal.}
\label{sec:respuesta_causal}

Los resultados de la simulación permiten responder con precisión a la pregunta
planteada al inicio de esta sección.
El componente~$\mathrm{ACT}$ del índice de desempeño vigente crea un incentivo
al subregistro estratégico que es cuantitativamente real en~$T=1$--$2$
($\Delta P\approx+0.11$), pero estructuralmente autolimitante: la topología del
grafo de seriaciones agota el inventario de UEAs de alta~$a_i$ no bloqueadas
en las primeras inscripciones, y el déficit acumulado en grado de avance~$\mathrm{GA}$
anula y revierte la ventaja antes del cuarto trimestre.
\textbf{El parámetro~$c_s\approx30$~cr/trim observado históricamente no puede
atribuirse de manera sostenida a la optimización racional del índice~$P$:
un estudiante que siguiera la estrategia conservadora de manera persistente
titularía~$4.8$~trimestres más tarde que uno que siguiera el baseline, y con
una probabilidad de egreso~$29$~puntos porcentuales inferior.}
La causa primaria de~$c_s\approx30$~cr/trim es la restricción topológica:
las cadenas de seriación bloquean la inscripción de créditos que el estudiante,
de otro modo, podría cursar.
El índice~$P$ amplifica la penalización de esas cadenas pero no la origina.

%──────────────────────────────────────────────────────
\subsection{El índice alternativo~\texorpdfstring{$h$}{h}: diseño de base entrópica}
\label{sec:indice_h}

El mecanismo vigente también incorpora las calificaciones en los
componentes~$\mathrm{PN}$ y~$\mathrm{PTN}$ (peso combinado~$c+e=0.2$).
Como muestran los datos del cuerpo docente CBI (\S\ref{sec:perfil_docente}),
la eficacia por docente presenta un coeficiente de variación mediano de~$0.21$,
con el~34\% de las UEAs con tres o más docentes superando~$\mathrm{CV}>0.30$;
en las cinco UEAs de mayor riesgo el rango es~$\mathrm{CV}\in[0.37,\,0.47]$.
Esto implica que la componente de calificaciones de~$P$ acumula ruido docente:
estudiantes con historial de acceso a docentes de alta eficacia obtienen~$\mathrm{PN}$
más alto, lo que mejora~$P$ y les garantiza acceso preferente en trimestres
subsiguientes, amplificando la inequidad generada por la heterogeneidad del pool
docente.

La propuesta de índice alternativo~$h$ \cite{prelacion_cslm} parte del
principio de máxima neutralidad informativa: dada la información disponible
sobre el ritmo de avance, la ponderación entre el estimador instantáneo y el
histórico debe maximizar la entropía de Shannon sobre el espacio de pesos.
Se definen la eficiencia trimestral
\begin{equation}
  \eta_t = \frac{C_{A,t-1}}{C_{\max}},
  \label{eq:eta}
\end{equation}
donde $C_{\max}=45$~cr/trim es el máximo inscribible, y su promedio histórico
\begin{equation}
  \langle\eta\rangle_t = \frac{1}{t}\sum_{n=2}^{t}\frac{C_{A,n-1}}{C_{\max}};
  \label{eq:eta_bar}
\end{equation}
ambos son estimadores insesgados de una misma cantidad~$\theta$ (el avance
esperado para el siguiente trimestre), de modo que el estimador lineal
$\theta_w = w\,\eta + (1-w)\langle\eta\rangle$ con~$w$ distribuido
uniformemente sobre~$[0,1]$ por el principio de máxima entropía produce
$w=1/2$, es decir
\begin{equation}
  h = \frac{\eta_t + \langle\eta\rangle_t}{2}.
  \label{eq:indice_h}
\end{equation}
El índice~$h$ prescinde de calificaciones y de grado de avance: solo observa
créditos aprobados sobre créditos máximos inscribibles, eliminando así la
componente de ruido docente presente en~$P$.

La Figura~\ref{fig:h_vs_P} compara las trayectorias de~$h$ y~$P$ para cuatro
planes representativos y muestra la dispersión de ambos índices en el
trimestre~12.
En general, $P>h$ para los planes con mayor fracción de calificaciones
satisfactorias, mientras que~$h$ y~$P$ convergen hacia el período de egreso
conforme el grado de avance se aproxima a~1.

\begin{figure}[H]\centering
\includegraphics[width=\linewidth]{../../Output/figuras_correlaciones/fig_32d_h_vs_P.pdf}
\caption{Comparación del índice de desempeño vigente~$P$ (Acuerdo~551.4.4.1-2015)
  con el índice alternativo~$h=(\eta_t+\langle\eta\rangle)/2$,
  donde~$\eta_t=C_{A,t-1}/45$ es la fracción del máximo de créditos inscribibles
  aprobada el trimestre anterior y~$\langle\eta\rangle$ su promedio acumulado.
  \emph{Panel superior:} trayectorias medias de~$P$ (línea continua) y~$h$
  (línea discontinua) bajo el baseline para cuatro planes representativos.
  En todos los planes~$P>h$ en etapas tempranas porque~$P$ incorpora las
  calificaciones ($\mathrm{PN}$, $\mathrm{PTN}$) y el Grado de Avance
  $\mathrm{GA}=C_A/C_P$ (fracción de créditos del plan acreditados); ambos
  índices convergen hacia el egreso conforme~$\mathrm{GA}\to1$.
  \emph{Panel inferior:} dispersión de~$P$ frente a~$h$ en el trimestre~12
  para los diez planes; la diagonal representa acuerdo perfecto.
  Los planes por encima de la diagonal tienen~$P>h$: el componente de
  calificaciones eleva~$P$ respecto al indicador de puro ritmo crediticio~$h$.}
\label{fig:h_vs_P}
\end{figure}

%──────────────────────────────────────────────────────
\subsection{Estrategias de inscripción bajo~\texorpdfstring{$h$}{h}: extinción del incentivo al subregistro}
\label{sec:h_estrategias}

La diferencia estructural entre~$P$ y~$h$ tiene consecuencias directas sobre
los incentivos de inscripción que cada índice genera.
Como~$\eta_t = C_{A,t-1}/C_\mathrm{max}$ mide créditos aprobados en valor
absoluto ---no una razón inscrito/aprobado--- la única forma de incrementar~$\eta$
es aprobar más créditos, lo que requiere inscribir más UEAs de alta probabilidad
de aprobación~$a_i$.
La estrategia que maximiza~$h$ por trimestre es, por tanto, la opuesta a la que
maximiza el componente~$\mathrm{ACT}$ de~$P$: inscribir la carga máxima posible
seleccionando UEAs en orden descendente de~$a_i$, estrategia denominada aquí
\emph{$h$-óptima}~($c^\ast=38$~cr/trim, selección por~$a_i$ descendente)
(nótese que $c^\ast=38$~cr/trim es el tope de inscripción de la
estrategia agresiva ---el ritmo de referencia del índice~$P$ vigente---,
distinto del denominador fijo $C_{\max}=45$~cr/trim que
aparece en la definición de $\eta_t$ en~\eqref{eq:eta}).

Para cuantificar esta diferencia se simulan cuatro estrategias de inscripción
sobre la misma cohorte~($S=3\,000$ estudiantes/plan, mismo vector de fragilidad
latente~$\theta_s$): la agresiva~($c^\ast=38$~cr/trim, orden topológico), el
baseline~($c_s$ calibrado, orden topológico), la conservadora~($24$~cr/trim,
fácil-primero) y la~$h$-óptima~($38$~cr/trim, fácil-primero).

La Figura~\ref{fig:h_trayectorias} confirma que el orden de las trayectorias
de~$h$ es monótono en todo~$t$:
$h_\mathrm{agresiva} > h_{h\text{-óptima}} \geq h_\mathrm{baseline}
> h_\mathrm{conservadora}$.
El estudiante racional que maximizara~$h$ no encontraría ventaja alguna en
reducir su carga de inscripción; al contrario, la reducción penaliza~$h$
inmediata y permanentemente.

\begin{figure}[H]\centering
\includegraphics[width=0.96\linewidth]{../../Output/figuras_correlaciones/fig_35b_h_trayectorias.pdf}
\caption{Trayectorias del índice alternativo~$h=(\eta_t+\langle\eta\rangle)/2$
  (donde~$\eta_t=C_{A,t-1}/45$ es la fracción del máximo inscribible de créditos
  aprobada el trimestre anterior) con intervalo intercuartílico (IQR, rango P25--P75,
  zona sombreada) para cuatro planes representativos bajo cuatro estrategias.
  \emph{Agresiva} (azul claro punteado): $c^\ast=38$~cr/trim, orden topológico.
  \emph{Baseline} (azul oscuro sólido): carga histórica~$c_s$, orden topológico.
  \emph{Conservadora} (rojo discontinuo): 24~cr/trim, UEAs de mayor~$a_i$ primero.
  \emph{$h$-óptima} (verde): $c^\ast=38$~cr/trim, UEAs de mayor~$a_i$ primero.
  El orden entre estrategias es monótono y se mantiene en todo el horizonte:
  conservadora~$<$ baseline~$<$ $h$-óptima~$<$ agresiva en~$h$ para todo~$t$.
  A diferencia del índice~$P$, no existe ninguna ventana temporal en que inscribir
  menos eleve~$h$.}
\label{fig:h_trayectorias}
\end{figure}

\medskip\noindent\textit{Coherencia de~$h$ con el egreso.}
La Figura~\ref{fig:h_coherencia} muestra~$\mathrm{E}[T\mid\text{tit}]$ en función
del valor medio de~$h$ en~$T=8$.
La correlación es marcadamente más coherente que la observada para~$P$ en la
Figura~\ref{fig:trampa_scatter}: mayor~$h$ se asocia con un egreso más rápido
para agresiva, baseline y conservadora.
La única excepción es la estrategia~$h$-óptima: con~$c^\ast=38$~cr/trim igual
a la agresiva pero con selección fácil-primero en vez de topológica, egresa
$\mathrm{E}[T\mid\text{tit}]_{h\text{-opt}}=19.3$~trimestres (excl. Metalurgia),
$2.1$~trimestres más tarde que la agresiva~($17.2$~trim) aunque con
idéntica carga máxima.
La flecha que conecta ambas estrategias en la figura cuantifica el costo de
ignorar el orden topológico: liberar prerrequisitos antes de acumular créditos
fáciles recorta $2.1$~trimestres del tiempo al egreso.

\begin{figure}[H]\centering
\includegraphics[width=0.80\linewidth]{../../Output/figuras_correlaciones/fig_35c_h_coherencia.pdf}
\caption{Coherencia del índice~$h=(\eta_t+\langle\eta\rangle)/2$ con la velocidad
  de egreso.
  El eje horizontal es el valor medio de~$h$ en el trimestre~8; el eje vertical
  es el tiempo promedio de egreso condicionado al egreso
  $\mathrm{E}[T\mid\text{tit}]$ (en trimestres; menor es mejor).
  Las cuatro estrategias son: \emph{agresiva} (triángulos azul claro:
  $c^\ast=38$~cr/trim, orden topológico), \emph{baseline} (círculos azul oscuro:
  carga histórica~$c_s$, orden topológico), \emph{conservadora} (cuadrados rojos:
  24~cr/trim, UEAs de mayor~$a_i$ primero) y \emph{$h$-óptima} (rombos verdes:
  $c^\ast=38$~cr/trim, UEAs de mayor~$a_i$ primero).
  Las flechas conectan la estrategia agresiva con la~$h$-óptima del mismo plan:
  idéntica carga máxima~$c^\ast=38$~cr/trim, pero distinto criterio de selección
  (topológico vs.\ fácil-primero).
  El descenso en~$h$ y el incremento en~$\mathrm{E}[T\mid\text{tit}]$ a lo largo
  de la flecha cuantifican el costo de ignorar el orden topológico
  ($2.1$~trimestres de retraso promedio, excl.\ Metalurgia).
  Metalurgia se excluye porque $c^\ast=38 < c_s=41$~cr/trim impide que los
  estudiantes egresen dentro del horizonte de simulación en dos estrategias.}
\label{fig:h_coherencia}
\end{figure}

\medskip\noindent\textit{Penalización topológica bajo~$h$.}
Aunque el índice~$h$ elimina el incentivo al subregistro, no elimina la
penalización topológica.
Un estudiante bloqueado por seriaciones aprueba menos créditos por trimestre,
lo que reduce~$\eta_t$ de manera análoga a como la seriación deprimía~$\mathrm{GA}$
en~$P$.
La restricción topológica impone un techo sobre los créditos aprobables por
trimestre que ningún índice basado en flujo crediticio puede suprimir; la
penalización queda cuantificada en la sección~\ref{sec:comparacion_indices}.

%──────────────────────────────────────────────────────
\subsection{Comparación~\texorpdfstring{$P$}{P} versus~\texorpdfstring{$h$}{h}: incentivos, egreso y restricción topológica}
\label{sec:comparacion_indices}

Las secciones precedentes examinaron~$P$ y~$h$ por separado; aquí se contrastan
directamente sobre las mismas cohortes y estrategias.
Tres diferencias estructurales emergen con claridad.

\textit{Incentivos de inscripción.}
La Figura~\ref{fig:incentivos_P_vs_h} resume el contraste central: la diferencia
$\Delta\,\text{índice}(t) = \text{índice(conservadora)}-\text{índice(baseline)}$
es positiva bajo~$P$ durante~$T\in[1,2.5]$ ($\Delta P(2)=+0.111$) y estrictamente
negativa bajo~$h$ en todo~$t\geq1$ ($\Delta h(2)=-0.043$, $\Delta h(12)=-0.171$).
El componente~$\mathrm{ACT}$ de~$P$ premia la razón créditos aprobados sobre créditos
inscritos: reducir el denominador ---inscribir menos--- eleva la razón incluso cuando
el numerador cae.
El índice~$h=(\eta+\langle\eta\rangle)/2$ mide créditos aprobados sobre un denominador
fijo ($C_\mathrm{max}=45$~cr/trim): la única manera de subir~$h$ es aprobar más créditos
absolutos, alineando el incentivo racional con la velocidad de avance en el plan.

\begin{figure}[H]\centering
\includegraphics[width=0.86\linewidth]{../../Output/figuras_correlaciones/fig_35a_incentivos_P_vs_h.pdf}
\caption{Contraste de incentivos de inscripción entre el índice vigente~$P$
  (Acuerdo~551.4.4.1-2015) y el índice alternativo~$h=(\eta_t+\langle\eta\rangle)/2$.
  Las curvas muestran la diferencia
  $\Delta\,\text{índice}(t)=\text{índice(conservadora)}-\text{índice(baseline)}$,
  promediada sobre los diez planes, entre la estrategia \emph{conservadora}
  (24~cr/trim, UEAs de mayor tasa de aprobación~$a_i$ primero) y el \emph{baseline}
  (carga histórica~$c_s$, orden topológico).
  Línea discontinua roja:~$\Delta P$; línea sólida verde:~$\Delta h$.
  La zona sombreada delimita la ventana~$T\in[1,2.5]$ en que~$\Delta P>0$
  ($\Delta P(2)=+0.111$): el componente~$\mathrm{ACT}=C_{A}/C_{I}$ de~$P$
  premia inscribir menos porque reduce el denominador~$C_I$ sin reducir
  proporcionalmente~$C_A$.
  La curva~$\Delta h$ es estrictamente negativa para todo~$t\geq1$
  ($\Delta h(2)=-0.043$): como~$h$ mide créditos aprobados sobre el denominador
  fijo~$C_{\max}=45$~cr/trim, inscribir menos solo puede reducir~$h$.}
\label{fig:incentivos_P_vs_h}
\end{figure}

\textit{Coherencia con el egreso.}
Bajo~$P$, la estrategia conservadora puede acumular un índice comparable al del
baseline en~$T=8$ con tiempos al egreso entre~$4$ y~$5$~trimestres mayores
(Figura~\ref{fig:trampa_scatter}): el índice puede ser inflado sin que la velocidad
de avance lo justifique.
Bajo~$h$, el ordenamiento por estrategia es coherente con el egreso en todo
momento: conservadora~$<$ baseline~$<$ $h$-óptima~$<$ agresiva tanto en~$h(T=8)$
como en velocidad de egreso (Figura~\ref{fig:h_coherencia}), salvo la excepción
de que la estrategia~$h$-óptima no respeta el orden topológico y egresa~$2.1$
trimestres más tarde que la agresiva con idéntica carga.

\textit{Penalización topológica compartida.}
Ambos índices son sensibles al bloqueo topológico porque ambos dependen del flujo
crediticio: en~$P$ el bloqueo deprime~$\mathrm{GA}=C_A/C_P$ y, en etapas avanzadas
con~$b=0.35$, arrastra~$P$ hacia abajo; en~$h$ el bloqueo reduce~$\eta_t=C_{A,t-1}/45$
y el promedio acumulado~$\langle\eta\rangle$, deprimiendo~$h$ de manera análoga.
La Figura~\ref{fig:h_componentes_elo} muestra esta dinámica para el plan Electrónica:
la separación entre estrategias en~$\eta_t$ replica, con distinta escala, la
separación en~$\mathrm{GA}$ observada bajo~$P$.
El mecanismo de penalización topológica es, por tanto, independiente de la fórmula
del índice: opera a través del flujo crediticio bruto, y ningún índice basado en
ese flujo puede ignorarlo mientras las seriaciones permanezcan en el plan.

\begin{figure}[H]\centering
\includegraphics[width=0.96\linewidth]{../../Output/figuras_correlaciones/fig_35d_h_componentes_elo.pdf}
\caption{Descomposición del índice alternativo~$h=(\eta_t+\langle\eta\rangle)/2$
  en el plan Electrónica bajo cuatro estrategias de inscripción:
  \emph{agresiva} (azul claro: $c^\ast=38$~cr/trim, orden topológico),
  \emph{baseline} (azul oscuro: carga~$c_s$, orden topológico),
  \emph{conservadora} (rojo: 24~cr/trim, UEAs de mayor~$a_i$ primero) y
  \emph{$h$-óptima} (verde: $c^\ast=38$~cr/trim, UEAs de mayor~$a_i$ primero).
  \emph{Panel superior, izquierda a derecha:}
  (1)~eficiencia trimestral~$\eta_t=C_{A,t-1}/45$, fracción de los 45~cr/trim
  máximos inscribibles aprobada en el trimestre anterior;
  (2)~promedio acumulado~$\langle\eta\rangle=t^{-1}\sum_{n=2}^{t}\eta_n$;
  (3)~índice~$h$ resultante.
  La separación entre estrategias en~$\eta_t$ desde~$T=2$ es el análogo bajo~$h$
  del diferencial en el Grado de Avance~$\mathrm{GA}=C_A/C_P$ (fracción de
  créditos del plan acreditados) que genera la doble penalización bajo~$P$:
  las seriaciones imponen un techo sobre~$C_A$ que ningún índice crediticio puede
  suprimir.
  \emph{Panel inferior:} comparación directa de~$h$ y del índice vigente~$P$
  (Acuerdo~551.4.4.1-2015) en el baseline; $P>h$ sistemáticamente porque~$P$
  incorpora las calificaciones (componentes~$\mathrm{PN}$ = Promedio Normalizado,
  $\mathrm{PTN}$ = Promedio Total Normalizado) y la ponderación variable del
  Grado de Avance.}
\label{fig:h_componentes_elo}
\end{figure}

\textbf{En síntesis:} $P$ crea un incentivo mensurable al subregistro estratégico
($\Delta P(2)=+0.111$) que desaparece antes del cuarto trimestre y resulta
contraproducente a largo plazo~($4.8$~trimestres adicionales, $-29$~pp de
tasa de egreso para la estrategia conservadora persistente).
El índice~$h$ elimina ese incentivo~($\Delta h<0$ en todo~$t$) y alinea el
comportamiento racional ---inscribir más, seleccionar UEAs de alta~$a_i$--- con
la velocidad de egreso.
Sin embargo, ninguno de los dos índices puede disolver la penalización topológica:
estudiantes bloqueados por seriaciones acumulan menos créditos por trimestre y
quedan rezagados en~$P$ o en~$h$ precisamente cuando el acceso preferente a los
grupos del cuello de botella sería más valioso.
La reforma del criterio de prelación es condición necesaria pero no suficiente;
la remoción de las seriaciones de mayor impacto sigue siendo la intervención que
ataca el mecanismo de fondo.

%══════════════════════════════════════════════════════
\section{Análisis global de complejidad sistémica}
\label{sec:complejidad}
%══════════════════════════════════════════════════════

Los resultados reportados en las secciones previas describen dimensiones parciales de un problema multi-causa.
El modelo de cascada cuantifica la probabilidad de completar la cadena de seriaciones
(\S\ref{sec:modelo}); el análisis de remoción estima el costo topológico de cada arista
(\S\ref{sec:efecto}); el análisis de índices caracteriza el incentivo al subregistro del criterio de
prelación (\S\ref{sec:comparacion_indices}); el perfil docente revela la heterogeneidad del insumo
educativo (\S\ref{sec:perfil_docente}); y las curvas de recuperabilidad mapean la resiliencia del
sistema ante el rezago acumulado (\S\ref{sec:correlaciones}).
Estas cinco perspectivas son complementarias, no redundantes.

Para integrarlas formalmente se construye un índice de complejidad sistémica~$C\in[0,1]$ como media
aritmética de cinco dimensiones normalizadas,
\begin{equation}
  C = \frac{D_1 + D_2 + D_3 + D_4 + D_5}{5},
  \label{eq:C_index}
\end{equation}
donde cada $D_k\in[0,1]$ se obtiene por normalización min-max de las métricas crudas a través de los diez
planes, de modo que el valor más complejo en cada dimensión recibe~1 y el menos complejo recibe~0.
Las cinco dimensiones son: $D_1$, complejidad topológica (profundidad máxima del grafo, densidad de
aristas, coeficiente de Gini de la centralidad de intermediación, fracción de nodos en el camino crítico,
entropía de la distribución de longitudes de caminos); $D_2$, complejidad probabilística (entropía de
Shannon de la función de masa de probabilidad del tiempo de egreso~$T|tit$, coeficiente de variación
$\sigma_T/E[T|tit]$, rango intercuartílico normalizado, tasa de baja definitiva, fracción de UEAs con
tasa de aprobación $a_i<0.60$); $D_3$, complejidad evaluativa (depresión crónica del índice~$P$ para la
estrategia agresiva al trimestre~12, magnitud de la trampa~ACT ---diferencia máxima $\Delta P_{\rm
cons}-\Delta P_{\rm base}$ en~$t\le 3$---, brecha $P$~vs.~$h$ al trimestre~12); $D_4$, complejidad
docente (coeficiente de variación inter-profesor de la eficiencia media por UEA, fracción de UEAs con
$CV>0.30$, rango medio de eficiencias, efecto de la varianza docente sobre~$\Delta E[T|tit]$); y $D_5$,
complejidad de recuperación, definida como la dificultad para recobrar trayectoria tras el primer
suspenso, e inversamente proporcional al spread de probabilidad de aprobación acumulada entre el primer y
el quinto intento.
El cómputo se realiza sobre los archivos de datos que alimentan los análisis seccionales (script~37,
\texttt{Analisis/37\_complejidad\_global.py}); las salidas quedan en \texttt{Output/complejidad\_*.csv}.

%------------------------------------------------------
\subsection{Ranking y arquetipos de complejidad}
\label{subsec:ranking_C}
%------------------------------------------------------

\begin{figure}[ht]
  \centering
  \includegraphics[width=\linewidth]{../../Output/figuras_correlaciones/fig_37b_heatmap_complejidad.pdf}
  \caption{Mapa de complejidad sistémica: scores normalizados $D_k\in[0,1]$ de las cinco dimensiones
    para cada uno de los diez planes de estudio CBI 2020, ordenados de mayor a menor índice~$C$
    (definido en la ecuación~\eqref{eq:C_index}).
    Cada celda muestra el score en la escala viridis (amarillo = complejidad máxima, azul = mínima).
    $D_1$: topológica; $D_2$: probabilística; $D_3$: evaluativa; $D_4$: docente; $D_5$: recuperación.
    La columna $C$ es la media de las cinco dimensiones.
    Fuente: script~37, datos UAM-A 16I--25O.}
  \label{fig:heatmap_complejidad}
\end{figure}

\begin{figure}[ht]
  \centering
  \includegraphics[width=\linewidth]{../../Output/figuras_correlaciones/fig_37a_radar_complejidad.pdf}
  \caption{Huella de complejidad sistémica por plan: gráfico de radar de las cinco dimensiones
    normalizadas $D_1$--$D_5$ (definidas en el texto y en el eje de la figura) para cada plan de estudio.
    El área del polígono es proporcional a la complejidad total; la forma revela el arquetipo
    (concentrada vs.\ difusa, dominancia docente vs.\ topológica).
    Escala radial: 0 (centro) a 1 (perímetro), con anillos de referencia en 0.25, 0.50 y 0.75.
    Fuente: script~37 a partir de grafos de seriación y diccionarios de UEA y profesor, 16I--25O.}
  \label{fig:radar_complejidad}
\end{figure}

\begin{figure}[ht]
  \centering
  \includegraphics[width=\linewidth]{../../Output/figuras_correlaciones/fig_37e_descomposicion_C.pdf}
  \caption{Descomposición del índice de complejidad~$C$ (definido en la ecuación~\eqref{eq:C_index})
    por plan y dimensión.
    Cada barra apilada muestra la contribución $D_k/5$ de cada dimensión al índice total.
    El rombo azul indica el valor de~$C$ y la etiqueta su magnitud.
    Las bandas horizontales delimitan las clases de complejidad:
    Baja ($C<0.40$, verde), Moderada ($0.40\le C<0.55$, azul), Alta ($0.55\le C<0.70$, naranja),
    Crítica ($C\ge 0.70$, rojo).
    $D_1$: topológica (azul oscuro); $D_2$: probabilística (violeta); $D_3$: evaluativa (naranja);
    $D_4$: docente (verde oscuro); $D_5$: recuperación (rojo).
    Planes ordenados de mayor (izquierda) a menor~$C$.
    Fuente: script~37, datos UAM-A 16I--25O.}
  \label{fig:descomposicion_C}
\end{figure}

Las figuras~\ref{fig:heatmap_complejidad}--\ref{fig:descomposicion_C} muestran que tres planes alcanzan
la clase \emph{Alta} ($C>0.55$): Computación ($C=0.617$), Química ($C=0.602$) y Mecánica ($C=0.550$).
Los seis planes restantes se clasifican como \emph{Moderada} ($C\in[0.46,0.49]$): Ambiental, Industrial,
Física, Electrónica, Civil y Eléctrica.
Metalurgia es el único plan con complejidad \emph{Baja} ($C=0.379$).

En los tres planes de mayor índice la dimensión dominante no es la topológica ---que corresponde a
Electrónica por su cadena EDO$\to$Circuitos~Eléctricos~I con profundidad de grafo~9--- sino la
docente~$D_4$: Computación y Química concentran la mayor heterogeneidad inter-profesor de la división y
el mayor efecto de varianza docente sobre~$\Delta E[T|tit]$; Mecánica supera a ambas en la dimensión
evaluativa~$D_3$, con la mayor depresión crónica del índice~$P$ y la trampa~ACT más pronunciada.
Electrónica, a pesar de ser el plan topológicamente más complejo (139 seriaciones, profundidad máxima~9)
y el de peor~$E[T|tit]=23.5$ trimestres, ocupa el séptimo lugar ($C=0.477$) porque su carga está
concentrada en una sola dimensión.
La complejidad multi-dimensional difusa es distinta de la carga concentrada en un único cuello de botella.

A partir del perfil de las cinco dimensiones se identifican cuatro arquetipos de complejidad.
El primero, \emph{trampa docente} (Computación, Química, Electrónica)%
\footnote{Electrónica se clasifica en este arquetipo por su
$D_4=0.66$ (la más alta entre los tres planes). Su dimensión
topológica $D_1=0.60$ también es elevada, pero el cuello de
botella topológico es altamente concentrado (una sola cadena de
tres aristas: Cálculo Integral~$\to$ EDO~$\to$ Circuitos Eléctricos~I),
mientras que la heterogeneidad docente afecta a
la totalidad del Tronco General.},
combina alta heterogeneidad
docente con fuerte dependencia de la asignación de profesores; la intervención prioritaria es la
homogeneización del insumo educativo mediante programas de acompañamiento y formación docente
coordinados con los departamentos, y una asignación de grupos que destine a los nodos críticos a
los profesores con mayor eficacia histórica verificable.
El segundo, \emph{trampa evaluativa} (Mecánica, Civil), exhibe la mayor distorsión del índice~$P$ y el
incentivo más intenso hacia la inscripción conservadora; la sustitución de~$P$ por~$h$ resuelve el
incentivo sin modificar el currículo.
El tercero, \emph{plan frágil} (Eléctrica, Física), está dominado por la dimensión probabilística:
la distribución del tiempo de egreso tiene alta varianza y la fracción de UEAs con $a_i<0.60$ es la
mayor de la división; la remoción de las aristas de mayor impacto marginal en~$\Delta E[T|tit]$ es la
intervención de fondo.
El cuarto, \emph{plan asistido} (Ambiental, Industrial), tiene la mayor capacidad de recuperación (baja
$D_5$), lo que modera el efecto de los otros mecanismos y coloca a estos planes en una posición de
intervención menos urgente.
Metalurgia constituye un caso especial: su baja complejidad sistémica enmascara la restricción
extraacadémica de la estancia industrial obligatoria, que el presente modelo no captura y que opera como
cuello de botella externo al grafo de seriaciones.

%------------------------------------------------------
\subsection{Acoplamiento entre dimensiones}
\label{subsec:acoplamiento}
%------------------------------------------------------

\begin{figure}[ht]
  \centering
  \includegraphics[width=0.85\linewidth]{../../Output/figuras_correlaciones/fig_37d_acoplamiento.pdf}
  \caption{Matriz de acoplamiento entre las cinco dimensiones de complejidad sistémica.
    Cada celda muestra la correlación de Spearman~$r_s$ calculada a través de los $N=10$ planes de
    estudio CBI 2020.
    Celdas rojas indican acoplamiento positivo (la complejidad en ambas dimensiones coexiste);
    celdas azules indican acoplamiento negativo (una dimensión alta tiende a asociarse con la otra baja).
    El asterisco~(*) señala valores con $|r_s|\ge 0.63$, umbral aproximado de significancia~$p<0.05$
    para $N=10$ sin corrección por comparaciones múltiples.
    $D_1$: topológica; $D_2$: probabilística; $D_3$: evaluativa; $D_4$: docente; $D_5$: recuperación.
    Fuente: script~37, datos UAM-A 16I--25O.}
  \label{fig:acoplamiento_dims}
\end{figure}

La figura~\ref{fig:acoplamiento_dims} revela que el acoplamiento inter-dimensional más notable es el
negativo entre~$D_5$ (Recuperación) y~$D_2$ (Probabilística; $r_s=-0.673$) y entre~$D_5$
y~$D_4$ (Docente; $r_s=-0.576$).
Este patrón indica que los planes con alta complejidad probabilística y docente tienen
simultáneamente menor capacidad de recuperación tras el rezago: los estudiantes enfrentan tasas de
aprobación bajas y alta heterogeneidad docente, y cuando acumulan suspensos la probabilidad de reaprobación
no mejora sustancialmente con intentos adicionales.
La consecuencia es una trampa de complejidad auto-reforzada en la que el rezago inicial reduce el acceso
a los grupos del cuello de botella (mecanismo topológico, $D_1$), el índice~$P$ penaliza a quien
inscribió menos (mecanismo evaluativo, $D_3$), y la capacidad de recuperación no basta para cerrar la
brecha (mecanismo probabilístico, $D_2$).
Esta dinámica emergente no es observable desde ninguna de las cinco dimensiones aisladas.

%------------------------------------------------------
\subsection{Validez predictiva del índice}
\label{subsec:validez_C}
%------------------------------------------------------

\begin{figure}[ht]
  \centering
  \includegraphics[width=0.92\linewidth]{../../Output/figuras_correlaciones/fig_37c_C_vs_desenlaces.pdf}
  \caption{Índice de complejidad sistémica~$C$ (definido en la ecuación~\eqref{eq:C_index}) frente a los
    desenlaces institucionales: panel superior, tiempo promedio de egreso $E[T|tit]$ en trimestres;
    panel inferior, tasa de egreso (\%).
    Cada punto representa un plan de estudio; el tamaño del punto es proporcional a la tasa de baja
    definitiva del plan; el color indica la dimensión dominante ($D_k$ con mayor score normalizado).
    $r_s$: correlación de Spearman ($N=10$, Metalurgia incluido); ninguna de las correlaciones alcanza significancia a
    $\alpha=0.05$.
    La anotación sobre Electrónica señala el caso de alta carga topológica concentrada en un único
    cuello de botella, con $C$ moderado pero $E[T|tit]$ máximo de la división.
    Fuente: script~37, datos UAM-A 16I--25O; $E[T|tit]$ oficial según tabla~\ref{tab:et_oficial}.}
  \label{fig:C_vs_desenlaces}
\end{figure}

La correlación de rango entre~$C$ y~$E[T|tit]$ es $r_s=-0.139$ y entre~$C$ y la tasa de egreso
$r_s=-0.164$ (figura~\ref{fig:C_vs_desenlaces}); ambas no significativas con $N=10$.
Esto no invalida el índice: refleja que la complejidad multi-dimensional difusa no equivale a peor
desenlace promedio.
Electrónica tiene el mayor~$E[T|tit]=23.5$ trimestres con $C=0.477$ (séptimo lugar) porque toda su
\emph{penalización sobre $E[T|\text{tit}]$} está
concentrada en una sola cadena topológica (EDO~$\to$ Circuitos
Eléctricos~I), aunque su perfil de complejidad sistémica también
incorpora una dimensión docente relevante ($D_4=0.66$, segunda
más alta de la División), no en la superposición de múltiples disfunciones.
Computación tiene el mayor~$C=0.617$ con un~$E[T|tit]=20.4$ trimestres moderado porque sus problemas
se distribuyen en cuatro dimensiones activas, ninguna de las cuales domina individualmente el desenlace.
La utilidad del índice~$C$ reside en identificar qué dimensiones requieren atención simultánea y bajo qué
arquetipo cae cada plan; no en predecir cuál plan tiene la peor trayectoria en una variable escalar.
Un plan \emph{Alta}~$C$ con dimensión docente dominante requiere una intervención en el cuerpo de
profesores; el mismo nivel de~$C$ con dimensión evaluativa dominante requiere reforma del criterio de
prelación; y ambas intervenciones son independientes de si el plan tiene el mayor~$E[T|tit]$ de la
división.

%══════════════════════════════════════════════════════
\section{Limitaciones del modelo}
\label{sec:limitaciones}
%══════════════════════════════════════════════════════

\textbf{Oferta y cupo.}
El modelo asume que cada estudiante puede inscribir todas las UEAs elegibles
hasta su límite de carga, sin restricciones de cupo ni de oferta trimestral.
Que $\ETtit^{\,\mathrm{mod}}$ reproduzca el promedio histórico con error menor
a medio trimestre sugiere que estas restricciones están, en promedio, absorbidas
por el parámetro~$\ell$ calibrado.

\textbf{Metalurgia.}
El valor $\ETtit=16.25$~trim para Metalurgia es una cota inferior.
La diferencia con el oficial (22.26~trim) refleja la disponibilidad limitada
de la estancia industrial (\uea{Trabajo en Planta Metalúrgica}), no ofertada
todos los trimestres.

\textbf{Umbrales de crédito.}
La puerta temporal de los cursos con umbral de créditos se calcula como
$\lceil\mathrm{umbral}/(45\cdot\bar{a})\rceil$, asumiendo carga máxima.
En el modelo calibrado ($c_s\approx 30$~cr/trim), los estudiantes acumulan
créditos más lentamente; la puerta abre a $t\approx 12$ cuando la tasa real
implicaría $t\approx 17$.
El sesgo es \emph{optimista} (habilita los cursos umbral antes de tiempo),
lo que puede subestimar ligeramente el efecto de esa categoría de restricciones.
El bajo valor de~$\ET{12}$ no se ve afectado por esta aproximación.

\textbf{Heterogeneidad docente.}
La tasa de aprobación~$a_i$ empleada en el modelo es la media histórica de todos los
profesores que han impartido la UEA~$i$.
El análisis de heterogeneidad intra-UEA (\S\ref{sec:dimensiones_adicionales}) muestra
que el 34\% de las UEAs CBI con tres o más docentes presenta un coeficiente de
varianza de la eficacia superior a~0.30.
El perfil del cuerpo docente (\S\ref{sec:perfil_docente}) complementa este dato a nivel
de actor: el 31.3\% de los 738~docentes CBI opera con~$a_d < 0.65$, y la mediana
de Ciencias Básicas ($\tilde{a}_d=0.670$) está por debajo del umbral de riesgo.
Esto implica que~$a_i$ no es una propiedad fija de la UEA sino un promedio que oculta
varianza inter-docente relevante.
En los nodos de mayor riesgo, con \textsc{cv}$\in[0.37,\,0.47]$ (categoría
\emph{Docente + intrínseca}), una intervención pedagógica que redujera el \textsc{cv}
a la mediana ($\approx0.21$) aumentaría~$a_i$ de forma heterogénea según el docente
asignado en cada trimestre.
El modelo actual no captura esta varianza y, por tanto, los beneficios de la
homogeneización docente no están reflejados en las proyecciones de~$\ETtit$.
El módulo~21 (\S\ref{sec:varianza_docente}) cuantifica este efecto mediante un
modelo extendido que sortea docente por UEA$\times$trimestre.

\subsection*{Alcance jurídico}
\label{sec:alcance_juridico}

El presente reporte es un diagnóstico técnico interno elaborado con fines de
mejora curricular y planificación institucional.
El análisis jurídico detallado de cada hallazgo y su respaldo en el Reglamento
de Estudios Superiores de la UAM (RES~UAM~\cite{RESUAM}) y en la Ley General de
Educación Superior (LGES~\cite{LGES2021}) se desarrolla en el documento
complementario \texttt{07\_analisis\_juridico.pdf}; en caso de divergencia de
interpretación, prevalece ese documento.
Los hallazgos del presente reporte son válidos dentro de las siguientes
limitaciones explícitas.

\textbf{Alcance estadístico-agregado, no individual.}
El modelo opera sobre distribuciones de cohortes.
El diagnóstico establece que la fracción de la cohorte que completa el plan en
12~trimestres es estadísticamente próxima a cero bajo las condiciones históricas
modeladas; el valor observado es~$1.52\%$
(IC~95\%~CP: $[1.1\%,\,2.0\%]$), consistente con el intervalo de predicción
del modelo.
Nada en este análisis implica que ningún estudiante individual no pueda egresar en
ese plazo.

\textbf{Condicionado a parámetros históricos.}
Los resultados están condicionados a la carga efectiva histórica~$c_s
\approx 30$~cr/trim calibrada sobre las cohortes~16I--25O.

\textbf{Sin efecto retroactivo.}
Este reporte no cuestiona la validez legal de ninguna resolución académica
adoptada con anterioridad a su fecha de publicación conforme al RES~UAM vigente.
Las bajas definitivas por acumulación de evaluaciones No Acreditadas (NA) se
rigen por los Arts.~16.V y~44.XI del RES~UAM, cuya legalidad no es objeto de
este análisis.

\textbf{Instrumento de mejora, no de sanción.}
El análisis docente (\S\ref{sec:perfil_docente}) reporta distribuciones estadísticas
agregadas por departamento; no constituye evaluación individual de ningún profesor
ni genera consecuencias laborales por sí mismo.

\textbf{Coordinación previa con Rectoría.}
El análisis del índice~$P$ (\S\ref{sec:comparacion_indices}) señala efectos no
deseados del Acuerdo~551.4.4.1-2015.
Dado que ese Acuerdo emana de la Rectoría, cualquier distribución amplia de este
reporte debe ir precedida de una comunicación formal al Rector que enmarque el
hallazgo como diagnóstico propositivo y no como impugnación del acuerdo vigente.
Presentar el análisis junto con las recomendaciones de acción de la
Sección~\ref{sec:discusion} ---que incluye la propuesta de índice~$h$ como
sustitución, no como crítica aislada--- es la condición mínima para que la
recepción institucional sea constructiva.

El módulo~26 (\texttt{26\_gestion\_riesgo\_legal.py}) identifica tres escenarios
de riesgo legal con probabilidad estimada, argumentos de defensa y acciones
preventivas recomendadas.
El primero es la impugnación estudiantil de baja definitiva (probabilidad media):
el argumento de defensa central es que el reporte es un diagnóstico prospectivo,
no una declaración de nulidad retroactiva de bajas ya ejecutadas; la acción
preventiva crítica es documentar en actas de Junta que el reporte fue recibido
como diagnóstico y que la División aprobó un plan de acción.
El segundo es el cuestionamiento por parte de acreditadores (probabilidad baja):
los marcos del Consejo de Acreditación de la Enseñanza de la Ingeniería (CACEI) y del Consejo para la Acreditación de la Educación Superior (COPAES) valoran el autodiagnóstico con mejora continua, y el
reporte presentado con las recomendaciones de la Sección~\ref{sec:discusion}
constituye evidencia positiva ante el próximo ciclo de acreditación.
El tercero es la intervención administrativa desde Rectoría (probabilidad
baja-media): la División debe informar proactivamente a Rectoría antes de cualquier
distribución amplia del reporte, coordinando su publicación con el anuncio
simultáneo del piloto de reforma en Electrónica y Química.



%══════════════════════════════════════════════════════
\section{Discusión general}
\label{sec:discusion}
%══════════════════════════════════════════════════════

El conjunto de resultados presentados en este reporte no describe una colección de defectos
curriculares independientes, sino una arquitectura de restricciones que se refuerzan mutuamente.
La comprensión de esa arquitectura ---y de la jerarquía que la organiza--- es el requisito
previo para diseñar intervenciones que no resulten en mejoras parciales que dejen intactos los
mecanismos de fondo.
El mecanismo primario es el desajuste de carga: el currículo 2020 fue diseñado para un
estudiante que inscribe~45~cr/trim, pero la población real lleva históricamente entre 27 y
33~cr/trim, y esa brecha es condición suficiente para que el egreso a doce trimestres sea
estadísticamente improbable en los diez programas, con independencia de cualquier otra
variable.
Todos los demás mecanismos operan sobre este fondo y lo agudizan; ninguno lo genera.
La implicación de política es directa: cualquier intervención que no aborde la capacidad
efectiva de inscripción ---sea incrementando los créditos que los estudiantes pueden y desean
llevar, sea reduciendo el mínimo requerido para graduación, sea acortando el tiempo que
las seriaciones mantienen bloqueados a los estudiantes--- opera aguas abajo del problema
principal.

Las seriaciones tienen, sin embargo, un papel amplificador que varía cualitativamente entre
planes.
En ocho de los diez programas la restricción dominante es la carga: incluso sin ninguna
seriación, los estudiantes tardarían entre 18 y 21~trimestres en acumular los créditos
necesarios.
En Electrónica, en cambio, la profundidad topológica de su grafo de seriación
(nueve niveles) constituye por sí misma una cota inferior de permanencia que ninguna
mejora en la carga puede eliminar sin reforma curricular.
Metalurgia representa el caso límite inverso: su alta carga efectiva histórica
($c_s=41$~cr/trim) reduce $T_{\min}^{\,\mathrm{carga}}=15$~trimestres por debajo del
bloqueo topológico (ocho niveles, $T_{\min}^{\,\mathrm{topo}}=9$), haciendo de él
el único plan donde la remoción total de seriaciones produce eficiencia terminal
en~12~trimestres ($\mathrm{ET}(12)=10.4\,\text{\%}$).
Esto define dos regímenes de intervención cualitativamente distintos.
Para los ocho planes dominados por carga, la remoción de seriaciones es la segunda
prioridad ---real, rentable, pero secundaria--- y debe ir acompañada de medidas que eleven
la inscripción efectiva hacia el máximo establecido.
Para Electrónica, la remoción es condición necesaria, y la evidencia señala un portfolio
quirúrgico mínimo: la cadena Cálculo Integral $\to$ EDO $\to$ Circuitos Eléctricos~I
y la arista Diseño de Sistemas Electrónicos $\to$ Control Digital concentran el 80\% de
la ganancia máxima con sólo ocho seriaciones (15\% del total), con un índice de
subaditividad $\mathcal{A}=0.52$ que indica que el beneficio real de eliminar las ocho
simultáneamente es la mitad de lo que sugeriría su suma de efectos marginales.
Para Química, cinco seriaciones (9\% del total) reproducen el mismo 80\% de ganancia.
La concentración del efecto en tan pocas aristas hace que el costo político de la
revisión curricular sea considerablemente menor que la magnitud del problema sugiere.

La heterogeneidad docente emerge del análisis como el factor transversal más subestimado.
Las cinco UEAs de mayor riesgo topológico ---Cálculo Diferencial, Cálculo Integral, EDO,
Cinemática y Dinámica de Partículas y Dinámica del Cuerpo Rígido--- presentan todas
coeficientes de variación inter-profesor entre 0.37 y 0.47 en eficacia de aprobación, lo que
las clasifica como \emph{Docente + intrínseca}: la baja tasa de aprobación no es una
propiedad fija del contenido sino una magnitud que puede moverse sustancialmente con una
asignación diferente de profesores.
El rango histórico de tasas de aprobación entre los docentes que han impartido esa UEA
va de 3\% a 100\% en los nodos más heterogéneos, lo que ilustra que la eficacia
observada del currículo depende marcadamente de decisiones de programación y asignación
que son ajenas al diseño del plan de estudios.
El análisis de complejidad sistémica confirma que tres de los cuatro planes más complejos
---Computación, Química y Electrónica--- tienen la dimensión docente~$D_4$ como dominante,
no la topológica.
Esto sugiere que una fracción significativa del impacto atribuido a las seriaciones en la
tasa de aprobación está mediado por la heterogeneidad del profesorado que las cubre; la
remoción de una arista sin intervención docente simultánea puede producir ganancias menores
de las que predice el modelo, que asume tasas de aprobación históricas invariantes.

El criterio de prelación vigente, índice~$P$ (Acuerdo~551.4.4.1-2015), agrega una
perturbación evaluativa sobre el mecanismo topológico que amplifica el daño sin generarlo.
El componente~$\mathrm{ACT}$ (peso~$d=0.40$, el mayor de la fórmula) crea un incentivo
real al subregistro durante los primeros dos a dos y medio trimestres
($\Delta P(T=2)=+0.111$ para la estrategia conservadora de~24~cr/trim frente al baseline),
pero el incentivo se autoextingue porque el grafo de seriaciones agota las UEAs de alta
aprobación accesibles sin prerrequisito y el déficit en Grado de Avance~GA invierte la
ventaja antes del cuarto trimestre.
El resultado final de seguir la estrategia conservadora de forma persistente es devastador:
4.8~trimestres adicionales de permanencia y 29~puntos porcentuales menos de egreso
respecto al baseline.
La evidencia no sugiere, por tanto, que la carga histórica de 27--41~cr/trim sea consecuencia
del comportamiento racional inducido por~$P$; la restricción topológica explica ese valor
sin necesidad de invocar racionalidad estratégica del alumnado.
La perturbación evaluativa opera sin embargo como un ruido adicional que penaliza
diferenciadamente a los estudiantes bloqueados por seriaciones: en Electrónica, la brecha de
índice~$P$ entre bloqueados y desbloqueados en~$T=12$ alcanza~$\Delta P=0.21$~puntos, lo
que les impide acceder a los grupos del cuello de botella que necesitarían para desbloquearse.
El índice alternativo~$h=(\eta_t+\langle\eta\rangle)/2$ elimina ese incentivo con
una modificación reglamentaria que no requiere reforma curricular, y produce un ordenamiento
entre estrategias de inscripción coherente con la velocidad de egreso en todos los
trimestres desde el primero.
Su implementación constituye la intervención de menor costo institucional y de impacto
mediato más claro de las identificadas en este análisis; es también la única que puede
comenzar con una modificación del Acuerdo institucional vigente sin esperar la aprobación
de cambios en el plan de estudios.

\subsection{Convergencia metodológica}
\label{sec:convergencia}

\begin{okbox}
\textbf{Los cuatro métodos de análisis producen de forma independiente el mismo diagnóstico:
la brecha entre el plazo del plan de estudios y el tiempo real de egreso
es estructural, medible y corregible.}
El modelo topológico-estocástico de cascada, el análisis de remoción combinada,
las correlaciones estadísticas multinivel y el índice de complejidad sistémica~$C$
no comparten supuestos, datos de entrada ni procedimientos de estimación.
La coherencia entre métodos independientes es evidencia de que el diagnóstico no
es artefacto de ningún supuesto particular, sino que refleja una realidad estructural del currículo.
\end{okbox}

El modelo de cascada establece las cotas de permanencia bajo parámetros históricos;
el análisis de remoción identifica las aristas específicas y su valor marginal; las
correlaciones estadísticas vinculan los rasgos estructurales del currículo con los
desenlaces a escala de División; y el índice~$C$ sintetiza los cinco mecanismos en un
mapa de complejidad que orienta la priorización de intervenciones por arquetipo.
Que cuatro metodologías independientes señalen los mismos planes, los mismos nodos y
las mismas intervenciones de mayor retorno constituye la validación cruzada más robusta
disponible dentro del alcance de los datos.

La priorización de las acciones recomendadas sigue la jerarquía de los mecanismos
identificados.
En primer término, la intervención quirúrgica en seriaciones de alto impacto en los planes
del cuadrante quirúrgico ---Electrónica y Química--- produce la mayor ganancia en~$\ETtit$
con el menor número de modificaciones al plan de estudios; el modelo predice reducciones de
2.2 y 1.5~trimestres, respectivamente, con portfolios de 8 y 5~aristas.
En segundo término, la reforma del criterio de prelación de~$P$ a~$h$ elimina el incentivo
al subregistro en todos los planes simultáneamente sin modificar ningún plan de estudios
y corrige la penalización diferencial de los estudiantes bloqueados; puede implementarse
antes de que concluya cualquier proceso de revisión curricular.
En tercer término, la intervención docente en los nodos clasificados como~\emph{Docente +
intrínseca} ---especialmente los cinco cubiertos por el Departamento de Ciencias Básicas
(\S\ref{sec:perfil_docente})--- tiene el potencial de reducir la tasa de aprobación de
primer intento sin ninguna modificación estructural al currículo.
En cuarto término, los planes clasificados como \emph{trampa docente} (Computación, Química,
Electrónica) requieren que las intervenciones anteriores se acompañen de protocolos de
homogeneización pedagógica, porque en ellos la remoción de seriaciones sin corrección
docente subestimará las ganancias predichas por el modelo.
En quinto término, los planes clasificados como \emph{plan frágil} (Eléctrica, Física),
donde la dimensión probabilística~$D_2$ es dominante, requieren reducir la fracción de
UEAs con $a_i<0.60$ como intervención primaria, lo que orienta el esfuerzo hacia los
nodos de alta frecuencia de aparición en los planes, no necesariamente los de mayor
profundidad topológica.
Los planes \emph{asistidos} (Ambiental, Industrial), con alta capacidad de recuperación
que modera parcialmente los otros mecanismos, pueden priorizarse en una segunda fase
sin pérdida sustancial de eficiencia del programa de reforma.
En todos los casos, la implementación debe ir acompañada de un sistema de seguimiento
basado en los mismos indicadores del análisis ---$\ETtit$, tasa de egreso, índice
$C$ por dimensión--- para detectar desviaciones entre los efectos predichos y los
observados y ajustar el portfolio de intervenciones antes del siguiente ciclo
de revisión curricular.




%══════════════════════════════════════════════════════
\section{Conclusiones}
\label{sec:conclusiones}
%══════════════════════════════════════════════════════

\textbf{Primero}: el currículo 2020 fue diseñado para un estudiante que
inscribe $\approx 45$~cr/trim; la población real inscribe históricamente
$\approx 27$--$33$~cr/trim.
Esa brecha es condición suficiente para que el egreso a 12~trimestres
sea estadísticamente improbable en los diez programas.
La evidencia histórica confirma el diagnóstico: $\ET{12}=1.52\%$
(IC~95\%: $[1.1\%,\,2.0\%]$) en las cohortes 2020--2021,
consistente con el intervalo de predicción del modelo bajo incertidumbre
paramétrica en~$c_s$.

\textbf{Segundo}: el mecanismo dominante en ocho de los diez programas es la
carga efectiva histórica, no la topología.
Con $c_s\approx 30$~cr/trim (la carga que los estudiantes realmente han
llevado), los créditos necesarios para egresar tardan entre 18 y 21~trimestres
en acumularse.
La profundidad topológica máxima es de 9~niveles (Computación y Electrónica),
imponiendo una cota de 10~trimestres que no alcanza por sí sola el umbral de~12.

\textbf{Tercero}: eliminar las seriaciones más costosas acortaría el egreso
entre 0.2~trim (Computación) y 2.3~trim (Electrónica).
La arista individual más impactante en la División es EDO $\to$ Circuitos
Eléctricos~I en Electrónica ($-1.07$~trim, $+2.1$~pp de egreso).
Sin embargo, incluso eliminar \emph{todas} las seriaciones no resuelve el
desajuste primario de carga en nueve planes: las seriaciones son la segunda
capa del problema, no la primera.

\textbf{Cuarto}: el análisis de correlaciones estadísticas (\S\ref{sec:correlaciones})
revela que el desempeño de un plan a escala de División depende principalmente de
dos rasgos del currículo: el número total de seriaciones
($r_s=-0.80$ con la tasa de egreso, $p=0.006$) y la fracción de UEAs con tasa
de aprobación inferior a~0.65
($r_s=+0.80$ con el trimestre de 50\% egresado, $p=0.005$).
Planes con muchas seriaciones y muchas materias difíciles tardan sistemáticamente
más y titúlan menos.
El plan Industrial ---pocas seriaciones, alta aprobación, baja profundidad--- es el
que más se aproxima al perfil óptimo observable en los datos actuales.

\textbf{Quinto}: a nivel de UEA, el riesgo topológico se concentra en una minoría
de nodos que combinan baja aprobación con alto alcance en el DAG.
Las cuatro UEAs de mayor riesgo ---Cálculo Diferencial, Cálculo Integral, Ecuaciones
Diferenciales Ordinarias y Cinemática y Dinámica de Partículas--- son obligatorias
en los diez planes con tasas de aprobación entre 0.51 y 0.57 y entre 308 y 488
UEAs descendientes.
El análisis de heterogeneidad docente (\S\ref{sec:dimensiones_adicionales}) revela
que las cinco UEAs del top-5 presentan \textsc{cv} de eficacia entre profesores en el
rango~$0.37$--$0.47$, lo que las clasifica como \emph{Docente + intrínseca}: la baja
aprobación no es solo propiedad del contenido, sino que está amplificada por la
dispersión pedagógica entre los docentes asignados.
La intervención más rentable combina por tanto \emph{reducción de seriaciones
salientes} con \emph{homogeneización de la práctica docente} en estos nodos.

\textbf{Sexto}: el impacto de las seriaciones sigue una distribución de Pareto
marcada: el 26.5\% de las aristas (132 de 499) concentra el 80\% del impacto
total sobre~$\ETtit$.
Esto hace posible concentrar la revisión curricular en un conjunto reducido de
seriaciones de alto impacto en lugar de revisar las más de 1\,200~relaciones del
currículo completo.

\textbf{Séptimo}: la remoción \emph{simultánea} de seriaciones es siempre
subaaditiva (§\ref{sec:dosis_respuesta}): el beneficio combinado de eliminar
$N$~aristas a la vez es una fracción~$\mathcal{A}\in[0.23,\,0.89]$ del resultado
que predice la suma de los efectos marginales individuales.
Sin embargo, la curva dosis-respuesta muestra saturación temprana en los planes más
críticos: en Electrónica se captura el 80\% de la ganancia máxima (2.18~trim) con
sólo 8~seriaciones (15\% del total); en Química, con apenas 5~seriaciones (9\%).
Las curvas de Lorenz (Figura~\ref{fig:remocion_resumen}D) cuantifican este efecto
de concentración para cada plan y permiten diseñar reformas quirúrgicas con tamaño
mínimo de intervención.

\textbf{Octavo}: el análisis del cuerpo docente CBI (738~profesores con actividad
en 16I--25O, §\ref{sec:perfil_docente}) añade un cuarto nivel de resolución al
diagnóstico.
La mediana divisional de aprobación docente es~$\tilde{a}_d=0.744$, pero el
departamento de Ciencias Básicas ---que cubre los nodos topológicamente más
críticos--- opera con la mediana de aprobación docente más baja de CBI y la mayor
dispersión (\S\ref{sec:perfil_docente}, tabla~\ref{tab:dept_docentes}).
El 31.3\% de los docentes CBI presenta~$a_d<0.65$, el mismo umbral de riesgo
alto que identifica las UEAs cuello de botella.
En las UEAs de alto riesgo, la eficacia mediana del docente \emph{sobre esa UEA
específica} cae entre~0.50 y~0.58, y los coeficientes de variación
intra-UEA (0.38--0.48) implican que el rango de aprobación esperada según el
docente asignado va de~3\% a~100\% en la misma asignatura.
El perfil docente confirma que el riesgo de los nodos críticos no es solo
intrínseco al contenido sino está amplificado por la heterogeneidad pedagógica
del profesorado que los cubre habitualmente.

\textbf{Noveno}: el déficit de Ritmo de Avance~$\mathrm{RA}=\min(C_A/(38\cdot T),1)<1$
es universal y estructural: más del~98\% de los estudiantes lo presentan desde
el tercer trimestre en nueve de los diez planes, independientemente de su
historial individual.
Las seriaciones generan una \emph{doble penalización} sobre el índice de
desempeño~$P$ (Acuerdo~551.4.4.1-2015):
deprimen el Grado de Avance~$\mathrm{GA}=C_A/C_P$ al bloquear créditos que el
estudiante podría inscribir, y simultáneamente reservan el acceso preferente
a los grupos del cuello de botella a quienes ya han acreditado la UEA bloqueante.
En los planes de perfil quirúrgico, la brecha entre estudiantes bloqueados y
desbloqueados alcanza~$\Delta P\approx0.21$~puntos en~$T=12$; al eliminar todas
las seriaciones se recuperan hasta~$0.27$~puntos de índice en Electrónica
y~$0.23$ en Química.

\textbf{Décimo}: el componente~$\mathrm{ACT}$ (razón de créditos aprobados sobre
créditos inscritos en los dos últimos trimestres, peso~$d=0.40$ en~$P$, el mayor
de todos) crea un incentivo real al subregistro estratégico durante~$T\in[1,2.5]$
($\Delta P(2)=+0.111$ para la estrategia conservadora de~24~cr/trim frente al
baseline de~$c_s$~cr/trim), pero el incentivo es estructuralmente autolimitante:
el grafo de seriaciones agota el inventario de UEAs de alta aprobación~$a_i$ no
bloqueadas en las primeras inscripciones y el déficit acumulado en~$\mathrm{GA}$
invierte la ventaja antes del cuarto trimestre.
\textbf{La carga histórica~$c_s\approx30$~cr/trim no puede atribuirse a la
optimización racional del índice~$P$}: un estudiante que siguiera la estrategia
conservadora de forma persistente titularía~$4.8$~trimestres más tarde con una
tasa de egreso~$29$~puntos porcentuales menor que el baseline.
La restricción topológica, no el comportamiento estratégico inducido por~$\mathrm{ACT}$,
es la causa primaria de~$c_s<38$~cr/trim.

\textbf{Undécimo}: el índice alternativo~$h=(\eta_t+\langle\eta\rangle)/2$,
donde~$\eta_t=C_{A,t-1}/45$ mide la fracción del máximo inscribible de créditos
aprobada el trimestre anterior y~$\langle\eta\rangle$ su promedio acumulado, elimina
el incentivo al subregistro~($\Delta h<0$ para todo~$t\geq1$;
$\Delta h(2)=-0.043$) y produce un ordenamiento entre estrategias monótonamente
coherente con la velocidad de egreso.
La estrategia~$h$-óptima ($c^\ast=38$~cr/trim con selección por~$a_i$ descendente)
cuantifica el valor de la información topológica: ignorar el orden del grafo de
seriaciones y preferir las UEAs fáciles cuesta~$2.1$~trimestres de egreso respecto
a la estrategia agresiva con idéntica carga ($\mathrm{E}[T\mid\text{tit}]_{h\text{-opt}}
= 19.3$~trim frente a~$17.2$~trim de la agresiva).
Sin embargo, ningún índice basado en flujo crediticio suprime la penalización
topológica: mientras las seriaciones permanezcan en el plan, reducirán~$C_A$
disponible y deprimirán~$P$ y~$h$ de forma análoga.
La reforma del criterio de prelación es condición necesaria pero no suficiente;
debe acompañar, no sustituir, la reforma curricular.

\textbf{Duodécimo}: el análisis global de complejidad sistémica
(\S\ref{sec:complejidad}) sintetiza los cinco mecanismos parciales en el índice
$C=\tfrac{1}{5}(D_1+D_2+D_3+D_4+D_5)\in[0,1]$ y clasifica los diez planes en cuatro
arquetipos de intervención diferenciada.
Los tres planes de mayor complejidad ---Computación ($C=0.617$), Química ($C=0.602$)
y Mecánica ($C=0.550$), clasificados~\emph{Alta}--- tienen como dimensión dominante
la docente~$D_4$ (Computación y Química) o la evaluativa~$D_3$ (Mecánica), no la
topológica.
Electrónica, que registra el peor $\ETtit=23.5$~trimestres de la División, ocupa
el séptimo lugar ($C=0.477$) porque su carga está concentrada en un único cuello de
botella topológico; la correlación entre el índice~$C$ y~$\ETtit$ es
$r_s=-0.139$ (no significativa, $N=10$), precisamente porque la complejidad
multi-dimensional difusa y la carga concentrada en un solo mecanismo son fenómenos
cualitativamente distintos que requieren respuestas distintas.
El acoplamiento negativo entre la dimensión de recuperación~$D_5$ y la dimensión
docente~$D_4$ ($r_s=-0.576$) y probabilística~$D_2$ ($r_s=-0.673$) identifica una
trampa auto-reforzada: los planes con mayor heterogeneidad docente y mayor varianza en
el tiempo de egreso tienen simultáneamente menor capacidad de reaprobación, de modo
que el rezago inicial se convierte en una condición difícilmente reversible para una
fracción significativa del alumnado.
Esta dinámica emergente no es observable desde ninguna de las cinco dimensiones
aisladas; emerge únicamente en el análisis conjunto y refuerza la conclusión de que
las reformas sectoriales parciales son insuficientes.

Para que los estudiantes puedan titular en el plazo del plan se requiere actuar
sobre cinco frentes de forma simultánea:
el \emph{topológico} (reducir las cadenas de seriación, priorizando las
132~aristas de mayor impacto);
el \emph{de aprobación} (apoyo académico focalizado en los nodos de la
Tabla~\ref{tab:top10_riesgo}, cuyas tasas de aprobación inferiores a~0.65
son el factor dominante en ocho de los diez planes);
el \emph{de carga} (incentivar o habilitar cargas cercanas a los~45~cr/trim
normales o hacia el máximo de~63~cr/trim);
el \emph{docente} (homogeneización de la práctica pedagógica en los nodos
\emph{Docente + intrínseca}, priorizando el departamento de Ciencias Básicas,
que con 272~docentes y~$\tilde{a}_d=0.670$ opera como la unidad de menor
eficacia de la División y cubre todos los nodos críticos del Tronco General;
el 31.3\% de los docentes CBI con~$a_d<0.65$ representa el margen de mejora
más accesible sin reforma curricular); y el \emph{evaluativo}
(reformar el criterio de prelación para eliminar el incentivo al subregistro,
sustituyendo o reequilibrando el componente~$\mathrm{ACT}$ por un indicador de
base entrópica como~$h$ que alinee el comportamiento racional del alumnado con la
velocidad de avance en el plan).
Los cinco son necesarios; ninguno es suficiente de forma aislada.

%\begin{okbox}
%El sistema de análisis está disponible en
%\texttt{Admin\_duties/Metanalisis/Analisis/}.
%Dado el grafo de seriaciones de cualquier plan propuesto y los datos históricos de
%aprobación, el sistema predice en minutos la distribución del tiempo de egreso, la
%tasa de egreso y el impacto marginal de cada seriación.
%El índice de riesgo topológico~$\rho_v$~(\ref{eq:risk}) y el ranking de impacto
%marginal están disponibles en formato CSV en \texttt{Output/R\_analisis/}
%para apoyar directamente las sesiones de revisión curricular.
%\end{okbox}



%══════════════════════════════════════════════════════
\clearpage
%\section*{Referencias}

{\small
\bibliographystyle{unsrt}
\bibliography{../../Analisis/referencias}
\addcontentsline{toc}{section}{Referencias}

}

\clearpage
\section*{Glosario de abreviaturas y notación}
\addcontentsline{toc}{section}{Glosario de abreviaturas y notación}
%══════════════════════════════════════════════════════

\subsection*{Siglas institucionales y de proceso académico}

\begin{tabularx}{\linewidth}{@{}lX@{}}
\toprule
Sigla & Significado \\
\midrule
AIC     & Criterio de Información de Akaike; usado para seleccionar el modelo
          de ajuste (Michaelis-Menten vs.\ exponencial) en la curva
          dosis-respuesta de remoción. \\[4pt]
BH      & Corrección de Benjamini-Hochberg; controla la tasa de falsos
          descubrimientos en comparaciones múltiples de correlación. \\[4pt]
CBI     & División de Ciencias Básicas e Ingeniería, UAM-Azcapotzalco. \\[4pt]
CSH     & División de Ciencias Sociales y Humanidades, UAM-Azcapotzalco. \\[4pt]
CYAD    & División de Ciencias y Artes para el Diseño, UAM-Azcapotzalco. \\[4pt]
DAG     & \emph{Directed Acyclic Graph} (grafo acíclico dirigido); estructura
          de datos que representa las relaciones de seriación entre UEAs. \\[4pt]
DFS     & \emph{Depth-First Search} (búsqueda en profundidad); algoritmo para
          calcular la profundidad topológica y el número de descendientes
          de cada nodo en el DAG. \\[4pt]
D+I     & \emph{Docente + Intrínseca}; clasificación de riesgo de UEA con
          $\mathrm{CV}>0.30$ y tasa de aprobación $a_i < 0.65$:
          la dificultad es tanto inherente al contenido como dependiente
          del docente asignado. \\[4pt]
ET      & Eficiencia Terminal; en el reporte acompañada siempre de un
          subíndice temporal, p.\,ej.\ $\ET{12}$ o $\ETtit$
          (véase Notación matemática). \\[4pt]
GEG     & Grupos de Evaluación Global; campo del diccionario de UEAs que
          registra los resultados de evaluación ordinaria de cada asignatura
          (base del cómputo de tasas de aprobación). \\[4pt]
LGES    & Ley General de Educación Superior (México). \\[4pt]
MM      & Modelo de Michaelis-Menten; curva de saturación
          $G(N)=V_{\max}N/(K_M+N)$ ajustada a la ganancia acumulada
          de remoción de seriaciones. \\[4pt]
NA      & No Acreditado; resultado de un intento fallido en una UEA.
          Acumular cinco~NA en la misma UEA produce baja definitiva
          del programa (RES~UAM Arts.~16.V y~44.XI). \\[4pt]
PD      & \emph{Principalmente Docente}; clasificación de riesgo de UEA con
          $\mathrm{CV}>0.30$ y $a_i\geq 0.65$: el contenido no es
          intrínsecamente difícil, pero la práctica docente diferenciada
          amplifica la varianza en los resultados. \\[4pt]
RES~UAM & Reglamento de Estudios Superiores de la Universidad Autónoma
          Metropolitana. \\[4pt]
TG      & Tronco General (compartido por todos los planes de CBI);
          incluye Cálculo Diferencial, Cálculo Integral, EDO,
          Cinemática, Física y Química. \\[4pt]
TGNA    & Tronco General Nuevo Actualizado; plan piloto analizado en el
          Reporte TGNA-UAM (2026) que sirve como caso de referencia
          metodológica. \\[4pt]
UAM-A   & Universidad Autónoma Metropolitana, Unidad Azcapotzalco. \\[4pt]
UEA     & Unidad de Enseñanza-Aprendizaje; la asignatura o materia en la
          nomenclatura UAM.  Identificada por un código de siete dígitos. \\
\bottomrule
\end{tabularx}

\bigskip
\subsection*{Notación matemática principal}

\begin{tabularx}{\linewidth}{@{}lX@{}}
\toprule
Símbolo & Definición \\
\midrule
$\ET{k}$
  & Eficiencia Terminal al trimestre~$k$: fracción de la cohorte de ingreso
    que ha completado todos los créditos y egresado en a lo sumo $k$~trimestres. \\[4pt]
$\ETtit = \mathbb{E}[T\mid\mathrm{tit}]$
  & Tiempo esperado de egreso en trimestres, condicionado a que el
    estudiante eventualmente egrese. \\[4pt]
$a_i$
  & Tasa de aprobación histórica de la UEA~$i$: probabilidad de acreditar
    la UEA en un intento dado, estimada como el promedio de aprobación
    en los grupos GEG del periodo 16I--25O. \\[4pt]
$c_s$
  & Carga semestral efectiva histórica (créditos/trimestre); se calibra
    como $c_s = \mathrm{cr\_total}/\ell$ a partir del tiempo oficial de
    egreso. Carga normal del plan: 45~cr/trim; máximo: 63~cr/trim. \\[4pt]
$\ell$
  & Duración efectiva de la trayectoria en trimestres, calculada como
    $\ell = \mathrm{cr\_total}/c_s$; es el denominador de la carga
    efectiva y el parámetro de calibración principal del modelo. \\[4pt]
$\Pacc$
  & Probabilidad de acceso a una UEA dada la cadena de seriaciones previas:
    $\Pacc = \prod_{j\in\mathrm{seriaciones}} P_{\leq k_j}(j)$
    donde $P_{\leq k}(j)$ es la probabilidad acumulada de haber aprobado
    $j$ en a lo sumo $k$ intentos. \\[4pt]
$\rho_v$
  & Índice de riesgo topológico enriquecido de la UEA~$v$
    (ecuación~\ref{eq:risk}); combina dificultad intrínseca, posición
    topológica en el DAG, heterogeneidad docente y probabilidad de
    expulsión. \\[4pt]
$\mathcal{A}(N)$
  & Índice de aditividad al eliminar $N$~aristas simultáneamente
    (ecuación~\ref{eq:aditividad}); ratio entre la ganancia combinada
    observada y la suma de los efectos marginales individuales.
    $\mathcal{A}<1$ indica subaditividad. \\[4pt]
$K_M$
  & Constante de Michaelis del modelo MM: número de aristas necesarias
    para capturar el 50\% de la ganancia máxima asintótica
    ($N_{50}^{\mathrm{MM}}=K_M$). \\[4pt]
$V_{\max}$
  & Ganancia máxima asintótica del modelo MM; no necesariamente
    realizable con el número finito de seriaciones del plan. \\[4pt]
CV
  & Coeficiente de Variación ($= \sigma/\mu$); usado para cuantificar
    la heterogeneidad de la eficacia entre los docentes que han
    impartido una UEA. \\[4pt]
$\delta a = a_i(\mathrm{rec})-a_i(\mathrm{eval})$
  & Diferencia de tasa de aprobación entre la evaluación de recuperación
    y la evaluación global ordinaria; mide si la segunda oportunidad
    institucional mejora los resultados. \\
\bottomrule
\end{tabularx}

\bigskip
\subsection*{Factores del análisis factorial de UEA (\S\ref{sec:factores_uea})}

\begin{tabularx}{\linewidth}{@{}lX@{}}
\toprule
Factor & Interpretación \\
\midrule
MR1 & Influencia topológica: carga alta en betweenness, grado de salida y
      número de descendientes.  Predice el impacto agregado del nodo
      sobre el sistema de seriaciones. \\[4pt]
MR2 & Dispersión docente: carga alta en CV y rango de eficacia entre
      profesores de la UEA.  Captura la componente pedagógica del riesgo,
      independiente de la dificultad intrínseca. \\[4pt]
MR3 & Calidad intrínseca: carga alta en eficacia, $a_i$ y eficiencia;
      carga negativa en spread$_{15}$.  Representa la facilidad/dificultad
      estructural del nodo independiente de su posición en el DAG. \\[4pt]
MR4 & Embeddedness curricular: carga positiva en profundidad y pagerank;
      carga negativa en número de planes.  Captura nodos especializados:
      profundos, presentes en pocos programas, de alta influencia local. \\[4pt]
MR5 & Recuperabilidad y peso crediticio: carga alta en fracción de
      recuperación y créditos.  Identifica UEAs donde los estudiantes
      realizan intentos adicionales y con mayor peso en la trayectoria. \\
\bottomrule
\end{tabularx}


%%══════════════════════════════════════════════════════
%\section{Comunicación institucional del diagnóstico}
%\label{sec:comunicacion}
%%══════════════════════════════════════════════════════
%
%El análisis técnico de las secciones precedentes produce un diagnóstico preciso
%pero de formulación exigente: $\ET{12}$ (eficiencia terminal a 12 trimestres)
%estadísticamente próxima a cero, una brecha de carga de 15~cr/trim entre diseño y realidad, un 31\%
%del profesorado por debajo del umbral de riesgo y una cadena de seriaciones
%que bloquea el egreso en el plazo del plan en los diez programas.
%Presentar estos resultados sin una estrategia de comunicación diferenciada
%según la audiencia es, en sí mismo, un riesgo para la reforma: la respuesta más
%probable de un órgano colegiado ante una formulación acusatoria es la
%impugnación del modelo, no la revisión del plan.
%Esta sección desarrolla la estrategia de comunicación diferenciada y el análisis
%de los actores institucionales relevantes.
%
%\subsection{Versiones diferenciadas por audiencia}
%\label{sec:comunicacion_estrategica}
%
%El módulo~24 (\texttt{24\_comunicacion\_estrategica.py}) genera tres versiones del
%hallazgo central para tres audiencias distintas, disponibles en
%\texttt{Output/comunicacion/}.
%Cada versión incluye un resumen ejecutivo de máximo 400 palabras, una tabla de
%hallazgos adaptada y los tres mensajes clave con las respuestas recomendadas ante
%las preguntas difíciles anticipadas.
%
%La \textbf{versión técnica}, destinada a jefes de departamento con perfil
%cuantitativo, emplea el encuadre completo de la versión~5.2 del reporte: el
%desajuste de carga como diagnóstico primario, $\ET{12}$ estadísticamente próxima
%a cero con intervalo de confianza (IC) explícito, los dos mecanismos estructurales,
%la cascada estática y el modelo estocástico completo.
%
%La \textbf{versión institucional}, para Junta de la División de Ciencias Básicas
%e Ingeniería (CBI) y directores, reformula el hallazgo en términos de
%\emph{brecha entre el plazo del plan y el tiempo alcanzable bajo condiciones
%históricas}, evitando expresamente los términos <<imposible>> y
%<<$\ET{12}=0$>>.
%Enfatiza que el 27\% de los prerrequisitos concentra el 80\% del beneficio
%reformable y que Electrónica y Química son candidatos a reforma quirúrgica con
%solo ocho y cinco ajustes, respectivamente.
%
%La \textbf{versión de política pública}, para Rectoría, acreditadores y la Asociación Nacional de Universidades e Instituciones de Educación Superior (ANUIES),
%enmarca los resultados en el contexto OCDE (brecha entre el plazo del plan y el
%real en las ingenierías mexicanas), cita el marco de la Ley General de Educación
%Superior (LGES~\cite{LGES2021}, Arts.~10 y~40--41) y posiciona a la
%UAM-Azcapotzalco como institución que identifica proactivamente sus problemas
%estructurales y diseña soluciones basadas en evidencia.
%
%\subsection{Análisis de actores institucionales y gestión de la resistencia}
%\label{sec:stakeholders}
%
%El análisis docente (\S\ref{sec:perfil_docente}, tabla~\ref{tab:dept_docentes})
%señala que el departamento de Ciencias Básicas es la unidad de menor eficacia y
%mayor dispersión de la División.
%Este señalamiento, presentado sin preparación contextual, puede producir
%resistencia defensiva en Junta que bloquee la deliberación sobre el fondo antes
%de que comience.
%La secuencia de presentación importa tanto como el contenido.
%
%El módulo~25 (\texttt{25\_analisis\_stakeholders.py}) construye una matriz de
%ocho actores institucionales con posición probable, fuente del interés, palanca
%de influencia e intervención recomendada.
%Las coaliciones mínimas necesarias para cada uno de los cuatro frentes de
%intervención en Consejo Divisional están especificadas en
%\texttt{Output/comunicacion/stakeholders.md}.
%
%La estrategia recomendada opera en cinco fases.
%La primera consiste en sesiones individuales con coordinadores de plan durante
%las semanas~1--2, con el objetivo de validar los datos de seriaciones plan por plan
%y sin abordar aún el análisis docente.
%La segunda, en la semana~3, presenta a los jefes de departamento los hallazgos de
%heterogeneidad docente comenzando por Ingeniería Geotécnica
%(CV$=0.93$, Departamento de Materiales) antes que por Cálculo Diferencial
%(CV$=0.48$, Departamento de Ciencias Básicas), estableciendo que la heterogeneidad
%docente es un fenómeno divisional y no concentrado en un solo departamento.
%La tercera fase, en la semana~4, es una sesión bilateral con el Jefe del
%Departamento de Ciencias Básicas en la que se ofrece co-autoría del protocolo
%de homogeneización pedagógica como oportunidad de liderazgo, no como señalamiento
%negativo.
%La cuarta fase, en la semana~6, es la presentación ante la Junta de CBI con la
%versión institucional del diagnóstico y la propuesta de reforma en tres etapas
%progresivas.
%La quinta fase, durante las semanas~7--10, cubre la negociación sindical con el
%programa de formación docente enmarcado como convenio de desarrollo académico,
%no como evaluación punitiva.
%
%%══════════════════════════════════════════════════════
%\section{Plan de implementación}
%\label{sec:implementacion}
%%══════════════════════════════════════════════════════
%
%Los cuatro frentes de intervención identificados en la Sección~\ref{sec:conclusiones}
%---topológico, de aprobación, de carga y docente--- son necesarios de forma
%simultánea, pero la capacidad institucional de la División no permite ejecutarlos
%en paralelo con la misma intensidad desde el primer ciclo.
%La reforma curricular formal en la UAM-Azcapotzalco requiere en la práctica más
%de dos años de tramitación; la intervención docente en el Departamento de Ciencias
%Básicas exige coordinación sindical.
%Esta sección presenta la hoja de ruta de implementación que respeta estas
%restricciones, el modelo de institucionalización del sistema de análisis y el
%modelo de trayectorias de abandono para los estudiantes que no egresan.
%
%\subsection{Hoja de ruta por fases (horizonte 2026--2028)}
%\label{sec:hoja_de_ruta}
%
%El módulo~27 (\texttt{27\_hoja\_de\_ruta.py}) modela la capacidad institucional
%como restricción explícita: máximo una reforma curricular formal por ciclo
%académico anual y máximo dos intervenciones docentes piloto simultáneas por
%departamento.
%Bajo esas restricciones, la secuencia óptima que maximiza la reducción de~$\ETtit$
%(tiempo esperado de egreso condicional al egreso) en el menor tiempo posible
%es la siguiente, en un horizonte de cuatro ciclos anuales.
%
%En el \textbf{Ciclo~1} (2026) se lanza el programa de tutorías en las Unidades
%de Enseñanza-Aprendizaje (UEA) críticas del Frente de Aprobación ---acción
%que no requiere aprobación en Consejo Divisional--- con la métrica de éxito de
%un incremento $\geq 5$ puntos porcentuales (pp) en la tasa de aprobación~$a_i$
%de Cálculo Diferencial.
%Simultáneamente se inicia el proceso formal de revisión curricular en Electrónica
%y Química (Frente Topológico), sometiendo ante la Comisión de Planes y Programas
%los portfolios de $8 + 5$ aristas identificados en la Sección~\ref{sec:portfolio}.
%
%En el \textbf{Ciclo~2} (2026--2027) se implementa la reforma aprobada en elo y
%qui, se lanza el piloto docente en el Departamento de Ciencias Básicas (dos
%cursos piloto, Frente Docente) y se inicia la segunda ronda de revisión curricular
%en amb, ele y met.
%Este ciclo marca el \textbf{primer cambio detectable} en~$\ETtit$: se espera una
%reducción $\geq 0.3$~trim en Electrónica o Química a partir de los primeros
%datos de la cohorte que curse el plan reformado.
%
%En el \textbf{Ciclo~3} (2027) se aprueba la reforma en amb, ele y met, se lanza
%el incentivo de carga académica (Frente de Carga, mediante convenio con Servicios
%Escolares) y se amplía el programa docente a todos los grupos de alto coeficiente
%de variación (CV) de eficacia.
%La métrica de seguimiento es $\ETtit(\mathrm{CBI})\leq 19.5$~trim.
%
%En el \textbf{Ciclo~4} (2027--2028) se completa la reforma topológica en los
%cinco planes restantes (fis, mec, civ, com, ind).
%La métrica de cierre de fase es $\ETtit(\mathrm{CBI})\leq 18.5$~trim y
%$\ET{18}\geq 30\%$ en Electrónica [donde $\ET{18}$ es la eficiencia terminal
%a 18~trimestres].
%
%La simulación del impacto sobre la cohorte activa ($N=6\,562$~estudiantes)
%indica que al final del Ciclo~1 aproximadamente~$1\,348$~estudiantes de los
%planes elo y qui estarán bajo planes con seriaciones revisadas, y que al final
%del Ciclo~3 los~$6\,562$~activos estarán bajo al menos un frente de intervención
%activo.
%El diagrama Gantt completo está en \texttt{Output/figuras/fig\_gantt.pdf} y los
%datos de impacto proyectado en \texttt{Output/implementacion/hoja\_de\_ruta.md}.
%
%\subsection{Institucionalización del sistema de análisis}
%\label{sec:institucionalizacion}
%
%El sistema de análisis depende actualmente de la continuidad de una sola persona
%y carece de gobernanza, protocolos de actualización ni validación
%institucionalizada.
%Para que los hallazgos de este reporte sean reproducibles en ciclos futuros y
%para que las decisiones curriculares puedan apoyarse en datos actualizados,
%el sistema debe transferirse a una unidad institucional permanente.
%
%El módulo~28 (\texttt{28\_institucionalizacion.py}) genera los tres componentes
%necesarios para esta transferencia.
%El primero es el \textbf{manual de operación}
%(\texttt{Output/docs/manual\_operacion.md}), que especifica los requisitos exactos
%de software, el procedimiento de actualización trimestral, el protocolo de
%validación del calibrador (tolerancia $\leq 1$~trim) y el \emph{checklist}
%de integridad de datos previo a cada corrida.
%El segundo es el \textbf{protocolo de gobernanza}
%(\texttt{Output/docs/gobernanza.md}), que asigna roles institucionales
%---propietario técnico, coordinador de datos, validador, custodio del
%repositorio---, define la frecuencia de actualización, el proceso de validación
%cruzada con Servicios Escolares y el procedimiento para incorporar cambios
%curriculares aprobados.
%El tercero es la \textbf{suite de tests automatizados} (\texttt{Analisis/tests/}),
%con tres módulos que verifican en cada ejecución la integridad del grafo acíclico
%dirigido (DAG) de seriaciones, la coherencia de las tasas de aprobación
%($a_i\in(0,1]$ para todas las UEAs), la convergencia del calibrador (error
%$<1$~trim en todos los planes comparables) y la reproducibilidad de~$\ETtit$
%con semilla fija.
%
%El rol institucional recomendado para la propiedad del sistema es la Coordinación
%del Nuevo Tronco General de Asignaturas (TGNA), que ya es propietaria de la
%metodología de referencia~\cite{TGNA2026} y tiene el mandato institucional más
%cercano al seguimiento curricular continuo que este tipo de análisis requiere.
%
%\subsection{Modelo de trayectorias de abandono}
%\label{sec:modelo_abandono}
%
%El análisis de las secciones precedentes optimiza intervenciones para los
%estudiantes que eventualmente egresan.
%Sin embargo, el~$36.3\%$ de la cohorte activa ($\approx 2\,380$~estudiantes)
%está proyectado a causar baja definitiva por acumulación de cinco evaluaciones
%No Acreditadas (NA) en alguna UEA, y ese grupo representa el problema de mayor
%urgencia social del diagnóstico.
%
%El módulo~29 (\texttt{29\_modelo\_abandono.py}) extiende el sistema para
%caracterizar y modelar las trayectorias de abandono.
%Para cada estudiante simulado se registran en cada trimestre los NA acumulados
%por UEA, la posición en el DAG y el trimestre de baja, cuando esta ocurre;
%esto permite estimar curvas de supervivencia tipo Kaplan-Meier de abandono por
%plan y ajustar regresiones de Cox sobre rasgos de trayectoria.
%
%Los hallazgos preliminares muestran que el~$50\%$ de las bajas se concentra
%entre los trimestres~4 y~8, y que las UEAs del Tronco General ---Cálculo
%Diferencial, Ecuaciones Diferenciales Ordinarias (EDO) y Cinemática y Dinámica
%de Partículas--- son las que más frecuentemente acumulan los NA que desencadenan
%la baja definitiva.
%
%El módulo simula tres intervenciones diseñadas para reducir la tasa de baja del
%$36\%$ al~$25\%$ como objetivo conservador.
%La \textbf{intervención temprana} ---tutorías en los trimestres~2--4--- proyecta
%un incremento de~$+8$~pp en la tasa de aprobación de las cinco UEAs de mayor
%alcance topológico durante los trimestres críticos de abandono, con un costo
%estimado de aproximadamente 20 grupos adicionales de nivelación por trimestre.
%La \textbf{intervención de recuperación} propone extender el límite de evaluaciones
%No Acreditadas más allá del máximo actual para las UEAs con probabilidad histórica
%de expulsión $P_{\mathrm{exp}}>0.05$; por afectar los Arts.~16.V y~44.XI del
%RES~UAM, requiere iniciativa ante el Consejo Académico de la institución, no
%solo una decisión divisional.
%La \textbf{intervención de carga} aumenta en un~$20\%$ la carga efectiva
%inscrita (de~30 a~36~cr/trim), con mayor impacto sobre la tasa de egreso
%eventual que sobre la tasa de baja.
%
%Para cada intervención se reportan la reducción esperada de la tasa de baja,
%el número de estudiantes de la cohorte activa directamente beneficiados y el
%costo estimado en grupos adicionales.
%Los resultados completos están disponibles en
%\texttt{Output/abandono\_trayectorias.csv} y
%\texttt{Output/abandono\_intervenciones.csv}.
%
%══════════════════════════════════════════════════════
\clearpage
\section*{Apéndice A\@.~Inventario de scripts de análisis}
\addcontentsline{toc}{section}{Apéndice A. Inventario de scripts de análisis}
%══════════════════════════════════════════════════════

Los scripts residen en \texttt{Admin\_duties/Metanalisis/Analisis/}.
Los prefijos numéricos indican el orden de ejecución; los scripts
R~\cite{RCoreTeam2024} son independientes entre sí y pueden ejecutarse en
cualquier orden una vez que los scripts Python (05--10) hayan generado sus
salidas CSV en \texttt{Output/datos/}.
El sistema comprende 21~módulos Python de análisis (incluido el módulo
compartido \texttt{simulacion.py}), 6~scripts R de análisis estadístico
consolidados, 5~scripts de gestión institucional en \texttt{delivery/}
y 4~utilidades en \texttt{utils/}.

\subsection*{Módulo compartido}

\texttt{simulacion.py}
es el núcleo importable del sistema: implementa el motor Monte Carlo trimestre
a trimestre y expone las funciones \texttt{simulate}, \texttt{metrics},
\texttt{calibrar}, \texttt{load\_plan}, \texttt{uea\_stats}, \texttt{uea\_p}
y \texttt{heterocaps}.
Centraliza las constantes $T_{\min}=12$, $T_{\max}=30$, $k_{\max}=5$~intentos,
$c_0=45$~cr/trim y $\sigma_0=0.6$, así como las rutas canónicas
\texttt{GRAFOS}, \texttt{DATOS} y \texttt{OUT}.
Todos los scripts de análisis (05--37) lo importan; la modificación de
cualquier parámetro en este módulo propaga automáticamente a todo el sistema.

\subsection*{Núcleo del análisis (scripts 01--10, 18, 30)}

\texttt{01\_extraer\_seriaciones.py}
construye el grafo canónico de seriaciones de los diez planes~2020:
valida los JSON por plan en \texttt{Grafos/}, asigna las tasas de aprobación
desde \texttt{diccionario\_ueas.json} y produce el grafo unificado
\texttt{seriaciones\_unificado.json} más tres CSV de resumen en \texttt{Output/}.

\texttt{02\_estadisticas\_ueas.py}
implementa el modelo de cascada estática: calcula la probabilidad de acceso
acumulada $P_{\mathrm{reach}}(u)$ para cada UEA obligatoria bajo los modelos
DAG y camino-crítico, y produce \texttt{cascade\_por\_uea.csv}.

\texttt{03\_generar\_diagramas.py}
genera diagramas de seriación interactivos (SVG/HTML) con layout Sugiyama
simplificado: rectángulos coloreados según $P_{\mathrm{reach}}$, flechas
de seriación directa y de corregistro etiquetadas con la tasa de aprobación.

\texttt{04\_modelo\_eficiencia\_terminal.py}
calcula la eficiencia terminal por plan bajo el supuesto de independencia entre
UEAs, usando los intentos disponibles $k_u(T)$ y los vectores de probabilidad
acumulada del diccionario; sirve como referencia analítica para la calibración.

\texttt{05\_modelo\_topologico.py}
es el núcleo del análisis: simulación Monte Carlo trimestre a trimestre con
inscripción proporcional, renuncia, aprobación con fragilidad latente
compartida~(\ref{eq:fragilidad}),
$a_{s,i}=\Phi(\Phi^{-1}(a_i)\sqrt{1+\sigma^2}+\sigma\theta_s)$,
límite reglamentario de cinco oportunidades y baja definitiva;
calibra la carga efectiva histórica~$\ell$ por bisección imponiendo
$\ETtit^{\,\mathrm{mod}}(p)=\ETtit^{\,\mathrm{ofic}}(p)$
(ecuación~\ref{eq:calibracion}).

\texttt{06\_figuras\_topologico.py}
produce las figuras de publicación del modelo estocástico: curvas~$\mathrm{ET}(t)$
por plan, distribución de percentiles del tiempo al egreso, calibración
$\mathrm{E}[T|\text{tit}]$ modelo~vs.~oficial y sensibilidad a~$\sigma$.

\texttt{07\_efecto\_seriaciones\_todos.py}
cuantifica el efecto marginal de cada arista en los diez planes: elimina
una a una todas las seriaciones y corregistros entre obligatorias y mide el
cambio en $\mathrm{ET}(12)$, tasa de graduación y~$\mathrm{E}[T|\text{tit}]$
respecto al baseline calibrado usando el mismo seed para reducir varianza.

\texttt{10\_remocion\_combinada.py}
calcula la curva dosis-respuesta de remoción simultánea: para cada plan
elimina las top-$N$ aristas ordenadas por impacto marginal, ejecuta~$K$
réplicas Monte Carlo con semillas independientes y registra la ganancia
observada, la predicción aditiva y el índice de subaditividad~$A(N)$.

\texttt{18\_analisis\_profesores.py}
produce las estadísticas del cuerpo docente CBI (1\,810~profesores,
periodos~16I--25O): distribución de la eficacia~$a_d$ por departamento,
fracción de docentes en el umbral de riesgo~$a_d < 0.65$ y cruce con
los nodos críticos del ranking de riesgo topológico.

\texttt{30\_incertidumbre\_ET12.py}
cuantifica la incertidumbre en $\ET{12}$ mediante IC~Monte Carlo exacto
(Clopper-Pearson, $S=20\,000$ por plan) y análisis de sensibilidad a
$c_s\pm\Delta$ para $\Delta\in\{-2,\ldots,+4\}$~cr/trim.
Resultado: $\ET{12}\in[0.000\%,\,0.018\%]$~(IC~MC~95\%) bajo el modelo;
$1.52\%$~(IC~95\%~CP: $[1.11\%,\,2.04\%]$) históricamente.
La brecha de $1.5\,\text{pp}$ cuantifica el efecto de la heterogeneidad individual
de carga no modelada (\S\ref{sec:hallazgo}).
Salida: \texttt{Output/datos/simulacion/incertidumbre\_ET12.csv},
        \texttt{Output/figuras/fig\_incertidumbre\_ET12.pdf}.

\subsection*{Módulos de análisis cuantitativo (scripts 20--24)}

\texttt{20\_descomposicion\_carga.py}
descompone la carga efectiva histórica~$c_s$ en componente estructural
(máximo inscribible según la topología con aprobación perfecta) y componente
conductual (residual), y corre el modelo Monte Carlo bajo tres escenarios
para cuantificar cuánto del retraso es reforma-topológica versus conductual.

\texttt{21\_modelo\_varianza\_docente.py}
incorpora varianza docente explícita en el modelo estocástico: en cada
trimestre sortea un docente del pool histórico de la UEA y usa su tasa
de aprobación individual; reporta $\ETtit$ e intervalos de confianza~$95\%$
bajo cuatro condiciones de asignación docente.

\texttt{22\_optimizador\_portfolio.py}
implementa tres estrategias de selección de portfolios de remoción para los
planes con curvas dosis-respuesta no-monótonas: greedy puro, greedy con
backtracking ($k=5$) y simulated annealing; identifica el subconjunto óptimo
de~80 aristas con mayor beneficio neto y reporta~$G_{\max}$ por plan.

\texttt{23\_modelo\_aditividad.py}
ajusta por OLS el modelo predictivo del índice de subaditividad~$\mathcal{A}(N)$
sobre cinco predictores topológicos del subconjunto de aristas eliminadas
(fracción de origen compartido, fracción de destino compartido, diversidad de
fuentes, intermediación media, profundidad media); valida por leave-one-plan-out.
Salida: \texttt{Output/R\_analisis/modelo\_aditividad\_coefs.csv}.

\texttt{24\_modelos\_saturacion.py}
genera las figuras del ajuste de saturación (\S\ref{sec:modelos_sat}):
diez paneles con $G(N)$ observado y curvas Michaelis-Menten / exponencial
ajustadas, con modelo preferido por AIC y marcas~$V_{\max}$, $N_{80}$;
segundo panel de comparación de parámetros.
Salida: \texttt{Output/figuras\_correlaciones/fig\_sat\_curvas.pdf},
        \texttt{fig\_sat\_params.pdf}.

\subsection*{Módulos de índices de desempeño y prelación (\S\ref{sec:indice_desempenio})}

\texttt{31\_indice\_desempenio\_simulacion.py}
simula trayectorias del índice~$P$ vigente (Acuerdo~551.4.4.1-2015) y del
índice alternativo~$h$ bajo el currículo seriado y bajo remoción de las
aristas de mayor impacto; cuantifica la doble penalización topológica de~$P$
($\mathrm{RA}<1$ crónico, $\mathrm{GA}$ bloqueado en los cuellos de Electrónica)
y la extinción del incentivo al subregistro bajo~$h$.
Salida: \texttt{Output/datos/indices/indice\_P\_trayectorias.csv},
        \texttt{indice\_P\_componentes.csv}, \texttt{indice\_P\_bloqueados.csv}.

\texttt{33\_estrategias\_inscripcion.py}
simula tres estrategias de inscripción (agresiva, baseline, conservadora) y
mide su efecto sobre~$P$, $h$, $\mathrm{RA}$, $\mathrm{GA}$ y~$\mathrm{ACT}$;
demuestra que la estrategia conservadora (24~cr/trim, orden fácil-primero)
maximiza $\mathrm{ACT}$ en $T\in[1,2]$ al costo de reducir el avance absoluto,
corroborando el incentivo al subregistro.
Salida: \texttt{Output/datos/indices/estrategias\_trayectorias.csv}.

\texttt{35\_indice\_h\_estrategias.py}
simula cuatro estrategias bajo la lógica del índice~$h$ (agresiva, baseline,
conservadora y $h$-óptima: máxima carga en orden fácil-primero);
muestra que bajo~$h$ no existe estrategia de subregistro rentable
---el único camino de mejora es aprobar créditos absolutos--- y que
la estrategia $h$-óptima se traduce en~$\ETtit=17.2$~trim frente a~$19.3$~trim
de la conservadora.
Salida: \texttt{Output/datos/indices/indice\_h\_trayectorias.csv}.

\texttt{04\_indices.R}
produce las figuras del análisis del índice~$P$ (trayectorias $P(t)$ por escenario,
$\mathrm{RA}$ crónico, doble penalización $\Delta P$, comparación estructural
$P$~vs.~$h$), las figuras de estrategias de inscripción (ventaja efímera de
$\mathrm{ACT}$, trampa ACT en el plano $\mathrm{E}[T|\text{tit}]$~vs.~$\bar{P}$,
$\Delta P(t)$ con cruce por cero) y el análisis comparativo completo
$P$~vs.~$h$ (extinción del incentivo al subregistro, coherencia de~$h$ con el
egreso y penalización topológica en Electrónica).
Consolida los scripts previamente independientes
\texttt{32\_analisis\_indice\_P.R}, \texttt{34\_analisis\_estrategias.R}
y~\texttt{36\_analisis\_indice\_h.R}.

\subsection*{Módulo de complejidad sistémica (\S\ref{sec:complejidad})}

\texttt{37\_complejidad\_global.py}
construye el índice compuesto de complejidad sistémica~$C\in[0,1]$ sobre
cinco dimensiones normalizadas: D1~topológica (profundidad DAG, densidad,
concentración de betweenness), D2~probabilística (entropía del tiempo de
egreso, dispersión, dropout), D3~evaluativa (depresión de~$P$, trampa ACT,
brecha $P$--$h$), D4~docente (heterogeneidad intra-UEA, rango, impacto
simulado) y D5~recuperación (ganancia por reintentos,
$P_{\mathrm{exp}}$).
Clasifica los diez planes en tres rangos (Alta, Moderada, Baja) y
calcula la matriz de acoplamiento entre dimensiones.
Salida: \texttt{Output/datos/complejidad/complejidad\_global.csv},
        \texttt{complejidad\_normalizada.csv}, \texttt{complejidad\_coupling.csv}.

\texttt{38\_visualizacion\_complejidad.R}
produce cinco figuras del análisis de complejidad: diagrama de radar de las
cinco dimensiones por plan, mapa de calor de dimensiones normalizadas, índice~$C$
frente a~$\ETtit$ y tasa de egreso, matriz de acoplamiento inter-dimensional
(Spearman) y descomposición apilada de~$C$ por plan.
Salida: \texttt{Output/figuras\_correlaciones/fig\_37a}--\texttt{fig\_37e.pdf}.

\subsection*{Scripts R de análisis estadístico consolidados}

Los dieciocho scripts R independientes originales (08--19 y 23R) fueron
consolidados en cinco archivos temáticos; cada archivo puede ejecutarse
de forma autónoma una vez disponibles los CSV de simulación en
\texttt{Output/datos/}.

\texttt{01\_correlaciones.R}
realiza el análisis estadístico de correlaciones a tres niveles
---plan ($N=10$), UEA ($N=689$) y arista ($N=499$)---:
correlaciones de Spearman con prueba de permutaciones, regresión
multivariada e identificación de los predictores estructurales
de~$\mathrm{ET}$; genera los dispersogramas plan~vs.~$\mathrm{ET}$
(Fig.~\ref{fig:scatter_plan}), el mapa de calor de correlaciones por
par de variables, y profundiza con correlaciones parciales de Spearman
controlando covariables de acceso acumulado según la recursión de
cascada~(\ref{eq:cascada}), análisis factorial exploratorio con
rotación oblimin~\cite{Horn1965} ($N=606$~UEAs, 5~factores por análisis
paralelo) y predictores a nivel de arista.
Consolida \texttt{08\_correlaciones\_R.R}, \texttt{09\_visualizaciones.R}
y \texttt{19\_correlaciones\_profundas.R}.

\texttt{02\_remocion.R}
ajusta los modelos de saturación (Michaelis-Menten~\cite{Michaelis1913},
exponencial y lineal) sobre las curvas dosis-respuesta de cada plan,
selecciona el mejor ajuste por AIC, estima $N_{50}$, $N_{80}$ y~$N_{95}$
y computa el índice de aditividad~$A(N)$; genera las figuras de remoción
combinada (dosis-respuesta, portfolio \textit{lollipop}, mapa de aditividad
y clasificación estratégica de los diez planes) y valida el modelo
de aditividad OLS de \texttt{23\_modelo\_aditividad.py}.
Consolida \texttt{11\_analisis\_remocion.R}, \texttt{12\_figuras\_adicionales.R}
y \texttt{23\_modelo\_aditividad.R}.

\texttt{03\_riesgo\_ueas.R}
cuantifica la recuperabilidad de cada UEA mediante
$\mathrm{spread}=P(k\leq 5)-P(k=1)$ y clasifica los nodos en los cuadrantes
recuperable, persistentemente difícil e intrínsecamente fácil; calcula el
gap de recuperación $\delta a_i$ y su correlación con la centralidad
topológica; estima la varianza docente intra-UEA y la probabilidad de
expulsión definitiva~$P_{\mathrm{exp}}$; y construye el índice de riesgo
topológico enriquecido
$\rho_v=(1-a_i)\times\mathrm{betw}\times\ln(1+\mathrm{desc})
\times(1+\mathrm{CV})\times(1+P_{\mathrm{exp}})$,
generando el ranking definitivo de UEAs de alto riesgo
(Tabla~\ref{tab:top10_riesgo}) y los mapas de riesgo por plan.
Consolida \texttt{13\_recoverability.R}, \texttt{14\_gap\_recuperacion.R},
\texttt{15\_heterogeneidad\_docente.R}, \texttt{16\_prob\_expulsion.R}
y \texttt{17\_risk\_score\_enriquecido.R}.

\texttt{06\_infografias.R}
genera las versiones modo oscuro (\texttt{\_info.pdf}) y modo claro
(\texttt{\_lm.pdf}) de las figuras destinadas a las cinco infografías:
clasificación estratégica de planes por cuadrante, radar de complejidad
sistémica, portfolio de remoción, curvas dosis-respuesta y comparación
de índices $P$~vs.~$h$.
Salida: \texttt{Output/figuras\_correlaciones/fig\_17\_clasificacion\_planes\_info.pdf},
        \texttt{fig\_17\_clasificacion\_planes\_lm.pdf} y equivalentes.

\subsection*{Scripts de análisis institucional y gestión (\texttt{delivery/})}

Estos cinco scripts residen en \texttt{Analisis/delivery/} y producen
documentos de gestión y comunicación; no forman parte del pipeline
estadístico.

\texttt{24\_comunicacion\_estrategica.py}
genera tres versiones del hallazgo central para tres audiencias (técnica,
institucional, política pública): resúmenes $\leq 400$~palabras, tablas de
hallazgos, tres mensajes clave y preguntas difíciles anticipadas con
respuestas recomendadas.

\texttt{25\_analisis\_stakeholders.py}
construye la matriz de ocho actores institucionales con posición, fuente de
interés, palanca de influencia e intervención recomendada; especifica coaliciones
mínimas por frente de reforma y la secuencia óptima de presentación en cinco fases.

\texttt{26\_gestion\_riesgo\_legal.py}
identifica tres escenarios de riesgo legal con probabilidad estimada, argumentos
de defensa y acciones preventivas; genera las cláusulas de alcance para
el~\S\ref{sec:limitaciones} del reporte (\texttt{Output/legal/}).

\texttt{27\_hoja\_de\_ruta.py}
modela la capacidad institucional como restricción explícita y genera la hoja
de ruta en tres ciclos anuales (2026--2028) con métricas de éxito, diagrama
Gantt e impacto estimado sobre la cohorte activa.

\texttt{28\_institucionalizacion.py}
genera el manual de operación, el protocolo de gobernanza con roles
institucionales definidos, y una suite de tests automatizados
(\texttt{Analisis/tests/}) que verifican integridad, calibración y
reproducibilidad en cada actualización trimestral.

\texttt{29\_modelo\_abandono.py}
extiende el modelo Monte Carlo para caracterizar trayectorias de abandono:
curvas Kaplan-Meier de baja por plan, $t_{50}^{\mathrm{baja}}$ y simulación
de tres intervenciones (tutorías tempranas, recuperación extendida, incremento
de carga) con estimación del número de estudiantes beneficiados y costo
en grupos adicionales.

\subsection*{Utilidades (\texttt{utils/})}

Los cuatro scripts en \texttt{Analisis/utils/} son herramientas de
mantenimiento de calidad textual; no generan datos ni figuras.

\texttt{corrector\_estilo.py}
detecta anglicismos y variaciones terminológicas inconsistentes en el
archivo \texttt{.tex} del reporte, preservando comandos \LaTeX\ y
expresiones matemáticas.

\texttt{fix\_terminologia.py}
aplica un glosario de sustituciones terminológicas canónicas al texto
del reporte (por ejemplo, normalizar variantes de ``seriación'',
``tasa de aprobación'' y ``eficiencia terminal'').

\texttt{fix\_reglamentario.py}
verifica que las referencias al Reglamento de Estudios y al Acuerdo~551
usen la numeración oficial vigente y aplica correcciones en bloque.

\texttt{check\_ortografia\_scripts.py}
revisa los comentarios y cadenas de texto de los propios scripts Python
y~R en busca de inconsistencias ortográficas y terminológicas.

\end{document}
